% Utilisation d’un logiciel de publication — Informatique 3ème — Cours signé OnBuch+
% Programme officiel MINESEC. Compiler avec : tectonic N08-utilisation-logiciel-publication.tex
\newcommand{\DOCMATIERE}{Informatique}
\newcommand{\DOCNIVEAU}{3ème}
\newcommand{\DOCMODULE}{Informatique – classe de 3ème}
\newcommand{\DOCLECON}{N08}
\newcommand{\DOCTITRE}{Utilisation d’un logiciel de publication}
% OnBuch+ — préambule « cours signés » ENRICHI (compilé avec tectonic / XeLaTeX)
% Système visuel « Pop » d'OnBuch+ : fond crème (jamais blanc), bordures encre
% épaisses, ombres « dures » décalées, titres Archivo Black en capitales.
% Version MATHÉMATIQUES : graphes (pgfplots/tikz), boîtes pédagogiques
% (définition, propriété, méthode, attention, le savais-tu, exercice résolu,
% exercice, TP, bilan), page de garde et signature OnBuch+ ancrée en bas.
%
% Macros attendues dans main.tex AVANT \input{preamble} :
%   \DOCMATIERE  \DOCNIVEAU  \DOCMODULE  \DOCLECON (numéro)  \DOCTITRE
\documentclass[11pt]{article}

\usepackage{fontspec}
\usepackage[french]{babel}
\usepackage[a4paper,margin=2cm,top=3.4cm,bottom=2.8cm,headheight=1.4cm,footskip=1.3cm]{geometry}
\usepackage{amsmath,amssymb,amsfonts}
\usepackage{mathtools} % \xmapsto, \coloneqq... souvent utilisés par les modèles
\usepackage{cancel} % simplifications barrées (bilans de forces)
% \abs{z} (module, valeur absolue) et \norm{u} (norme), utilisés par les modèles
\providecommand{\abs}{}\let\abs\relax\DeclarePairedDelimiter{\abs}{\lvert}{\rvert}
\providecommand{\norm}{}\let\norm\relax\DeclarePairedDelimiter{\norm}{\lVert}{\rVert}
\usepackage[table]{xcolor} % [table] : fournit aussi \rowcolors (alternance de lignes)
\usepackage{pagecolor}
\usepackage{tikz}
\usetikzlibrary{arrows.meta,positioning,calc,shapes.geometric,shapes.misc,shapes.arrows,shapes.symbols,decorations.pathmorphing,decorations.pathreplacing,decorations.markings,patterns,shadows,mindmap,trees,fit,backgrounds,matrix,intersections,angles,quotes}
\usepackage{pgfplots}
\usepackage[version=4]{mhchem} % équations nucléaires \ce{^{235}_{92}U}
\usepackage[european,siunitx]{circuitikz} % schémas électriques
\pgfplotsset{compat=1.18}
\usepgfplotslibrary{fillbetween,groupplots} % aires entre courbes (calcul intégral)
\usepackage{enumitem}
\usepackage{fancyhdr}
% [most] charge toutes les bibliothèques tcolorbox (listings, minted, raster,
% documentation...) et gonfle inutilement la mémoire TeX sur un document long
% (~300 boîtes) : on ne charge que ce qui est réellement utilisé.
\usepackage[skins,breakable]{tcolorbox}
\usepackage{graphicx}
\usepackage{colortbl}
\usepackage{booktabs}
\usepackage{tabularx}
\usepackage{diagbox} % case d'en-tête diagonale des tableaux à double entrée
% Colonnes extensibles centrée / gauche / droite (C, L, R), souvent utilisées par les modèles
\newcolumntype{C}{>{\centering\arraybackslash}X}
\newcolumntype{L}{>{\raggedright\arraybackslash}X}
\newcolumntype{R}{>{\raggedleft\arraybackslash}X}
\usepackage{array}
\usepackage{multirow}
\usepackage{multicol}
\usepackage{siunitx}
% parse-numbers=false : accepte aussi \SI{2,5 \times 10^{-3}}{...}
\sisetup{locale=FR,output-decimal-marker={,},group-minimum-digits=4,parse-numbers=false}
\DeclareSIUnit\ohm{\ensuremath{\symup{\Omega}}} % U+2126 absent de la police
\DeclareSIUnit\division{div} % réglages d'oscilloscope (ms/div, V/div)
\DeclareSIUnit\tr{tr} % tours (vitesse de rotation, ex. \SI{1500}{\tr\per\minute})
\DeclareSIUnit\bit{bit}
\DeclareSIUnit\byte{o} % octet
\DeclareSIUnit\inch{po} % pouce
\DeclareSIUnit\ppp{ppp} % points par pouce (résolution d'image)
\usepackage{qrcode}
% Blocs de code source (PHP, C, HTML, JavaScript, SQL) — spécifique à la
% matière Informatique : listings + habillage aux couleurs OnBuch+ (voir le
% style \lstset ci-dessous, à côté des autres boîtes pédagogiques).
\usepackage{listings}
% \degree : commande fréquemment utilisée par les modèles pour un angle ou une
% température en dehors de \SI{}{} (non fournie par les paquets chargés).
\providecommand{\degree}{^{\circ}}
% \qed (fin de démonstration), utilisé par les modèles sans amsthm
\providecommand{\qed}{\hfill$\square$}
% \barre : double barre des valeurs interdites dans un tableau de variations (tkz-tab)
\providecommand{\barre}{\ensuremath{\|}}
% environnement proof (amsthm non chargé) utilisé par les modèles
\ifdefined\proof\else\newenvironment{proof}[1][Démonstration]{\par\noindent\textit{#1.}\ }{\qed\par}\fi
\usepackage{lastpage}
\usepackage{hyperref}
\hypersetup{hidelinks}

% ============================================================
% Palette « Pop » OnBuch+ — mêmes hex que la classe P (plus_ui.dart)
% ============================================================
\definecolor{poporange}{HTML}{F59321}
\definecolor{popdark}{HTML}{E07A0C}
\definecolor{popcream}{HTML}{FFF6E8}
\definecolor{popink}{HTML}{1C1714}
\definecolor{poppurple}{HTML}{7A5AE0}
\definecolor{popgreen}{HTML}{2E9E6A}
\definecolor{poppink}{HTML}{E0457B}
\definecolor{popblue}{HTML}{2F7DE1}
\definecolor{popgold}{HTML}{C98A12}
\definecolor{popmuted}{HTML}{8A7F76}
\definecolor{popline}{HTML}{EADFD2}
\definecolor{popgray}{HTML}{8A7F76}
\colorlet{popgrey}{popgray}
% Alias tolérants (le modèle nomme parfois les couleurs différemment)
\colorlet{popwhite}{white}
\colorlet{popblack}{popink}
\colorlet{popred}{poppink}
\colorlet{popyellow}{popgold}
% Teintes claires pour fonds de tableaux / zones
\colorlet{popblueL}{popblue!10!white}
\colorlet{poporangeL}{poporange!14!white}
\colorlet{popgreenL}{popgreen!12!white}
\colorlet{poppurpleL}{poppurple!12!white}
\colorlet{poppinkL}{poppink!12!white}
\colorlet{popgoldL}{popgold!14!white}

\pagecolor{popcream}

% ============================================================
% Typographie
% ============================================================
\defaultfontfeatures{Ligatures=TeX}
\setmainfont{PlusJakartaSans-Medium.ttf}[Path=../../../fonts/,BoldFont=PlusJakartaSans-Bold.ttf]
% Maths en sans-serif (Fira Math) avec chiffres et lettres droites de Plus
% Jakarta, pour que les formules se fondent dans le texte.
\usepackage[math-style=ISO,bold-style=ISO]{unicode-math}
\setmathfont{FiraMath-Regular.otf}[Path=../../../fonts/]
\setmathfont{PlusJakartaSans-Medium.ttf}[Path=../../../fonts/,range={up/num,up/{Latin,latin},\mathup}]
\usepackage{icomma} % virgule décimale sans espace en mode math
% Caractères mathématiques Unicode absents des polices de texte (Plus Jakarta,
% Archivo Black) : rendus par leur équivalent mathématique, en texte comme en
% formule — sinon ils s'affichent en carré vide (ex. « Arithmétique dans ℤ »).
\usepackage{newunicodechar}
\newunicodechar{ℝ}{\ensuremath{\mathbb{R}}}
\newunicodechar{ℤ}{\ensuremath{\mathbb{Z}}}
\newunicodechar{ℂ}{\ensuremath{\mathbb{C}}}
\newunicodechar{ℕ}{\ensuremath{\mathbb{N}}}
\newunicodechar{ℚ}{\ensuremath{\mathbb{Q}}}
\newunicodechar{𝔻}{\ensuremath{\mathbb{D}}}
\newunicodechar{α}{\ensuremath{\alpha}}
\newunicodechar{β}{\ensuremath{\beta}}
\newunicodechar{γ}{\ensuremath{\gamma}}
\newunicodechar{δ}{\ensuremath{\delta}}
\newunicodechar{ε}{\ensuremath{\varepsilon}}
\newunicodechar{θ}{\ensuremath{\theta}}
\newunicodechar{λ}{\ensuremath{\lambda}}
\newunicodechar{μ}{\ensuremath{\mu}}
\newunicodechar{σ}{\ensuremath{\sigma}}
\newunicodechar{τ}{\ensuremath{\tau}}
\newunicodechar{φ}{\ensuremath{\varphi}}
\newunicodechar{ψ}{\ensuremath{\psi}}
\newunicodechar{ω}{\ensuremath{\omega}}
\newunicodechar{Σ}{\ensuremath{\Sigma}}
% majuscules produites par \MakeUppercase dans les titres (α-aminés -> Α)
\newunicodechar{Α}{\ensuremath{\alpha}}
\newunicodechar{Β}{\ensuremath{\beta}}
\newunicodechar{Γ}{\ensuremath{\Gamma}}
\newunicodechar{Δ}{\ensuremath{\Delta}}
\newunicodechar{Ε}{\ensuremath{\varepsilon}}
\newunicodechar{Θ}{\ensuremath{\Theta}}
\newunicodechar{Λ}{\ensuremath{\Lambda}}
\newunicodechar{Μ}{\ensuremath{\mu}}
\newunicodechar{Τ}{\ensuremath{\tau}}
\newunicodechar{Φ}{\ensuremath{\Phi}}
\newunicodechar{Ψ}{\ensuremath{\Psi}}
\newunicodechar{Ω}{\ensuremath{\Omega}}
\newunicodechar{↦}{\ensuremath{\mapsto}}
\newunicodechar{∈}{\ensuremath{\in}}
\newunicodechar{∉}{\ensuremath{\notin}}
\newunicodechar{∧}{\ensuremath{\wedge}}
\newunicodechar{⇔}{\ensuremath{\Leftrightarrow}}
\newunicodechar{⟺}{\ensuremath{\Longleftrightarrow}}
\newunicodechar{⇒}{\ensuremath{\Rightarrow}}
\newunicodechar{⟹}{\ensuremath{\Longrightarrow}}
\newunicodechar{∘}{\ensuremath{\circ}}
\newunicodechar{∪}{\ensuremath{\cup}}
\newunicodechar{∩}{\ensuremath{\cap}}
\newunicodechar{≡}{\ensuremath{\equiv}}
\newunicodechar{⊂}{\ensuremath{\subset}}
\newunicodechar{⊥}{\ensuremath{\perp}}
\newunicodechar{∃}{\ensuremath{\exists}}
\newunicodechar{∀}{\ensuremath{\forall}}
\newunicodechar{∛}{\ensuremath{\sqrt[3]{\,}}}
\newunicodechar{∜}{\ensuremath{\sqrt[4]{\,}}}
\newunicodechar{⃗}{\ensuremath{{}^{\to}}}
\newunicodechar{ⁿ}{\ifmmode^{n}\else\textsuperscript{\ensuremath{n}}\fi}
\newunicodechar{ˣ}{\ifmmode^{x}\else\textsuperscript{\ensuremath{x}}\fi}
\newunicodechar{ᵇ}{\ifmmode^{b}\else\textsuperscript{\ensuremath{b}}\fi}
\newunicodechar{ʳ}{\ifmmode^{r}\else\textsuperscript{\ensuremath{r}}\fi}
\newunicodechar{ᵘ}{\ifmmode^{u}\else\textsuperscript{\ensuremath{u}}\fi}
\newunicodechar{ʸ}{\ifmmode^{y}\else\textsuperscript{\ensuremath{y}}\fi}
\newunicodechar{ᵃ}{\ifmmode^{a}\else\textsuperscript{\ensuremath{a}}\fi}
\newunicodechar{ᵉ}{\ifmmode^{e}\else\textsuperscript{\ensuremath{e}}\fi}
\newunicodechar{ᵏ}{\ifmmode^{k}\else\textsuperscript{\ensuremath{k}}\fi}
\newunicodechar{ᵖ}{\ifmmode^{p}\else\textsuperscript{\ensuremath{p}}\fi}
\newunicodechar{ᴺ}{\ifmmode^{N}\else\textsuperscript{\ensuremath{N}}\fi}
\newunicodechar{ᶜ}{\ifmmode^{c}\else\textsuperscript{\ensuremath{c}}\fi}
\newunicodechar{ʰ}{\ifmmode^{h}\else\textsuperscript{\ensuremath{h}}\fi}
\newunicodechar{ᵠ}{\ifmmode^{q}\else\textsuperscript{\ensuremath{q}}\fi}
\newunicodechar{ₙ}{\ifmmode_{n}\else\textsubscript{\ensuremath{n}}\fi}
\newunicodechar{ₐ}{\ifmmode_{a}\else\textsubscript{\ensuremath{a}}\fi}
\newunicodechar{ᵢ}{\ifmmode_{i}\else\textsubscript{\ensuremath{i}}\fi}
\newunicodechar{ᵣ}{\ifmmode_{r}\else\textsubscript{\ensuremath{r}}\fi}
\newunicodechar{ₖ}{\ifmmode_{k}\else\textsubscript{\ensuremath{k}}\fi}
\newunicodechar{ᵦ}{\ifmmode_{\beta}\else\textsubscript{\ensuremath{\beta}}\fi}
\newunicodechar{ₓ}{\ifmmode_{x}\else\textsubscript{\ensuremath{x}}\fi}
\newfontfamily\popdisplay{ArchivoBlack-Regular.ttf}[Path=../../../fonts/]
\newfontfamily\poplogo{SpaceGrotesk-Bold.ttf}[Path=../../../fonts/]
\newfontfamily\popsemi{PlusJakartaSans-SemiBold.ttf}[Path=../../../fonts/]
\newfontfamily\popxbold{PlusJakartaSans-ExtraBold.ttf}[Path=../../../fonts/]
\newfontfamily\popxboldls{PlusJakartaSans-ExtraBold.ttf}[Path=../../../fonts/,LetterSpace=3.0]

\setlist[enumerate]{leftmargin=1.7em,itemsep=5pt,topsep=4pt}
\setlist[itemize]{leftmargin=1.7em,itemsep=4pt,topsep=4pt,label={\color{poporange}\textbullet}}
\setlength{\parindent}{0pt}
\setlength{\parskip}{5pt}
\renewcommand{\arraystretch}{1.25}

% pgfplots : style Pop par défaut
\pgfplotsset{
  every axis/.append style={axis line style={popink,line width=1pt},tick style={popink},
    grid style={popline,line width=0.5pt},label style={font=\small},tick label style={font=\footnotesize}},
  popcurve/.style={poporange,line width=1.8pt,no markers,smooth},
}
% Même style disponible dans un \draw TikZ simple (hors pgfplots)
\tikzset{popcurve/.style={poporange,line width=1.8pt}}
\tikzset{popgrid/.style={popline,very thin}}
\colorlet{popcurve}{poporange} % le nom du style est parfois utilisé comme couleur

% ============================================================
% Pill / squiggle
% ============================================================
\newcommand{\poppill}[3]{%
  \tikz[baseline=-0.55ex]\node[draw=popink,line width=1.2pt,rounded corners=2.6pt,%
    fill=#2,inner xsep=6.5pt,inner ysep=3pt]{%
    \popxboldls\fontsize{7}{7}\selectfont\color{#3}\MakeUppercase{#1}};%
}
\newcommand{\popsquiggle}[1]{%
  \tikz[baseline=-0.4ex]{
    \draw[#1,line width=2.6pt,line cap=round]
      plot [smooth] coordinates {(0,0.09) (0.34,0.01) (0.68,0.21) (1.02,0.01) (1.36,0.10)};
  }%
}
% Mot-clé surligné (marqueur orange sous le texte)
\newcommand{\cle}[1]{{\popxbold\color{popdark}#1}}

% ============================================================
% En-tête / pied
% ============================================================
\pagestyle{fancy}
\fancyhf{}
\renewcommand{\headrulewidth}{0pt}
\renewcommand{\footrulewidth}{0pt}
\lhead{\raisebox{-0.15ex}{\poplogo\fontsize{13}{13}\selectfont\color{popink}OnBuch\color{poporange}+}}
\rhead{\poppill{\DOCMATIERE\ \textbullet\ \DOCNIVEAU\ \textbullet\ Leçon \DOCLECON}{white}{popink}}
\lfoot{\popxbold\fontsize{7.6}{7.6}\selectfont\color{popmuted}\MakeUppercase{Cours signé OnBuch+}\ \textbullet\ {\popsemi\DOCTITRE}}
\rfoot{\tikz[baseline=-0.6ex]\node[rounded corners=8.5pt,fill=popink,minimum height=17pt,minimum width=17pt,inner xsep=5pt,inner sep=0]{\popxbold\fontsize{8}{8}\selectfont\color{popcream}\thepage\,/\,\pageref*{LastPage}};}

% ============================================================
% Titres
% ============================================================
\newcommand{\chaptitre}[2]{%
  \begin{tcolorbox}[enhanced,colback=poporange,colframe=popink,boxrule=2.6pt,arc=11pt,
      drop shadow=popink,left=15pt,right=15pt,top=12pt,bottom=11pt,before skip=3pt,after skip=18pt]
    \poppill{#2}{popink}{popcream}\par\vspace{9pt}
    {\raggedright\hyphenpenalty=10000\exhyphenpenalty=10000\popdisplay\fontsize{22}{24}\selectfont\color{popcream}\MakeUppercase{#1}\par}
  \end{tcolorbox}
}
% Section numérotée : pastille orange avec le numéro + titre Archivo
\newcounter{popsec}
\newcommand{\coursec}[1]{%
  \stepcounter{popsec}\setcounter{popsub}{0}%
  \par\vspace{16pt}\noindent
  \begin{minipage}{\linewidth}
  \tikz[baseline=-0.7ex]\node[draw=popink,line width=1.6pt,fill=poporange,rounded corners=4pt,minimum size=22pt,inner sep=0]{\popdisplay\fontsize{12}{12}\selectfont\color{popcream}\thepopsec};%
  \hspace{8pt}{\popdisplay\fontsize{15}{17}\selectfont\color{popink}\MakeUppercase{#1}}\par
  \vspace{4pt}\noindent{\color{poporange}\rule{36pt}{3.6pt}}
  \end{minipage}\par\nopagebreak\vspace{8pt}
}
\newcounter{popsub}
\newcommand{\courssub}[1]{%
  \stepcounter{popsub}%
  \par\vspace{9pt}\noindent{\popxbold\fontsize{11.5}{13}\selectfont\color{popink}{\color{poporange}\thepopsec.\thepopsub}\hspace{6pt}#1}\par\nopagebreak\vspace{4pt}
}

% ============================================================
% Boîtes pédagogiques — même recette : fond blanc, bordure épaisse colorée,
% ombre dure encre, badge plein. Titre optionnel [#1].
% ============================================================
\tcbset{popbase/.style={breakable,enhanced,colback=white,boxrule=2.3pt,arc=9pt,
  shadow={3pt}{-3pt}{0pt}{popink},left=14pt,right=14pt,top=11pt,bottom=11pt,before skip=11pt,after skip=15pt,
  fontupper=\popsemi\fontsize{10.6}{14.5}\selectfont\color{popink},
  fontlower=\popsemi\fontsize{10.6}{14.5}\selectfont\color{popink}}}
% Coche dessinée (le glyphe ✓ n'existe pas dans les polices du document)
\newcommand{\popcheck}[1]{\tikz[baseline=-0.5ex]\draw[#1,line width=1.6pt,line cap=round,line join=round](-0.13,0.0)--(-0.04,-0.09)--(0.13,0.1);}
\providecommand{\coche}{\popcheck{popgreen}}
\providecommand{\resultat}[1]{\par\smallskip\noindent\fcolorbox{popgreen}{popgreenL}{\parbox{\dimexpr\linewidth-2\fboxsep-2\fboxrule\relax}{\textbf{Résultat.} #1}}\par\smallskip}
\newcommand{\popboxhead}[3]{\poppill{#1}{#2}{white}\ifx\relax#3\relax\else\hspace{6pt}{\popxbold\fontsize{10.6}{12}\selectfont\color{popink}#3}\fi\par\vspace{7pt}}

\newtcolorbox{aretenir}[1][]{popbase,colframe=popgreen,before upper={\popboxhead{À retenir}{popgreen}{#1}}}
\newtcolorbox{exemplebox}[1][]{popbase,colframe=poporange,before upper={\popboxhead{Exemple}{poporange}{#1}}}
\newtcolorbox{definition}[1][]{popbase,colframe=popblue,before upper={\popboxhead{Définition}{popblue}{#1}}}
\newtcolorbox{propriete}[1][]{popbase,colframe=poppurple,before upper={\popboxhead{Propriété}{poppurple}{#1}}}
\newtcolorbox{methode}[1][]{popbase,colframe=popgold,before upper={\popboxhead{Méthode}{popgold}{#1}}}
\newtcolorbox{attention}[1][]{popbase,colframe=poppink,before upper={\popboxhead{Attention}{poppink}{#1}}}
\newtcolorbox{experience}[1][]{popbase,colframe=popblue,colback=popblueL,before upper={\popboxhead{Expérience}{popblue}{#1}}}
% « Le savais-tu ? » : carte violette pleine (contexte, culture, Cameroun)
\newtcolorbox{savaistu}[1][]{popbase,colback=poppurple,colframe=popink,
  fontupper=\popsemi\fontsize{10.4}{14.2}\selectfont\color{white},
  before upper={\poppill{Le savais-tu\,?}{white}{poppurple}\ifx\relax#1\relax\else\hspace{6pt}{\popxbold\color{white}#1}\fi\par\vspace{7pt}}}
% Exercice résolu : énoncé en haut, \tcblower puis la solution sur fond crème
\newcounter{popexo}\newcounter{popexr}
\newtcolorbox{exoresolu}[1][]{popbase,colframe=poporange,colbacklower=popcream,
  segmentation style={popink,line width=1.2pt,dashed},
  before upper={\stepcounter{popexr}\popboxhead{Exercice résolu \thepopexr}{poporange}{#1}},
  before lower={\poppill{Solution}{popink}{popcream}\par\vspace{6pt}}}
% Exercice (non résolu en place) — la correction est dans la partie Corrigés
\newtcolorbox{exercice}[1][]{popbase,colframe=popink,
  before upper={\stepcounter{popexo}\popboxhead{Exercice \thepopexo}{popink}{#1}}}
% Corrigé
\newtcolorbox{corrige}[1][]{popbase,colframe=popgreen,colback=popgreenL,
  before upper={\popboxhead{Corrigé}{popgreen}{#1}}}
% Objectifs / prérequis / bilan
\newtcolorbox{objectifs}{popbase,colframe=popink,colback=poporangeL,
  before upper={\popboxhead{Objectifs de la leçon}{popink}{}}}
\newtcolorbox{prerequis}{popbase,colframe=popink,colback=popblueL,
  before upper={\popboxhead{Prérequis}{popblue}{}}}
\newtcolorbox{bilan}{popbase,colframe=popink,colback=white,boxrule=2.8pt,
  before upper={\popboxhead{Fiche bilan}{poporange}{L'essentiel à connaître pour l'examen}}}
% Figure encadrée avec légende
\newtcolorbox{popfigure}[1][]{enhanced,breakable=false,colback=white,colframe=popline,boxrule=1.4pt,arc=8pt,
  left=10pt,right=10pt,top=10pt,bottom=8pt,before skip=10pt,after skip=12pt,halign=center,
  lower separated=false,halign lower=center,
  fontlower=\popsemi\fontsize{9}{11.5}\selectfont\color{popmuted},
  #1}
\newcommand{\legende}[1]{\par\vspace{6pt}{\popsemi\fontsize{9}{11.5}\selectfont\color{popmuted}#1}}

% ============================================================
% Bloc de code source — \begin{lstlisting}[language=Php]...\end{lstlisting}
% (ou language=C, HTML, JavaScript, SQL ; option omise = pseudo-code brut).
% Même recette visuelle que les figures (\popfigure) : cadre encre épais,
% fond clair (popcream), légende possible juste après via \legende{...}
% comme pour une figure (listings ne sait pas arrondir ses coins : on garde
% un cadre rectangulaire, cohérent avec les autres tableaux du document).
% CONTRAT : les modèles ne doivent utiliser QUE \begin{lstlisting}[...] tel
% quel (pas de \begin{tcblisting}, pas de \verb pour un bloc multi-lignes).
% ============================================================
\lstdefinestyle{poplst}{
  backgroundcolor=\color{popcream},
  basicstyle=\popsemi\fontsize{8.6}{11}\selectfont\color{popink},
  keywordstyle=\bfseries\color{popblue},
  commentstyle=\itshape\color{popmuted},
  stringstyle=\color{popgreen},
  numberstyle=\tiny\color{popmuted},
  identifierstyle=\color{popink},
  frame=l,
  rulecolor=\color{popink},
  framerule=1.6pt,
  framesep=7pt,
  framexleftmargin=7pt,
  framexrightmargin=7pt,
  xleftmargin=2pt,
  xrightmargin=2pt,
  breaklines=true,
  breakatwhitespace=false,
  showstringspaces=false,
  showspaces=false,
  showtabs=false,
  tabsize=2,
  columns=flexible,
  keepspaces=true,
  numbers=none,
  aboveskip=10pt,
  belowskip=10pt,
  captionpos=b,
}
\lstset{style=poplst, language=PHP}
% La distribution TeX embarquée ne fournit pas le fichier de langue
% JavaScript de listings (lstlang*.dat) : on la redéfinit nous-mêmes
% pour que \begin{lstlisting}[language=JavaScript] fonctionne.
\lstdefinelanguage{JavaScript}{
  keywords={break,case,catch,class,const,continue,debugger,default,delete,do,
    else,export,extends,finally,for,function,if,import,in,instanceof,let,new,
    return,static,super,switch,this,throw,try,typeof,var,void,while,with,yield,
    async,await,of,get,set},
  keywordstyle=\bfseries\color{popblue},
  ndkeywords={true,false,null,undefined,NaN,Infinity},
  ndkeywordstyle=\bfseries\color{poporange},
  identifierstyle=\color{popink},
  sensitive=true,
  comment=[l]{//},
  morecomment=[s]{/*}{*/},
  string=[b]",
  morestring=[b]',
  morestring=[b]`,
}
% Même limitation pour CSS et les fichiers de configuration (apacheconf,
% nginx, ini...) : ils ne sont pas fournis. CSS et un style générique de
% type "config" (commentaires # ou ;, pas de mots-clés) couvrent les besoins.
\lstdefinelanguage{CSS}{
  keywords={color,background,background-color,margin,padding,border,
    border-radius,font,font-size,font-weight,font-family,width,height,
    display,position,top,left,right,bottom,float,clear,text-align,
    line-height,list-style,cursor,overflow,box-shadow,transition,
    flex,grid,align-items,justify-content,z-index},
  keywordstyle=\bfseries\color{popblue},
  identifierstyle=\color{popink},
  sensitive=false,
  comment=[s]{/*}{*/},
  string=[b]",
  morestring=[b]',
}
\lstdefinelanguage{apacheconf}{
  sensitive=false,
  morecomment=[l]{\#},
}
\lstdefinelanguage{Apache}{
  sensitive=false,
  morecomment=[l]{\#},
}
\lstdefinelanguage{text}{sensitive=false}
\lstdefinelanguage{powershell}{
  keywords={New-Object,New-SmbShare,Get-Acl,Set-Acl,Get-ChildItem,Set-Content,
    Get-Content,Write-Host,ForEach-Object,Where-Object,Select-Object,
    Test-Path,Remove-Item,Copy-Item,Move-Item,Start-Process,Stop-Process,
    If,Else,ElseIf,ForEach,While,Function,Return,Try,Catch},
  keywordstyle=\bfseries\color{popblue},
  identifierstyle=\color{popink},
  sensitive=false,
  comment=[l]{\#},
  string=[b]",
  morestring=[b]',
}
\lstdefinelanguage{ini}{
  sensitive=false,
  morecomment=[l]{;},
  morecomment=[l]{\#},
}

% ============================================================
% Signature OnBuch+
% ============================================================
\newcommand{\poplogoinline}[1]{{\poplogo\fontsize{#1}{#1}\selectfont\color{popink}OnBuch\color{poporange}+}}
\newcommand{\signatureonbuch}{%
  \begin{tcolorbox}[enhanced,colback=popink,colframe=popink,boxrule=2.2pt,arc=10pt,
      drop shadow=poporange,left=14pt,right=14pt,top=10pt,bottom=10pt,before skip=0pt,after skip=0pt]
    \tikz[baseline=-0.7ex]\node[circle,fill=popgreen,draw=popcream,line width=1.3pt,minimum size=22pt,inner sep=0]{\popcheck{white}};%
    \hspace{9pt}\begin{minipage}[c]{0.62\linewidth}
      {\popxbold\fontsize{12}{13}\selectfont\color{popcream}Signé\ \textbullet\ {\poplogo OnBuch\color{poporange}+}}\par\vspace{2pt}
      {\raggedright\popsemi\fontsize{8}{10.5}\selectfont\color{popline}Ludovic A.\\\MakeUppercase{Conforme au programme officiel MINESEC}\par}
    \end{minipage}\hfill
    {\poplogo\fontsize{22}{22}\selectfont\color{popcream}OnBuch\color{poporange}+}
  \end{tcolorbox}%
}

% ============================================================
% Page de garde
% #1 titre · #2 matière · #3 niveau · #4 module · #5 n° leçon · #6 durée ·
% #7 description / objectifs (texte ou itemize)
% ============================================================
\newcommand{\pagegarde}[7]{%
  \thispagestyle{empty}
  \newgeometry{a4paper,margin=1.7cm,top=1.6cm,bottom=1.4cm}%
  \begin{tikzpicture}[remember picture,overlay]
    \fill[poporange!18] ([xshift=-3.2cm,yshift=-3.6cm]current page.north east) circle (4.6cm);
    \fill[poppurple!14] ([xshift=2.4cm,yshift=7.6cm]current page.south west) circle (3.8cm);
  \end{tikzpicture}%
  {\poplogo\fontsize{30}{30}\selectfont\color{popink}OnBuch\color{poporange}+}\hfill
  \raisebox{0.6ex}{\poppill{Cours signé \textbullet\ Premium}{popink}{popcream}}\par
  \vspace{1pt}\popsquiggle{poppurple}\par
  \vspace{8pt}
  \begin{tcolorbox}[enhanced,colback=poporange,colframe=popink,boxrule=2.8pt,arc=13pt,
      drop shadow=popink,left=18pt,right=18pt,top=12pt,bottom=13pt,after skip=10pt]
    \poppill{#2 \textbullet\ #3}{popink}{popcream}\hspace{5pt}\poppill{#4}{white}{popink}\par\vspace{8pt}
    {\popdisplay\fontsize{14}{14}\selectfont\color{popink}LEÇON #5}\par\vspace{4pt}
    {\raggedright\hyphenpenalty=10000\exhyphenpenalty=10000\popdisplay\fontsize{25}{27}\selectfont\color{popcream}\MakeUppercase{#1}\par}
    \vspace{8pt}
    {\popxbold\fontsize{9.5}{9.5}\selectfont\color{popink}\MakeUppercase{Durée conseillée : #6}}
  \end{tcolorbox}
  \begin{tcolorbox}[enhanced,colback=white,colframe=popink,boxrule=2.2pt,arc=10pt,
      drop shadow=popink,left=16pt,right=16pt,top=10pt,bottom=10pt,after skip=8pt]
    \poppill{Dans cette leçon}{poppurple}{white}\par\vspace{5pt}
    \popsemi\fontsize{9.3}{12.6}\selectfont\color{popink}#7
  \end{tcolorbox}
  \noindent\begin{minipage}[t]{0.487\linewidth}
    \begin{tcolorbox}[enhanced,colback=popgreen,colframe=popink,boxrule=2.2pt,arc=10pt,
        drop shadow=popink,left=11pt,right=11pt,top=10pt,bottom=10pt,height=3.1cm,valign=center]
      \tcbox[colback=white,colframe=popink,boxrule=1.3pt,arc=5pt,boxsep=0pt,nobeforeafter,
        left=3pt,right=3pt,top=3pt,bottom=3pt,valign=center]{\qrcode[height=1.9cm]{https://play.google.com/store/apps/details?id=com.luvvix.onbuch&hl=fr}}%
      \hspace{8pt}\begin{minipage}[c]{0.5\linewidth}
        {\popdisplay\fontsize{11}{12.5}\selectfont\color{white}\MakeUppercase{Télécharge OnBuch}}\par\vspace{4pt}
        {\raggedright\popsemi\fontsize{8.4}{11}\selectfont\color{white}Gratuit \textbullet\ Google Play\\Tuteur IA, annales, concours, résultats\par}
      \end{minipage}
    \end{tcolorbox}
  \end{minipage}\hfill
  \begin{minipage}[t]{0.487\linewidth}
    \begin{tcolorbox}[enhanced,colback=poppurple,colframe=popink,boxrule=2.2pt,arc=10pt,
        drop shadow=popink,left=11pt,right=11pt,top=10pt,bottom=10pt,height=3.1cm,valign=center]
      \tikz[baseline=-0.7ex]\node[circle,fill=white,draw=popink,line width=1.3pt,minimum size=1.6cm,inner sep=0]{\popdisplay\fontsize{24}{24}\selectfont\color{poppurple}+};%
      \hspace{8pt}\begin{minipage}[c]{0.6\linewidth}
        {\popdisplay\fontsize{11}{12.5}\selectfont\color{white}\MakeUppercase{Passe à OnBuch+}}\par\vspace{4pt}
        {\raggedright\popsemi\fontsize{8.4}{11}\selectfont\color{white}Tous les cours signés, Tuteur IA illimité, fiches PDF — onglet OnBuch+ de l'app\par}
      \end{minipage}
    \end{tcolorbox}
  \end{minipage}
  \vspace{8pt}
  \popcontenu
  \vfill
  \signatureonbuch
  \restoregeometry
  \newpage
}

% Rangée « Ce cours contient » (page de garde)
\newcommand{\poptile}[3]{% #1 couleur #2 gros texte #3 légende
  \begin{tcolorbox}[enhanced,colback=#1,colframe=popink,boxrule=2pt,arc=9pt,shadow={2.5pt}{-2.5pt}{0pt}{popink},
      left=6pt,right=6pt,top=9pt,bottom=9pt,halign=center,valign=center,height=1.75cm,width=\linewidth,nobeforeafter]
    {\popdisplay\fontsize{12.5}{13}\selectfont\color{white}\MakeUppercase{#2}}\par\vspace{3pt}
    {\centering\popsemi\fontsize{7.8}{9.5}\selectfont\color{white}#3\par}
  \end{tcolorbox}}
\newcommand{\popcontenu}{%
  \par\noindent{\popxboldls\fontsize{7.5}{7.5}\selectfont\color{popmuted}CE COURS SIGNÉ CONTIENT}\par\vspace{7pt}
  \noindent\begin{minipage}[t]{0.235\linewidth}\poptile{popblue}{Cours}{complet, illustré, pas à pas}\end{minipage}\hfill
  \begin{minipage}[t]{0.235\linewidth}\poptile{poporange}{Méthodes}{et exercices résolus}\end{minipage}\hfill
  \begin{minipage}[t]{0.235\linewidth}\poptile{poppink}{Exercices}{type examen + corrigés}\end{minipage}\hfill
  \begin{minipage}[t]{0.235\linewidth}\poptile{popgreen}{Fiche bilan}{l'essentiel à retenir}\end{minipage}\par}

% Sommaire
\newtcolorbox{sommaire}{enhanced,colback=white,colframe=popink,boxrule=2.2pt,arc=10pt,
  drop shadow=popink,left=18pt,right=18pt,top=13pt,bottom=13pt,before skip=6pt,after skip=20pt,
  before upper={\poppill{Sommaire}{poppurple}{white}\par\vspace{9pt}\popsemi\fontsize{10.6}{14.5}\selectfont\color{popink}}}

% Signature de fin de cours — poussée tout en bas de la dernière page
\newcommand{\findecours}{%
  \par\vfill\nopagebreak
  \begin{center}\popsquiggle{poporange}\end{center}\vspace{4pt}
  \signatureonbuch
}
\AtBeginDocument{\def\llbracket{\mathopen{[\mkern-3.5mu[}}\def\rrbracket{\mathclose{]\mkern-3.5mu]}}} % intervalles d'entiers (glyphe ⟦ absent de la police)
% Glyphes absents de Fira Math : redéfinis à partir de glyphes disponibles
\AtBeginDocument{%
  \def\setminus{\mathbin{\backslash}}\let\smallsetminus\setminus
  \def\vdots{\mathord{\vbox{\baselineskip3.2pt\lineskiplimit0pt\kern4pt\hbox{.}\hbox{.}\hbox{.}}}}%
  \def\ddots{\mathinner{\mkern1mu\raise6.5pt\hbox{.}\mkern2mu\raise3.6pt\hbox{.}\mkern2mu\raise0.7pt\hbox{.}\mkern1mu}}%
  \def\triangle{\mathord{\scalebox{0.9}{$\symup{\Delta}$}}}%
  \def\nabla{\mathord{\raisebox{\depth}{\scalebox{1}[-1]{$\symup{\Delta}$}}}}%
  \def\checkmark{\popcheck{popdark}}%
  \def\star{\mathbin{\ast}}\def\bigstar{\mathord{\ast}}%
  \def\diamond{\mathbin{\scalebox{0.6}{$\Diamond$}}}%
  \def\bigcup{\mathop{\mathchoice{\raisebox{-0.3em}{\scalebox{1.7}{$\cup$}}}{\scalebox{1.25}{$\cup$}}{\cup}{\cup}}\displaylimits}%
  \def\bigcap{\mathop{\mathchoice{\raisebox{-0.3em}{\scalebox{1.7}{$\cap$}}}{\scalebox{1.25}{$\cap$}}{\cap}{\cap}}\displaylimits}%
  \def\lesseqqgtr{\mathrel{\lessgtr}}\def\bigoplus{\mathop{\oplus}\displaylimits}%
}
\providecommand{\ssi}{si et seulement si} % abréviation fréquente
\AtBeginDocument{\providecommand{\wideparen}{\overparen}} % arc de cercle

\begin{document}

\pagegarde{Utilisation d’un logiciel de publication}{Informatique}{3ème}{Informatique – classe de 3ème}{N08}{2 séances (fiche de progression harmonisée)}{Cette leçon initie les élèves à la Publication Assistée par Ordinateur (PAO) : découverte de l'interface d'un logiciel de publication (Microsoft Publisher ou équivalent libre) et production concrète d'un document de communication (affiche, dépliant, carte de visite). Elle développe l'autonomie numérique des élèves et leur capacité à justifier leurs choix de mise en page.
\vspace{4pt}
\begin{multicols}{2}\small
\begin{itemize}
\item À la fin de cette leçon, je sais définir la PAO et citer au moins deux logiciels de publication
\item À la fin de cette leçon, je sais distinguer un logiciel de PAO d'un traitement de texte
\item À la fin de cette leçon, je sais identifier et nommer les zones principales de l'interface d'un logiciel de PAO (ruban, zone de travail, repères, pages)
\item À la fin de cette leçon, je sais créer une publication à partir d'un modèle ou d'une page vierge
\item À la fin de cette leçon, je sais insérer et mettre en forme des zones de texte, des images et des formes dans une publication
\item À la fin de cette leçon, je sais utiliser les repères, les marges et les colonnes pour organiser une page
\end{itemize}\end{multicols}}

\begin{sommaire}
\begin{multicols}{2}
\begin{enumerate}
\item Introduction à la PAO
\item Lancement et interface du logiciel de PAO
\item Création d'une publication : page et structure
\item Insertion et mise en forme des objets
\item Finalisation, exportation et impression
\item Démarche de production et justification
\item Activité d'intégration
\item Méthodes et exercices résolus
\item Exercices
\item Corrigés des exercices
\item Fiche bilan
\end{enumerate}
\end{multicols}
\end{sommaire}

\begin{objectifs}
\begin{itemize}
\item À la fin de cette leçon, je sais définir la PAO et citer au moins deux logiciels de publication
\item À la fin de cette leçon, je sais distinguer un logiciel de PAO d'un traitement de texte
\item À la fin de cette leçon, je sais identifier et nommer les zones principales de l'interface d'un logiciel de PAO (ruban, zone de travail, repères, pages)
\item À la fin de cette leçon, je sais créer une publication à partir d'un modèle ou d'une page vierge
\item À la fin de cette leçon, je sais insérer et mettre en forme des zones de texte, des images et des formes dans une publication
\item À la fin de cette leçon, je sais utiliser les repères, les marges et les colonnes pour organiser une page
\item À la fin de cette leçon, je sais enregistrer, exporter en PDF et imprimer une publication
\item À la fin de cette leçon, je sais rédiger et justifier la démarche suivie pour produire un objet de PAO adapté à un besoin concret
\end{itemize}
\end{objectifs}

\begin{prerequis}
\begin{itemize}
\item Notions de base sur l'environnement Windows (fenêtres, fichiers, dossiers) vues en 4ème
\item Utilisation élémentaire d'un traitement de texte (saisie, mise en forme simple) vue en 4ème
\item Insertion d'images dans un document (leçon sur le traitement de texte, début d'année de 3ème)
\item Notions sur les périphériques d'entrée/sortie (clavier, souris, imprimante)
\end{itemize}
\end{prerequis}

\newpage

\coursec{Introduction à la PAO}

L'association des jeunes de ton quartier à Yaoundé organise un tournoi de football inter-quartiers et te confie la réalisation d'une affiche attrayante ainsi que d'un dépliant présentant le programme. Tu ouvres ton traitement de texte habituel, tu insères une photo du stade, tu écris le titre en grand, tu ajoutes des colonnes pour les horaires... et c'est le drame : l'image saute à la page suivante quand tu appuies sur \texttt{Entrée}, les colonnes se déforment, le texte ne contourne pas la photo comme tu le veux. Le résultat est terne, rigide, loin de l'affiche professionnelle que tu imagines. \textbf{Comment produire un document combinant textes, images et couleurs avec une précision millimétrique, sans que les éléments ne « bougent » au moindre ajout de ligne ?} C'est tout l'objet de la \textbf{Publication Assistée par Ordinateur (PAO)}.

\courssub{Qu'est-ce que la Publication Assistée par Ordinateur (PAO) ?}

\begin{definition}[Publication Assistée par Ordinateur (PAO)]
La \textbf{Publication Assistée par Ordinateur (PAO)} désigne l'ensemble des techniques informatiques permettant de concevoir, de composer et de mettre en page des documents destinés à être imprimés ou diffusés numériquement. Contrairement à la simple saisie de texte, la PAO offre un contrôle total sur la \textbf{typographie}, la \textbf{géométrie de la page} (marges, colonnes, gabarits), la \textbf{gestion des couleurs} (quadrichromie CMJN, tons directs Pantone) et le \textbf{positionnement absolu} des blocs textuels et graphiques.
\end{definition}

En termes simples : la PAO, c'est l'art de \textbf{maquetter} une page comme un architecte dessine un plan. Tu ne « tapes » pas seulement du texte ; tu \textbf{places} des objets (cadres de texte, images, formes, tableaux) aux coordonnées exactes $(x, y)$ sur une page de dimensions définies (A4, A5, A3, format personnalisé). Le logiciel de PAO gère la \textbf{résolution} des images (300 ppp pour l'impression), les \textbf{fonds perdus} (zone excédentaire rognée à la coupe) et les \textbf{repères de coupe} indispensables à l'imprimeur.

\courssub{La PAO dans notre quotidien camerounais}

Tu croises des objets de PAO tous les jours sans forcément le savoir. Voici des exemples concrets au Cameroun :

\begin{itemize}
    \item \textbf{Affiches de campagne électorale ou événementielle} : collées sur les poteaux électriques à Carrefour Warda ou au Rond-Point Deïdo, elles combinent photos du candidat, slogans en gros caractères, logos de parti et mentions légales en petit bas de page.
    \item \textbf{Menus de restaurant} : au \emph{Chalet 54} ou chez \emph{Le Bœuf}, le menu plastifié A4 plié en deux ou trois volets, avec photos des plats, prix alignés à la virgule, en-tête au logo de l'établissement.
    \item \textbf{Faire-part de mariage, baptême, décès} : format carte (10×15 cm) ou DL (99×210 mm), papier texturé, dorure à chaud, police calligraphiée pour les noms, police lisible pour l'adresse et l'heure à la mairie de Nkolndongo ou à l'église Saint-Jean-Baptiste de Mvolyé.
    \item \textbf{Programmes de cérémonies officielles} : fête de la jeunesse du 11 février, remise des diplômes à l'Université de Yaoundé I, avec emploi du temps en tableau, photos des lauréats, armoiries de la République.
    \item \textbf{Journaux scolaires et bulletins paroissiaux} : \emph{Le Lycéen} au Lycée Leclerc, \emph{L'Écho de la Paroisse} à Saint-Paul de Nsam, mis en page en colonnes, avec titraille, chapô, corps de texte, légendes de photos, pagination.
    \item \textbf{Cartes de visite} : format standard 85×55 mm, recto/verso, logo vectoriel, coordonnées alignées sur grille de base, coins arrondis ou carrés, vernis sélectif.
\end{itemize}

Tous ces documents partagent une caractéristique : \textbf{la mise en page prime sur la saisie}. Le contenu est « versé » dans une structure visuelle préétablie (le \emph{gabarit} ou \emph{maquette}).

\courssub{PAO vs Traitement de texte : deux logiques différentes}

C'est la distinction fondamentale pour ne plus jamais galérer avec une image qui « saute » au mauvais endroit.

\begin{propriete}[Différence structurelle : Flux continu vs Blocs positionnés]
\begin{itemize}
    \item \textbf{Traitement de texte (Word, Writer, Google Docs) : logique de \emph{flux continu} (flow).} Le texte s'écrit de gauche à droite, de haut en bas, page après page. Les images sont des « caractères géants » ancrés au paragraphe : si tu insères une ligne avant, tout le contenu suivant se décale. Idéal pour : rapports, dissertations, lettres, romans, comptes rendus.
    \item \textbf{PAO (Publisher, Scribus, InDesign, Canva) : logique de \emph{blocs indépendants} (frames/objets).} La page est une toile vierge. Tu crées un \textbf{cadre de texte} ici, un \textbf{cadre image} là, une \textbf{forme vectorielle} ailleurs. Chaque objet a ses coordonnées $(x, y)$, sa largeur, sa hauteur, son ordre de superposition (calque). Tu peux faire « couler » le texte d'un cadre à l'autre (chaînage), contourner une image de forme irrégulière (contour de détourage), superposer une transparence sur une photo. Idéal pour : affiches, dépliants, magazines, cartes de visite, packaging, signalétique.
\end{itemize}
\end{propriete}

\begin{popfigure}
\begin{tikzpicture}[
    scale=0.85,
    every node/.style={font=\small},
    frame/.style={draw=popdark, thick, rounded corners=3pt, fill=popcream},
    arrow/.style={->, >=Stealth, thick, popblue},
    labelstyle/.style={font=\bfseries\large, popdark, anchor=north}
]
% --- PAGE TRAITEMENT DE TEXTE (GAUCHE) ---
\begin{scope}[xshift=0cm]
    % Page frame
    \draw[frame] (0,0) rectangle (7,10);
    \node[labelstyle] at (3.5, 10.3) {Traitement de texte (Flux continu)};
    % Text flow representation
    \foreach \y in {9.2, 8.5, 7.8, 7.1, 6.4, 5.7, 5.0, 4.3, 3.6, 2.9, 2.2, 1.5, 0.8} {
        \draw[popmuted, line width=1.2pt] (0.8, \y) -- (6.2, \y);
    }
    % Paragraph breaks
    \draw[poporange, line width=2pt] (0.8, 5.7) -- (6.2, 5.7);
    \draw[poporange, line width=2pt] (0.8, 2.9) -- (6.2, 2.9);
    % Image "anchored" in flow - moves with text
    \draw[fill=popblueL, draw=popblue, thick] (1.5, 5.0) rectangle (5.5, 3.2);
    \node[popdark] at (3.5, 4.1) {Image ancrée au paragraphe};
    \draw[poporange, dashed, thick] (3.5, 3.2) -- (3.5, 2.9) node[midway, right, poporange, font=\tiny] {Si j'ajoute du texte avant...};
    \draw[poporange, dashed, thick, ->, >=Stealth] (3.5, 2.9) -- (3.5, 2.2);
    % Margin guides
    \draw[dashed, popmuted, thin] (0.8, 0) -- (0.8, 10);
    \draw[dashed, popmuted, thin] (6.2, 0) -- (6.2, 10);
\end{scope}
% --- PAGE PAO (DROITE) ---
\begin{scope}[xshift=9.5cm]
    % Page frame
    \draw[frame] (0,0) rectangle (7,10);
    \node[labelstyle] at (3.5, 10.3) {PAO (Blocs positionnés librement)};
    % Margin guides (fixed)
    \draw[dashed, popgreen, thin] (0.8, 0) -- (0.8, 10);
    \draw[dashed, popgreen, thin] (6.2, 0) -- (6.2, 10);
    \draw[dashed, popgreen, thin] (0, 9.2) -- (7, 9.2);
    \draw[dashed, popgreen, thin] (0, 0.8) -- (7, 0.8);
    % Block 1: Title
    \draw[fill=poppinkL, draw=poppink, thick] (0.8, 9.0) rectangle (6.2, 7.8);
    \node[popdark, align=center] at (3.5, 8.4) {Cadre Titre\\(Fixe en haut)};
    % Block 2: Image (free position)
    \draw[fill=popblueL, draw=popblue, thick] (0.8, 7.2) rectangle (3.8, 4.8);
    \node[popdark, align=center] at (2.3, 6.0) {Cadre Image\\(Position exacte x,y)};
    % Wrap text indicator
    \draw[popgreen, decorate, decoration={snake, amplitude=1pt, segment length=4pt}] (3.8, 6.5) -- (4.5, 6.5);
    \node[popgreen, font=\tiny, right] at (4.5, 6.5) {Contour de texte};
    % Block 3: Text column 1
    \draw[fill=popgreenL, draw=popgreen, thick] (4.2, 7.2) rectangle (6.2, 4.8);
    \node[popdark, align=center] at (5.2, 6.0) {Colonne 1\\(Texte chainé)};
    % Block 4: Text column 2
    \draw[fill=popgreenL, draw=popgreen, thick] (0.8, 4.2) rectangle (3.0, 1.5);
    \node[popdark, align=center] at (1.9, 2.85) {Colonne 2\\(Suite du texte)};
    % Block 5: Sidebar / Box
    \draw[fill=popgoldL, draw=popgold, thick] (3.5, 4.2) rectangle (6.2, 1.5);
    \node[popdark, align=center] at (4.85, 2.85) {Encadré\\(Infos pratiques)};
    % Layer order indicator
    \node[poppurple, font=\tiny, rotate=90] at (6.5, 5) {Ordre des calques (z-index)};
\end{scope}
\end{tikzpicture}
\legende{Comparatif visuel : à gauche, le flux linéaire du traitement de texte où l'image est prisonnière du paragraphe ; à droite, la toile de PAO où chaque cadre (titre, image, colonnes, encadré) est un objet indépendant positionné avec précision.}
\end{popfigure}

\begin{center}
\begin{tabularx}{\linewidth}{|l|X|X|}
\hline
\rowcolor{popblueL}
\textbf{Critère} & \textbf{Traitement de texte (ex. Word, Writer)} & \textbf{PAO (ex. Publisher, Scribus, InDesign)} \\
\hline
\textbf{Modèle de mise en page} & Flux continu (linéaire) & Blocs / Cadres / Objets positionnés (libre) \\
\hline
\textbf{Positionnement des images} & Ancrage au paragraphe ou à la page (relatif) & Coordonnées absolues $(x,y)$, alignement sur grille, magnétisme \\
\hline
\textbf{Gestion des colonnes} & Colonnes de section (texte qui coule) & Cadres de texte indépendants, chaînables à volonté \\
\hline
\textbf{Repères d'impression} & Marges simples & Fond perdu, traits de coupe, repères de pliage, marques d'enregistrement \\
\hline
\textbf{Gestion des couleurs} & RVB (écran) principalement & CMJN (impression), Tons directs (Pantone), Prévisualisation séparation \\
\hline
\textbf{Typographie avancée} & Basique (ligatures rares, pas de crénage fin) & Crénage (kerning), approche (tracking), césure experte, ligatures contextuelles, chasse fixe \\
\hline
\textbf{Documents types} & Lettres, rapports, thèses, CV, courriers & Affiches, dépliants, magazines, journaux, cartes de visite, packaging, signalétique \\
\hline
\textbf{Exemples logiciels} & Microsoft Word, LibreOffice Writer, Google Docs & Microsoft Publisher, Scribus (libre), Canva (en ligne), Adobe InDesign, QuarkXPress, Affinity Publisher \\
\hline
\end{tabularx}
\end{center}

\coursec{Lancement et interface du logiciel de PAO}

\courssub{Les principaux logiciels de PAO accessibles en classe de 3ème}

Le choix dépend de ton budget, de ton système d'exploitation et de la complexité du projet.

\begin{itemize}
    \item \textbf{Microsoft Publisher} : inclus dans la suite Microsoft 365 (souvent installé sur les PC des cybercafés et lycées camerounais). Interface proche de Word (ruban), gabarits prédéfinis (affiches, dépliants, cartes). Idéal pour débuter en classe de 3ème.
    \item \textbf{Scribus} : \textbf{libre, gratuit, multiplateforme} (Windows, Linux, macOS). Puissant, gère le CMJN, les fonds perdus, les PDF/X-1a (norme imprimerie). Courbe d'apprentissage un peu plus raide, mais zéro coût de licence. Parfait pour un établissement scolaire sous Ubuntu ou pour un télécentre communautaire.
    \item \textbf{Canva} : \textbf{en ligne (SaaS)}, freemium. Très populaire au Cameroun pour les réseaux sociaux (affiches d'événements Facebook, stories Instagram). Interface glisser-déposer, bibliothèques de photos/éléments. \textbf{Attention} : export PDF imprimé de qualité professionnelle réservé à la version Pro ; gestion des fonds perdus limitée en version gratuite ; nécessite une connexion internet stable (CAMTEL, MTN, Orange).
    \item \textbf{Adobe InDesign} : \textbf{standard professionnel mondial}. Abonnement coûteux (Creative Cloud). Gestion experte des livres longs (livres scolaires, \emph{Cameroon Tribune}), scripts, XML, accessibilité PDF/UA. Cible : graphistes, maisons d'édition (Éditions CLE, Presses Universitaires de Yaoundé).
    \item \textbf{Autres} : Affinity Publisher (achat unique, pas d'abonnement, excellent rapport qualité/prix), QuarkXPress (historique), Lucidpress (en ligne).
\end{itemize}

\courssub{Découverte de l'espace de travail (ex. Microsoft Publisher / Scribus)}

À l'ouverture, l'interface présente des zones communes qu'il faut identifier pour travailler efficacement.

\begin{popfigure}
\begin{tikzpicture}[
    scale=0.75,
    every node/.style={font=\small},
    panel/.style={draw=popdark, thick, rounded corners=3pt, fill=popcream},
    title/.style={font=\bfseries, popdark, anchor=west},
    label/.style={font=\tiny, popmuted, anchor=west}
]
% Main Window
\draw[panel] (0,0) rectangle (16,10);
% Top Bar - Ribbon/Menu
\draw[fill=popblueL, draw=popblue, thick] (0.2, 9.2) rectangle (15.8, 9.8);
\node[title] at (0.5, 9.5) {Barre de titre / Menu Fichier (Nouveau, Ouvrir, Enregistrer, Exporter, Imprimer)};
\draw[fill=popblueL, draw=popblue, thick] (0.2, 8.4) rectangle (15.8, 9.1);
\node[title] at (0.5, 8.75) {Ruban / Barre d'outils (Accueil, Insertion, Mise en page, Affichage, Outils)};
% Left Panel - Tools
\draw[fill=popgreenL, draw=popgreen, thick] (0.2, 0.5) rectangle (2.0, 8.2);
\node[title, rotate=90] at (0.35, 4.35) {Panneau Outils};
\node[label, rotate=90] at (0.8, 4.35) {Pointeur, Cadre texte, Cadre image, Formes, Ligne, Tableau, Zoom, Pipette...};
% Right Panel - Properties/Pages/Layers
\draw[fill=poppinkL, draw=poppink, thick] (13.5, 0.5) rectangle (15.8, 8.2);
\node[title, rotate=-90] at (15.65, 4.35) {Panneaux latéraux};
\node[label, rotate=-90, align=center, text width=6cm] at (14.8, 4.35) {Propriétés (Remplissage, Contour, Texte) | Pages (Vignettes, Gabarits) | Calques (Ordre, Verrouillage) | Couleurs};
% Center - Page Area + Scratch Area
\draw[dashed, popmuted] (2.2, 0.5) rectangle (13.3, 8.2);
\node[popmuted, font=\tiny] at (7.75, 8.0) {Zone de travail (Pasteboard / Hors page) : zone de stockage temporaire des objets};
% Page representation
\draw[fill=white, draw=popdark, very thick] (4.5, 1.5) rectangle (11.0, 7.0);
% Margins
\draw[dashed, popgreen] (5.0, 1.5) rectangle (10.5, 7.0);
\node[popgreen, font=\tiny] at (7.75, 7.15) {Marges intérieures (zone de sécurité)};
% Bleed
\draw[dashed, poporange] (4.2, 1.2) rectangle (11.3, 7.3);
\node[poporange, font=\tiny] at (7.75, 1.0) {Fond perdu (Bleed) 3-5 mm};
% Rulers
\draw[fill=popmuted!18, draw=popmuted] (2.2, 7.0) rectangle (13.3, 7.3); % Horizontal ruler
\draw[fill=popmuted!18, draw=popmuted] (2.2, 1.5) rectangle (2.5, 7.0); % Vertical ruler
\node[popmuted, font=\tiny, rotate=90] at (2.1, 4.25) {Règle verticale (cm/mm)};
\node[popmuted, font=\tiny] at (7.75, 7.45) {Règle horizontale (cm/mm)};
% Guides example
\draw[poppurple, dashed, thick] (7.75, 1.5) -- (7.75, 7.0);
\draw[poppurple, dashed, thick] (4.5, 4.25) -- (11.0, 4.25);
\node[poppurple, font=\tiny] at (7.75, 4.0) {Guides de colonne / Grille};
% Status bar
\draw[fill=popmuted!18, draw=popmuted] (0.2, 0.0) rectangle (15.8, 0.4);
\node[label] at (0.5, 0.2) {Zoom: 100\% | Page 1/1 | CMJN | Langue: Français};
\end{tikzpicture}
\legende{Anatomie de l'interface type d'un logiciel de PAO (Publisher / Scribus) : zones de commande, zone de page avec repères (marges, fond perdu, guides), panneaux d'inspection et zone de travail (pasteboard).}
\end{popfigure}

\begin{aretenir}[Éléments clés de l'interface à maîtriser]
\begin{enumerate}
    \item \textbf{La zone de page (Page)} : Représentation visuelle de la feuille (A4, A3, etc.). C'est la zone d'impression finale.
    \item \textbf{La zone de travail / Pasteboard (Hors page)} : Espace gris autour de la page. Sert à stocker temporairement des images, des blocs de texte ou des variantes avant de les placer sur la page. \textbf{N'est pas imprimé}.
    \item \textbf{Les règles et repères (Guides)} : Règles horizontales/verticales (unités : mm, cm, pts). Permettent de tirer des \textbf{guides} (lignes bleues/vertes non imprimables) pour aligner les objets avec précision (magnétisme).
    \item \textbf{Les marges et le fond perdu (Bleed)} : Les marges définissent la zone de sécurité (texte/logo à \textgreater 5 mm du bord). Le fond perdu (3 à 5 mm) est la zone qui sera rognée à la coupe ; les images de fond doivent s'y étendre pour éviter un filet blanc.
    \item \textbf{Le panneau Pages / Gabarits (Master Pages)} : Gestion des pages (ajout, suppression, duplication) et surtout des \textbf{gabarits} (en-têtes, pieds de page, numérotation, grille de colonnes communs à plusieurs pages).
    \item \textbf{Le panneau Calques (Layers)} : Gestion de l'ordre de superposition (z-index), verrouillage, visibilité. Essentiel pour sélectionner un objet caché derrière un autre.
\end{enumerate}
\end{aretenir}

\coursec{Création d'une publication : page et structure}

\courssub{Démarrage : Nouveau document et gabarits}

\begin{methode}[Créer une publication structurée (ex. Affiche A3 ou Dépliant 3 volets)]
\begin{enumerate}
    \item \textbf{Fichier $\rightarrow$ Nouveau} : Choisir \emph{Publication vierge} ou un \emph{Modèle} (Flyer, Brochure, Affiche). Pour apprendre, préfère \textbf{Publication vierge}.
    \item \textbf{Configuration de la page} (\texttt{Mise en page $\rightarrow$ Taille} ou \texttt{Fichier $\rightarrow$ Paramètres de page}) :
        \begin{itemize}
            \item \textbf{Format} : A4 (210×297 mm), A3 (297×420 mm), A5 (148×210 mm), ou \emph{Personnalisé} (ex. Carte de visite 85×55 mm).
            \item \textbf{Orientation} : Portrait ou Paysage.
            \item \textbf{Marges} : Intérieures (sécurité texte) \textgreater 5 mm. Extérieures selon besoin.
            \item \textbf{Fond perdu (Bleed)} : \textbf{Activer} et régler à \textbf{3 mm} (minimum) ou 5 mm (recommandé imprimeur).
        \end{itemize}
    \item \textbf{Gabarits (Master Pages)} : Ouvrir le panneau \texttt{Pages} $\rightarrow$ double-cliquer sur \texttt{Gabarit A} (ou \texttt{Maître A}).
        \begin{itemize}
            \item Placer les éléments \textbf{répétitifs} : numéros de page (Insertion $\rightarrow$ Numéro de page), en-tête/pied de page (nom journal, date), grille de colonnes, repères de pliage.
            \item Ces éléments apparaissent automatiquement sur toutes les pages liées à ce gabarit. \textbf{Ne pas y mettre le contenu unique de chaque page}.
        \end{itemize}
    \item \textbf{Grille et guides} : \texttt{Mise en page $\rightarrow$ Guides} $\rightarrow$ \emph{Guides de colonnes} (ex. 3 colonnes, gouttière 4 mm) et \emph{Guides de lignes de base} (alignement du texte sur une grille invisible). Activer \texttt{Affichage $\rightarrow$ Aligner sur les guides}.
\end{enumerate}
\end{methode}

\courssub{Structure d'un dépliant (pliage)}

Un dépliant 2 volets (1 pli) ou 3 volets (2 plis) impose une \textbf{structure de pages précise} pour que le pli tombe juste et que l'ordre de lecture soit correct une fois plié.

\begin{popfigure}
\begin{tikzpicture}[
    scale=0.7,
    every node/.style={font=\small},
    page/.style={draw=popdark, thick, fill=popcream, rounded corners=2pt},
    fold/.style={draw=poporange, dashed, very thick},
    arrow/.style={->, >=Stealth, popblue, thick}
]
% RECTO (Avant impression)
\begin{scope}[yshift=0]
    \node[popdark, font=\bfseries\large] at (7, 6.5) {RECTO (Face avant - Vue à plat)};
    % Page 1 (Volet 1 - Couverture)
    \draw[page] (0.5, 0.5) rectangle (6.5, 5.5);
    \node[align=center, popdark] at (3.5, 3) {PAGE 1\\Volet 1 (Couverture)\\Titre, Image forte, Date};
    \node[poporange, font=\tiny] at (3.5, 0.2) {Bord droit = Pli 1};
    % Page 2 (Volet 2 - Intérieur gauche)
    \draw[page] (6.5, 0.5) rectangle (12.5, 5.5);
    \node[align=center, popdark] at (9.5, 3) {PAGE 2\\Volet 2 (Intérieur)\\Sommaire, Edito, Plan};
    \node[poporange, font=\tiny] at (9.5, 0.2) {Bord gauche = Pli 1 | Bord droit = Pli 2};
    % Page 3 (Volet 3 - Intérieur droit)
    \draw[page] (12.5, 0.5) rectangle (18.5, 5.5);
    \node[align=center, popdark] at (15.5, 3) {PAGE 3\\Volet 3 (Intérieur)\\Corps principal, Articles};
    \node[poporange, font=\tiny] at (15.5, 0.2) {Bord gauche = Pli 2};
    % Fold lines
    \draw[fold] (6.5, 0.5) -- (6.5, 5.5) node[midway, right, poporange, font=\small] {Pli 1 (Mountain fold)};
    \draw[fold] (12.5, 0.5) -- (12.5, 5.5) node[midway, right, poporange, font=\small] {Pli 2 (Valley fold)};
    % Repères de pliage (ticks)
    \foreach \x in {6.5, 12.5} {
        \draw[poporange, very thick] (\x, 0.5) -- (\x, 0.2);
        \draw[poporange, very thick] (\x, 5.5) -- (\x, 5.8);
    }
\end{scope}
% VERSO (Dos)
\begin{scope}[yshift=-7]
    \node[popdark, font=\bfseries\large] at (7, 6.5) {VERSO (Face arrière - Vue à plat)};
    % Page 4 (Volet 3 verso - Suite articles)
    \draw[page] (0.5, 0.5) rectangle (6.5, 5.5);
    \node[align=center, popdark] at (3.5, 3) {PAGE 4\\Volet 3 (Verso)\\Suite articles, Photos};
    % Page 5 (Volet 2 verso - Infos pratiques)
    \draw[page] (6.5, 0.5) rectangle (12.5, 5.5);
    \node[align=center, popdark] at (9.5, 3) {PAGE 5\\Volet 2 (Verso)\\Infos pratiques, Contacts,\\Horaires, Sponsors};
    % Page 6 (Volet 1 verso - Dos couverture / Adresse)
    \draw[page] (12.5, 0.5) rectangle (18.5, 5.5);
    \node[align=center, popdark] at (15.5, 3) {PAGE 6\\Volet 1 (Verso/Dos)\\Adresse retour, Logo,\\Code barre, Mentions légales};
    % Fold lines
    \draw[fold] (6.5, 0.5) -- (6.5, 5.5);
    \draw[fold] (12.5, 0.5) -- (12.5, 5.5);
    \foreach \x in {6.5, 12.5} {
        \draw[poporange, very thick] (\x, 0.5) -- (\x, 0.2);
        \draw[poporange, very thick] (\x, 5.5) -- (\x, 5.8);
    }
\end{scope}
% Reading order arrow
\draw[arrow] (19, 3) -- (19, -3.5) node[midway, right, popblue, font=\small] {Ordre de lecture une fois plié : 1 $\rightarrow$ 2 $\rightarrow$ 3 $\rightarrow$ 4 $\rightarrow$ 5 $\rightarrow$ 6};
\end{tikzpicture}
\legende{Structure d'un dépliant 3 volets (6 pages) : correspondance entre la vue à plat (imposition) et l'ordre de lecture plié. Les repères de pliage (traits orange) sont indispensables sur le gabarit.}
\end{popfigure}

\begin{attention}[Piège : Ordre des pages vs Imposition]
Ne confonds pas l'\textbf{ordre de création dans le logiciel} (souvent séquentiel 1, 2, 3, 4, 5, 6) et l'\textbf{imposition imprimeur}. En PAO (Publisher, Scribus, InDesign), tu travailles en \textbf{ordres de lecture logiques} (pages face à face : 1-2, 3-4, 5-6 ou pages simples 1, 2, 3...). C'est le logiciel (ou l'imprimeur via un logiciel d'imposition) qui réorganisera les pages pour l'impression recto-verso sur grande feuille (ex. format SRA3) avant pliage et massicotage. \textbf{Toi, tu conçois dans l'ordre de lecture.}
\end{attention}

\coursec{Insertion et mise en forme des objets}

En PAO, \textbf{tout est objet}. On ne tape pas directement sur la page, on crée un conteneur (cadre) puis on y verse le contenu.

\courssub{Les cadres de texte : création, chaînage et colonnes}

\begin{methode}[Maîtriser le texte en PAO]
\begin{enumerate}
    \item \textbf{Créer un cadre} : Outil \texttt{Cadre de texte} (icône \texttt{A} ou \texttt{T}) $\rightarrow$ Cliquer-glisser sur la page pour définir la zone.
    \item \textbf{Importer du texte} : \texttt{Fichier $\rightarrow$ Importer / Placer} (ou \texttt{Ctrl+D}) $\rightarrow$ Sélectionner un fichier \texttt{.txt}, \texttt{.odt}, \texttt{.docx}. Le texte entre dans le cadre. \textbf{Avantage} : les styles (Titre 1, Corps de texte) sont souvent conservés.
    \item \textbf{Chaînage (Linking) - \cle{Notion clé}} : Si le texte déborde (icône \texttt{+} ou flèche rouge en bas du cadre), cliquer sur l'icône de débordement $\rightarrow$ le curseur change (pichet/flèche) $\rightarrow$ cliquer sur un \textbf{autre cadre vide} (même page ou page suivante). Le texte « coule » du premier vers le second. \textbf{Rompre le chaînage} : clic droit sur le cadre $\rightarrow$ \emph{Rompre le lien}.
    \item \textbf{Colonnes dans un cadre} : Sélectionner le cadre $\rightarrow$ \texttt{Format $\rightarrow$ Colonnes} (ou Propriétés) $\rightarrow$ Nombre de colonnes, Largeur, Gouttière. Le texte se répartit automatiquement \textbf{à l'intérieur du même cadre}.
    \item \textbf{Contour de texte (Text Wrap)} : Sélectionner une \textbf{image} ou une \textbf{forme} $\rightarrow$ \texttt{Format $\rightarrow$ Contour de texte} (ou \texttt{Fenêtre $\rightarrow$ Contour de texte}) $\rightarrow$ Choisir : \emph{Carré}, \emph{Serré} (suit la forme), \emph{Au-dessus/Dessous}, \emph{Traversant}. Définir les \textbf{marges de contour} (espace entre image et texte).
\end{enumerate}
\end{methode}

\courssub{Les cadres d'image : importation, ajustement et résolution}

\begin{propriete}[Règles d'or pour les images en PAO (Prépresse)]
\begin{itemize}
    \item \textbf{Résolution} : \textbf{300 ppp (pixels par pouce / dpi)} à la \textbf{taille finale d'utilisation} sur la page. Une image 300 ppp réduite à 50\% devient 600 ppp (OK). Agrandie à 200\%, elle tombe à 150 ppp (floue à l'impression). \textbf{Vérifier} dans le panneau \texttt{Propriétés de l'image} / \texttt{Infos}.
    \item \textbf{Format} : \textbf{TIFF} (sans perte, CMJN), \textbf{PSD} (Photoshop, calques conservés), \textbf{PDF} (vectoriel + raster), \textbf{JPEG} (qualité max, mais compression avec pertes), \textbf{PNG} (RVB, transparence, pour écran ou impression numérique simple). \textbf{Éviter} : GIF, BMP, images copiées-collées depuis le web (72 ppp, RVB).
    \item \textbf{Mode couleur} : \textbf{CMJN} pour l'impression offset. \textbf{RVB} uniquement pour écran ou impression jet d'encre bureautique. Vérifier : \texttt{Fenêtre $\rightarrow$ Sortie $\rightarrow$ Aperçu séparation} (Scribus/InDesign).
    \item \textbf{Ajustement dans le cadre} : \emph{Ajuster le contenu au cadre} (déforme), \emph{Ajuster le cadre au contenu} (redimensionne le cadre), \emph{Remplir le cadre proportionnellement} (recadre, conserve ratio), \emph{Adapter le contenu proportionnellement} (tout voir, espace vide possible). \textbf{Privilégier « Remplir le cadre proportionnellement » puis déplacer l'image dans le cadre (double-clic pour déplacer le contenu indépendamment du cadre)}.
\end{itemize}
\end{propriete}

\courssub{Formes, lignes, tableaux et calques}

\begin{itemize}
    \item \textbf{Formes vectorielles} (Rectangle, Ellipse, Polygone, Étoile, Flèche) : Traits (épaisseur, pointillés, couleur), Remplissage (couleur unie, dégradé, motif, image). Utiles pour : fonds de couleur, encadrés, séparateurs, pictogrammes.
    \item \textbf{Tableaux} : \texttt{Insertion $\rightarrow$ Tableau}. Moins puissants qu'Excel/Calc (pas de formules complexes), mais utiles pour des grilles de tarifs, horaires. Conversion possible : Texte $\rightarrow$ Tableau (séparateur tabulation/virgule).
    \item \textbf{Calques (Layers)} : Chaque objet vit sur un calque. Par défaut : \emph{Calque 1}. Créer des calques nommés : \emph{Fond}, \emph{Images}, \emph{Texte}, \emph{Repères}. Permet de \textbf{verrouiller} le fond pour ne pas le déplacer par erreur, de masquer le texte pour vérifier les images, de gérer la transparence globale.
    \item \textbf{Alignement et Distribution} : Sélectionner plusieurs objets (Maj + clic) $\rightarrow$ \texttt{Organiser $\rightarrow$ Aligner} (Gauche, Centres, Haut, Bas) / \texttt{Distribuer} (Horizontalement, Verticalement). \textbf{Indispensable} pour une mise en page professionnelle.
\end{itemize}

\coursec{Finalisation, exportation et impression}

\courssub{La vérification avant envoi (Preflight / Contrôle prépresse)}

Avant d'exporter, le contrôle technique évite les mauvaises surprises chez l'imprimeur (SOPECAM, imprimerie de quartier, cybercafé avec traceur).

\begin{methode}[Check-list de vérification (Preflight)]
\begin{enumerate}
    \item \textbf{Texte} : Aucune \textbf{erreur orthographique} (correcteur + relecture humaine). Aucun \textbf{débordement de texte} (icône \texttt{+} rouge sur un cadre). Pas de \textbf{polices manquantes} (substitution invisible) : \texttt{Fichier $\rightarrow$ Polices} $\rightarrow$ tout doit être \emph{Installé} ou \emph{Incorporé}.
    \item \textbf{Images} : Toutes en \textbf{CMJN} (ou RVB assumé). Toutes à \textbf{$\ge$ 300 ppp} effective. Aucune image \textbf{liée manquante} (chemin cassé) : \texttt{Fichier $\rightarrow$ Gestionnaire d'images} $\rightarrow$ tout \emph{OK}. Pas d'image en \textbf{RVB vive} (bleu électrique, vert fluo) non reproductible en CMJN sans décalage.
    \item \textbf{Couleurs} : Vérifier les \textbf{tons directs (Pantone)} : sont-ils voulus ? (Coût +1 plaque). Sinon, convertir en \textbf{Quadrichromie (CMJN)}. Vérifier le \textbf{Noir} : \emph{Noir pur} (K=100\%) pour le texte fin ; \emph{Noir riche} (C=40 M=30 J=30 K=100) pour les grands aplats.
    \item \textbf{Géométrie} : \textbf{Fonds perdus} présents (images de fond débordent de 3-5 mm). \textbf{Zone de sécurité} respectée (texte/logo \textgreater 5 mm du bord de coupe). \textbf{Repères de coupe} (Crop marks) activés à l'export. \textbf{Repères de pliage} (Fold marks) si dépliant.
    \item \textbf{Épreuves} : Imprimer une \textbf{épreuve composite} (sur imprimante bureau) pour valider la mise en page, les plis, l'ordre des pages. Plier physiquement le dépliant pour tester.
\end{enumerate}
\end{methode}

\courssub{Exportation PDF : le standard d'échange}

Le \textbf{PDF (Portable Document Format)} est le format universel pour l'impression et la diffusion.

\begin{propriete}[Préréglages d'export PDF selon la destination]
\begin{center}
\begin{tabularx}{\linewidth}{|l|X|X|}
\hline
\rowcolor{popblueL}
\textbf{Destination} & \textbf{Préréglage / Standard} & \textbf{Options critiques à cocher} \\
\hline
\textbf{Imprimeur professionnel (Offset)} & \textbf{PDF/X-1a:2001} (ou PDF/X-3, PDF/X-4) & \textbf{Fonds perdus (3-5 mm)}, \textbf{Traits de coupe}, \textbf{Repères de pliage}, \textbf{Incorporer toutes les polices}, \textbf{Convertir les couleurs en CMJN (profil ISO Coated v2 / FOGRA39)}, \textbf{Résolution images 300 ppp}. \\
\hline
\textbf{Impression bureautique / Photocopie} & \textbf{Qualité supérieure / Impression} & Traits de coupe (optionnel), Fonds perdus, Incorporer polices. RVB accepté. \\
\hline
\textbf{Diffusion Écran / WhatsApp / Email / Site web} & \textbf{Plus petit taille de fichier / Standard} & \textbf{Compression images} (JPEG qualité 80-90\%), \textbf{RVB}, \textbf{Pas de traits de coupe}, \textbf{Pas de fonds perdus}, \textbf{Hyperliens actifs}, \textbf{Signets (sommaire)}. \\
\hline
\textbf{Archivage légal / Long terme} & \textbf{PDF/A-1b} (ou PDF/A-2u) & \textbf{Tout incorporé} (polices, profils ICC, images), \textbf{Pas de chiffrement}, \textbf{Pas de JavaScript}, \textbf{Métadonnées XMP} complètes. \\
\hline
\end{tabularx}
\end{center}
\end{propriete}

\begin{attention}[Erreur fatale : « Enregistrer sous PDF » depuis Word/Publisher vs « Exporter »]
Dans Publisher / Scribus / InDesign : utilise \textbf{Fichier $\rightarrow$ Exporter $\rightarrow$ PDF} (ou \texttt{Publier au format PDF}). \textbf{Pas} « Enregistrer sous » $\rightarrow$ PDF. L'export offre la boîte de dialogue complète avec les options prépresse (PDF/X, fonds perdus, traits de coupe, profils ICC). « Enregistrer sous » produit souvent un PDF « minimal » sans repères, sans fonds perdus, avec polices non incorporées $\rightarrow$ \textbf{Rejeté par l'imprimeur}.
\end{attention}

\courssub{Impression directe et « Pack and Go » (Collecte pour sortie)}

\begin{itemize}
    \item \textbf{Impression directe} : \texttt{Fichier $\rightarrow$ Imprimer}. Utile pour épreuve. Options : \emph{Séparations} (une page par couleur C, M, J, N + tons directs) pour vérifier les surimpressions ; \emph{Composite CMJN} (épreuve couleur unique) ; \emph{Marques d'imprimeur} (traits de coupe, barres de contrôle).
    \item \textbf{Collecte pour sortie / Pack and Go (Publisher) / Collect for Output (Scribus/InDesign)} : \texttt{Fichier $\rightarrow$ Enregistrer sous $\rightarrow$ Package / Collecte pour sortie}. Crée un dossier contenant : le fichier source \texttt{.pub} / \texttt{.sla} / \texttt{.indd}, \textbf{toutes les images liées} (copiées localement), \textbf{toutes les polices utilisées} (fichiers \texttt{.ttf}, \texttt{.otf}), un \textbf{rapport d'inventaire}. \textbf{Indispensable} pour transférer le dossier à l'imprimeur ou sur un autre poste sans casser les liens.
\end{itemize}

\coursec{Démarche de production et justification}

Produire une publication, ce n'est pas « faire joli », c'est \textbf{répondre à un besoin de communication} par une démarche rigoureuse.

\begin{methode}[Démarche type de production PAO (Pour le BEPC : savoir la décrire et la justifier)]
\begin{enumerate}
    \item \textbf{Analyse du besoin (Cahier des charges)} : 
        \begin{itemize}
            \item \textbf{Cible} : Qui lit ? (Jeunes, parents, autorités, clients).
            \item \textbf{Message clé} : Quelle est l'info \#1 à retenir en 3 sec ? (Date, Lieu, Prix, Appel à l'action).
            \item \textbf{Support final} : Affiche A3 murale ? Dépliant main à main ? PDF WhatsApp ? Carte de visite ? $\rightarrow$ Détermine format, orientation, papier, finitions.
            \item \textbf{Contraintes techniques} : Budget (couleurs ? papier ?), Délai, Logiciel dispo, Compétences équipe.
        \end{itemize}
    \item \textbf{Maquette papier (Rough / Wireframe)} : \textbf{Sur papier}, crayon. Dessiner la structure : zones titre, image, texte, logo, mentions légales. Tester les plis. \textbf{Valider avec le commanditaire} avant d'ouvrir l'ordinateur. Gain de temps énorme.
    \item \textbf{Choix du logiciel} : Justifier (ex. « Scribus car gratuit, gère CMJN/fonds perdus, pas d'internet requis au cybercafé »).
    \item \textbf{Préparation des ressources} : 
        \begin{itemize}
            \item Textes : Rédigés, corrigés, stylisés dans Writer/Word $\rightarrow$ export \texttt{.odt}/\texttt{.docx}.
            \item Images : Retouchées (recadrage, niveaux, conversion CMJN) dans GIMP/Photoshop $\rightarrow$ stockées dans dossier \texttt{/images} à 300 ppp.
            \item Logos : Version \textbf{vectorielle} (\texttt{.svg}, \texttt{.eps}, \texttt{.pdf}) impérative pour netteté à toute taille.
            \item Palette couleurs : Codes CMJN/Pantone exacts (charte graphique).
        \end{itemize}
    \item \textbf{Mise en page (Production)} : 
        \begin{enumerate}
            \item Configurer document (format, marges, fonds perdus, gabarits, grille).
            \item Créer la structure (cadres vides, guides, calques verrouillés pour fond).
            \item Importer textes $\rightarrow$ chaînage $\rightarrow$ application styles.
            \item Importer images $\rightarrow$ ajustement proportionnel $\rightarrow$ contour de texte.
            \item Affiner : crénage, césure, alignement, distribution, hiérarchie visuelle.
        \end{enumerate}
    \item \textbf{Contrôle qualité (Preflight)} : Check-list complète (texte, images, couleurs, géométrie, plis).
    \item \textbf{Export / Livraison} : PDF/X-1a pour imprimeur + PDF écran pour diffusion + Package (fichiers sources) pour archivage/modif future.
    \item \textbf{Rédaction de la justification (Compétence BEPC)} : Expliquer \textbf{POURQUOI} chaque choix technique a été fait (ex. « J'ai choisi Scribus car le budget est nul et l'imprimeur exige du PDF/X-1a avec fonds perdus 3 mm, ce que Word ne gère pas nativement »).
\end{enumerate}
\end{methode}

\courssub{Critères de qualité d'une publication (Rappel synthétique)}

\begin{aretenir}[Les 5 piliers d'une publication réussie]
\begin{enumerate}
    \item \textbf{Lisibilité} : police de corps adaptée (9–11 pt pour le corps de texte), interlignage suffisant (1,2× à 1,5× la taille de police), contraste fort texte/fond (noir sur blanc, blanc sur bleu marine), longueur de ligne idéale (45–75 caractères par ligne).
    \item \textbf{Équilibre et alignement} : utilisation d'une \textbf{grille de mise en page} (marges, gouttières, colonnes, lignes de base). Alignement gauche, centré, justifié ou droit \emph{de façon cohérente}. Règle d'or : \emph{tout élément doit avoir une relation visuelle avec un autre élément}.
    \item \textbf{Cohérence des couleurs} : palette limitée (2–3 couleurs principales + neutres). Respect du mode colorimétrique : \textbf{CMJN} pour l'impression (Cyan, Magenta, Jaune, Noir), \textbf{RVB} pour l'écran. Vérifier les \textbf{tons directs} (Pantone) si logo institutionnel (ex. : vert/rouge/jaune exacts du drapeau pour une affiche du 20 mai).
    \item \textbf{Hiérarchie de l'information} : le lecteur doit comprendre \textbf{qui, quoi, où, quand} en 3 secondes. Titraille (titre principal) $\rightarrow$ Chapô (résumé) $\rightarrow$ Corps de texte $\rightarrow$ Légendes $\rightarrow$ Mentions légales (imprimeur, dépôt légal, ISSN pour les journaux).
    \item \textbf{Qualité technique (Prépresse)} : images à \textbf{300 ppp} (pixels par pouce) à la taille finale d'impression (pas 72 ppp du web !), polices \textbf{vectorielles} ou embeddées dans le PDF, fonds perdus de \textbf{3 à 5 mm} de chaque côté, traits de coupe, pas de texte trop près du bord (zone de sécurité 5 mm).
\end{enumerate}
\end{aretenir}

\begin{exemplebox}[Identifier les documents relevant de la PAO]
Parmi la liste suivante, quels documents nécessitent impérativement un logiciel de PAO (plutôt qu'un traitement de texte) ?
\begin{enumerate}
    \item Une dissertation de 10 pages sur « Les causes de la déforestation au Cameroun ».
    \item Une affiche A3 pour la fête de la musique au quartier Mvog-Ada.
    \item Une lettre de motivation pour un stage à la CAMTEL.
    \item Un dépliant 3 volets (format A4 plié en 3) pour la cantine scolaire.
    \item Un CV d'une page.
    \item La une du journal scolaire \emph{Le Lycéen} (4 pages, colonnes, photos, publicités).
    \item Un faire-part de mariage format carte 12×17 cm avec dorure.
\end{enumerate}
\tcblower
\textbf{Analyse :}
\begin{itemize}
    \item \textbf{PAO obligatoire} : \textbf{2 (Affiche A3)}, \textbf{4 (Dépliant 3 volets : plis précis, recto/verso alignés)}, \textbf{6 (Une de journal : colonnes, chaînage, mise en page complexe)}, \textbf{7 (Faire-part : format non standard, finitions spéciales)}.
    \item \textbf{Traitement de texte suffisant} : \textbf{1 (Dissertation : flux continu, notes de bas de page, sommaire auto)}, \textbf{3 (Lettre : structure simple, en-tête/pied de page)}, \textbf{5 (CV : structure linéaire, bien que Publisher/Canva fassent de beaux CV)}.
\end{itemize}
\textbf{Critère décisionnel} : \emph{Ai-je besoin de positionner des blocs indépendants, de gérer des plis, des fonds perdus, du chaînage de texte ou du CMJN ?} $\rightarrow$ PAO.
\end{exemplebox}

\begin{attention}[Piège classique : « Word fait de la PAO »]
\textbf{Non.} Microsoft Word est un \textbf{traitement de texte}. Certes, depuis Word 2010, tu peux insérer des « Zones de texte », faire du « Positionnement : Devant le texte », gérer des colonnes, des sauts de section, et même exporter en PDF/XPS. \textbf{Mais} :
\begin{itemize}
    \item Pas de gestion native du \textbf{CMJN} (conversion RVB $\rightarrow$ CMJN approximative à l'export).
    \item Pas de \textbf{fonds perdus} automatiques (il faut ruser avec une page plus grande).
    \item Pas de \textbf{traits de coupe} ni repères d'imprimeur.
    \item Le \textbf{chaînage de cadres} (texte qui coule d'un cadre à un autre non adjacent) est laborieux et buggé.
    \item Le \textbf{moteur de composition} (césure, justification, crénage) reste inférieur à celui d'InDesign ou Scribus.
\end{itemize}
\textbf{Utilise Word pour écrire le texte, puis importe-le (Fichier $\rightarrow$ Placer / Importer) dans ton logiciel de PAO pour la mise en page finale.}
\end{attention}

\begin{savaistu}[Le Cameroun et la PAO professionnelle]
Le quotidien national \textbf{Cameroon Tribune} (créé en 1974) et son pendant anglophone sont entièrement mis en page avec des logiciels de PAO professionnels (historiquement QuarkXPress, aujourd'hui majoritairement \textbf{Adobe InDesign} sur serveur centralisé). Chaque soir, les journalistes envoient leurs articles (texte brut + métadonnées) via le système éditorial ; les maquettistes « versent » le texte dans les gabarits de la une, des pages sports, culture, etc., gèrent les photos haute résolution envoyées par les correspondants régionaux (Bertoua, Maroua, Bamenda), vérifient le \textbf{fer à repasser} (alignement des colonnes entre pages paires/impaires), génèrent le \textbf{PDF haute définition (PDF/X-1a)} envoyé à la rotative de la \textbf{SOPECAM} (Société de Presse et d'Édition du Cameroun) à Yaoundé. C'est une chaîne de production PAO industrielle : \textbf{plus de 50 000 exemplaires/jour}, imprimés en quadrichromie sur papier journal, distribués avant l'aube sur tout le triangle national.
\end{savaistu}

\begin{exoresolu}[Choix du logiciel selon le contexte]
\textbf{Énoncé :} L'association « Jeunesse Active de Nkolbisson » doit produire trois supports pour son assemblée générale annuelle : (1) le rapport d'activités de 20 pages (texte, quelques graphiques), (2) l'affiche A2 d'invitation (visuel fort, date, lieu, sponsors), (3) le programme pliable 2 volets (ordre du jour, bios des intervenants, plan d'accès). Le trésorier a un budget nul. Le cybercafé du coin a Windows 10 avec Office 2019 et une connexion internet instable. Quel(s) logiciel(s) recommandes-tu pour chaque support ? Justifie.

\tcblower
\textbf{Solution détaillée :}

\begin{center}
\begin{tabularx}{\linewidth}{|l|X|X|}
\hline
\rowcolor{popblueL}
\textbf{Support} & \textbf{Logiciel recommandé} & \textbf{Justification technique} \\
\hline
(1) Rapport d'activités 20 p. & \textbf{LibreOffice Writer} (ou Word) & Flux continu, sommaire auto, notes de bas de page, numérotation chapitres, export PDF/A pour archivage. Pas de mise en page complexe. \\
\hline
(2) Affiche A2 & \textbf{Scribus} (installé hors ligne) & Format A2 (420×594 mm) > limites Word (55,87 cm max). Gestion CMJN pour impression grand format. Blocs libres, fonds perdus 5 mm, traits de coupe. Gratuit, pas d'internet requis. \\
\hline
(3) Programme 2 volets & \textbf{Scribus} ou \textbf{Publisher} & Pli central précis (repères de pliage). Recto/verso aligné (registration). Chaînage texte entre volets. Publisher si déjà maîtrisé par l'équipe ; Scribus sinon (plus robuste pour l'imprimeur). \\
\hline
\end{tabularx}
\end{center}

\textbf{Workflow optimal :}
\begin{enumerate}
    \item Rédiger tout le texte dans \textbf{Writer} (correction orthographique, styles titres/paragraphes).
    \item Exporter en \texttt{.odt} ou \texttt{.docx}.
    \item Ouvrir \textbf{Scribus} $\rightarrow$ \texttt{Fichier > Importer > Texte} pour faire couler le rapport dans les cadres du gabarit « Rapport ».
    \item Créer l'affiche et le programme dans Scribus (nouvelles pages/documents), importer les textes nettoyés.
    \item Exporter les 3 PDF : \texttt{PDF/X-1a:2001} (imprimeur) + \texttt{PDF 1.4} (diffusion WhatsApp/Email).
\end{enumerate}
\textbf{Astuce :} Scribus portable (sur clé USB) permet de travailler au cybercafé sans droits admin.
\end{exoresolu}

\begin{exercice}[Mise en situation : L'affiche du tournoi inter-quartiers]
Tu dois réaliser l'affiche A3 (portrait) du tournoi de football de ton quartier. L'imprimeur demande un PDF/X-1a avec fonds perdus 3 mm et traits de coupe. Tu utilises Scribus.
\begin{enumerate}
    \item Liste les \textbf{étapes de configuration du document} (format, marges, fonds perdus, gabarit) avant de poser le moindre objet.
    \item Tu as une photo du stade (fichier \texttt{stade.jpg}, 4000×3000 px, RVB, 5 Mo). Tu veux qu'elle occupe tout le fond de l'affiche. \textbf{Quelles vérifications et transformations} effectues-tu sur ce fichier \textbf{avant} de l'importer dans Scribus ? (Logiciel suggéré : GIMP).
    \item Le titre « TOURNOI INTER-QUARTIERS » doit être en haut, centré, police \texttt{Anton} (téléchargée sur Google Fonts), 72 pt, blanc avec contour noir 2 pt. \textbf{Comment t'assures-tu que la police sera visible chez l'imprimeur ?}
    \item Le programme des matchs est un tableau horaire. Tu décides de le faire dans un \textbf{cadre de texte unique à 3 colonnes} (chaînage interne). \textbf{Comment procèdes-tu pour créer ce cadre et y faire couler le texte importé d'un fichier \texttt{.odt} ?}
    \item Pour l'export final, \textbf{coches-tu l'option « Inclure les repères de pliage » ?} Justifie.
\end{enumerate}
\end{exercice}

\begin{corrige}[Exercice : L'affiche du tournoi inter-quartiers]
\begin{enumerate}
    \item \textbf{Configuration document (Fichier > Nouveau > Propriétés du document) :}
        \begin{itemize}
            \item Format : \textbf{A3} (297 × 420 mm), Orientation : \textbf{Portrait}.
            \item Marges : \textbf{10 mm} minimum (zone de sécurité pour texte/logo).
            \item \textbf{Fond perdu (Bleed) : 3 mm} sur les 4 côtés (cocher « Utiliser les fonds perdus »).
            \item Gabarit : Créer un \textbf{Gabarit « Affiche »} avec guides verticaux (marges) et horizontaux (tiers pour zones titre/image/programme), calque « Fond » verrouillé en bas.
        \end{itemize}
    \item \textbf{Traitement image (GIMP) :}
        \begin{itemize}
            \item Vérifier résolution : 4000 px / (420 mm / 25,4) $\approx$ 242 ppp. \textbf{Insuffisant} pour du 300 ppp à taille réelle A3 (hauteur 420 mm). 
            \item \textbf{Action} : Ne pas agrandir (perte qualité). Soit recadrer (crop) sur zone d'intérêt pour augmenter ppp effectif, soit accepter légère perte (200-250 ppp toléré pour grand format vu de loin), soit trouver image source plus grande.
            \item \textbf{Conversion couleur} : \textbf{Image > Mode > CMJN} (profil ISO Coated v2 300\%). Vérifier les noirs riches.
            \item \textbf{Enregistrement} : \texttt{stade\_fond\_cmyk.tif} (sans compression LZW) ou \texttt{.jpg} qualité 100.
        \end{itemize}
    \item \textbf{Gestion police \texttt{Anton} :}
        \begin{itemize}
            \item Installer la police \texttt{.ttf} sur le poste (Windows/Fonts).
            \item Dans Scribus : \texttt{Fichier > Polices > Additionnelles} : vérifier que \texttt{Anton} est listée et \textbf{« Incorporer » cochée}.
            \item À l'export PDF/X-1a : onglet \textbf{Polices} $\rightarrow$ \textbf{Incorporer toutes les polices} (ou « Incorporer seulement les glyphes utilisés » pour réduire taille). \textbf{C'est cette incorporation qui garantit l'affichage chez l'imprimeur.}
        \end{itemize}
    \item \textbf{Cadre texte 3 colonnes :}
        \begin{enumerate}
            \item Outil \texttt{Cadre de texte} : tracer un grand cadre sous l'image (hauteur ~150 mm, largeur page moins marges).
            \item Avec le cadre sélectionné : \texttt{Propriétés (F2) > Texte > Colonnes} : Nombre = \textbf{3}, Gouttière = \textbf{4 mm}.
            \item \texttt{Fichier > Importer > Texte} : choisir \texttt{programme.odt}. Cliquer dans le cadre. Le texte remplit la colonne 1, puis 2, puis 3 automatiquement.
            \item Appliquer un \textbf{Style de paragraphe} « HoraireMatch » (police 10 pt, interlignage 1.3, tabulations pour aligner scores).
        \end{enumerate}
    \item \textbf{Repères de pliage :} \textbf{NON.} C'est une \textbf{affiche plate (A3)}, pas un dépliant. Il n'y a \textbf{aucun pli}. Cocher cette option ajouterait des traits de pliage inutiles et confus pour le massicotage. On coche uniquement : \textbf{Traits de coupe}, \textbf{Fonds perdus (3 mm)}, \textbf{Marques d'enregistrement} (optionnel mais pro).
\end{enumerate}
\end{corrige}

\coursec{Lancement et interface du logiciel de PAO}

Après avoir découvert \emph{ce qu'est la PAO} et \emph{à quoi elle sert} dans la section précédente, il est temps de passer à la pratique. Tout projet de publication — qu'il s'agisse d'une affiche pour la fête de fin d'année du collège, d'un dépliant pour la boutique du quartier ou d'une carte d'invitation pour une cérémonie familiale — commence par l'ouverture du logiciel et la compréhension de son espace de travail. Connaître chaque zone de l'interface, c'est gagner du temps, éviter les erreurs de placement et produire un document soigné du premier coup.

\courssub{Lancement du logiciel et écran d'accueil}

Sur les postes de travail de la salle informatique du collège (sous Windows), le logiciel de PAO de référence est \textbf{Microsoft Publisher}. Il s'installe généralement avec la suite bureautique Microsoft Office. Pour le lancer, la procédure standard est la suivante :

\begin{methode}[Lancer Publisher et créer une publication vierge en 4 étapes]
\begin{enumerate}
    \item Cliquer sur le bouton \textbf{Démarrer} (logo Windows en bas à gauche de l'écran) ou appuyer sur la touche \texttt{Windows} du clavier.
    \item Saisir \texttt{Publisher} dans la zone de recherche ; l'icône de l'application apparaît en haut des résultats.
    \item Cliquer sur l'icône \textbf{Microsoft Publisher} pour ouvrir le logiciel.
    \item Sur l'écran d'accueil qui s'affiche, choisir \textbf{Page vierge} (format A4 par défaut) pour partir d'une page blanche, ou sélectionner un \textbf{modèle prédéfini} (Affiche, Dépliant, Carte de visite, Bulletin, etc.) pour gagner du temps sur la mise en page.
\end{enumerate}
\end{methode}

\begin{attention}[Ne pas confondre zone de fond et page imprimable]
L'écran de travail présente une grande surface grise qui entoure la feuille blanche : c'est la \textbf{zone de fond} (ou plan de travail). \textbf{Seule la page blanche (la feuille) sera imprimée ou exportée en PDF.} Tout objet (cadre de texte, image, forme) placé \emph{en dehors} de cette page, dans la zone de fond, ne figurera \emph{pas} dans le document final. C'est l'erreur n°1 des débutants : ils croient avoir « perdu » leur titre ou leur photo, alors qu'il est simplement hors page.
\end{attention}

\begin{savaistu}[Alternative libre : Scribus]
Dans les établissements équipés de Linux (par exemple Ubuntu) ou pour un usage personnel sans licence Microsoft, \textbf{Scribus} est un logiciel de PAO libre, gratuit et multiplateforme. Son interface diffère un peu (barres d'outils flottantes, pas de ruban), mais les concepts restent identiques : page, repères, cadres de texte, images. Le lancement se fait via le menu \texttt{Applications → Graphisme → Scribus}.
\end{savaistu}

\courssub{Anatomie de l'interface : les grandes zones}

Une fois la publication ouverte (vierge ou issue d'un modèle), l'interface de Publisher se décompose en zones bien distinctes. La figure ci-dessous en donne une vue d'ensemble simplifiée et annotée.

\begin{popfigure}
\begin{tikzpicture}[
    zone/.style={rectangle, draw=popdark, thick, rounded corners=3pt, inner sep=5pt, font=\small, align=center},
    lab/.style={font=\small\bfseries, text=popblue, align=center}
]
% Barre de titre
\node[zone, fill=popblueL, minimum width=13.5cm, minimum height=0.9cm] (titre) at (0,5.6)
    {\textbf{Barre de titre} : nom du fichier (ex. \texttt{Affiche\_Fete.pub}) + boutons Réduire / Agrandir / Fermer};
% Barre d'accès rapide + ruban
\node[zone, fill=poporangeL, minimum width=13.5cm, minimum height=0.9cm] (acces) at (0,4.5)
    {\textbf{Barre d'outils Accès rapide} : Enregistrer (\texttt{Ctrl+S}), Annuler (\texttt{Ctrl+Z}), Rétablir};
\node[zone, fill=popgreenL, minimum width=13.5cm, minimum height=1.1cm] (ruban) at (0,3.3)
    {\textbf{Le Ruban} : Fichier \textbullet\ Accueil \textbullet\ Insertion \textbullet\ Mise en page \textbullet\ Conception \textbullet\ Publipostage \textbullet\ Révision \textbullet\ Affichage};
% Volet des pages
\node[zone, fill=poppurpleL, minimum width=3cm, minimum height=5.6cm] (volet) at (-5.2,-0.6)
    {\textbf{Volet des}\\\textbf{pages}\\[4pt]Miniatures\\Page 1\\Page 2\\...\\[4pt]\texttt{+} Ajouter};
% Zone de travail avec la page
\node[zone, fill=popcream, minimum width=9.6cm, minimum height=5.6cm] (travail) at (1.6,-0.6) {};
\node[lab, text=popdark] at (1.6,1.9) {\textbf{Zone de travail}};
% Page blanche
\node[rectangle, draw=popdark, thick, fill=white, minimum width=4.6cm, minimum height=4.6cm] (page) at (1.6,-0.7) {};
% Marges bleues
\draw[popblue, thick, dash pattern=on 4pt off 3pt] (-0.3,1.2) -- (3.5,1.2);
\draw[popblue, thick, dash pattern=on 4pt off 3pt] (-0.3,-2.6) -- (3.5,-2.6);
\draw[popblue, thick, dash pattern=on 4pt off 3pt] (-0.3,1.2) -- (-0.3,-2.6);
\draw[popblue, thick, dash pattern=on 4pt off 3pt] (3.5,1.2) -- (3.5,-2.6);
% Colonnes roses
\draw[poppink, thick, dash pattern=on 2pt off 2pt] (1.0,1.2) -- (1.0,-2.6);
\draw[poppink, thick, dash pattern=on 2pt off 2pt] (2.3,1.2) -- (2.3,-2.6);
\node[font=\footnotesize, align=center] at (1.6,-0.7) {Page A4\\(imprimable)\\[3pt]marges \textcolor{popblue}{bleues}\\colonnes \textcolor{poppink}{roses}};
\node[font=\footnotesize\itshape, text=popmuted] at (5.0,-2.9) {Zone de fond (non imprimable)};
% Barre d'état
\node[zone, fill=popgoldL, minimum width=13.5cm, minimum height=0.9cm] (etat) at (0,-4.0)
    {\textbf{Barre d'état} : Zoom 100\% \textbullet\ Page 1 / 2 \textbullet\ Position X = 5.2cm ; Y = 3.1cm \textbullet\ Mode d'affichage};
\end{tikzpicture}
\legende{Interface de Microsoft Publisher : zones principales. La page blanche (au centre) est encadrée par les marges bleues et les repères de colonnes roses. Le volet des pages (à gauche) permet de naviguer entre les pages. La barre d'état (en bas) affiche le zoom, le numéro de page et la position du pointeur.}
\end{popfigure}

\begin{definition}[Repère]
Un \cle{repère} est une \textbf{ligne non imprimable} (donc invisible sur le document final papier ou PDF) qui s'affiche à l'écran pour aider l'utilisateur à \textbf{aligner, positionner et dimensionner} les objets (cadres de texte, images, formes) avec précision. On distingue :
\begin{itemize}
    \item les \textbf{repères de marges} (bleus) : limites de la zone utile de la page ;
    \item les \textbf{repères de colonnes et de lignes} (roses) : division de la page en grille pour structurer la mise en page ;
    \item les \textbf{repères personnalisés} : lignes horizontales ou verticales que l'on tire soi-même depuis les règles.
\end{itemize}
\end{definition}

\courssub{Le ruban : les onglets et leurs fonctions principales}

Le \textbf{ruban} (en haut, sous la barre d'accès rapide) regroupe \emph{toutes} les commandes organisées par onglets. Chaque onglet correspond à une grande étape de la production. Le tableau ci-dessous associe chaque onglet à sa fonction majeure — un réflexe à acquérir pour ne pas chercher au mauvais endroit.

\begin{exemplebox}[Associer chaque onglet du ruban à sa fonction principale]
\begin{center}
\begin{tabularx}{\linewidth}{|l|X|X|}
\hline
\rowcolor{popblueL}
\textbf{Onglet} & \textbf{Fonction principale} & \textbf{Commandes représentatives} \\
\hline
\textbf{Fichier} & Gestion du fichier (Backstage) & Nouveau, Ouvrir, Enregistrer, Imprimer, Exporter en PDF, Partager, Options \\
\hline
\textbf{Accueil} & Mise en forme de base & Police, taille, couleur, alignement, puces, styles, Rechercher/Remplacer, Sélection \\
\hline
\textbf{Insertion} & \textbf{Ajouter des objets} sur la page & Cadre de texte, Image (fichier/en ligne), Formes, WordArt, Tableau, Symbole \\
\hline
\textbf{Mise en page} & \textbf{Structure de la page} & Marges, Orientation, Taille, Repères (marges/colonnes/lignes), Arrière-plan, Bordures \\
\hline
\textbf{Conception} & Apparence globale & Schémas de couleurs, Polices de thème, Modèles de page, Variantes \\
\hline
\textbf{Publipostage} & Fusion de données & Sélectionner destinataires, Insérer champ de fusion, Aperçu, Terminer et fusionner \\
\hline
\textbf{Révision} & Contrôle et collaboration & Orthographe, Recherche, Commentaires \\
\hline
\textbf{Affichage} & Confort de travail & Zoom, Règles, Repères, Volet des pages, Fenêtres \\
\hline
\end{tabularx}
\end{center}
\end{exemplebox}

\begin{aretenir}[Raccourcis indispensables]
\begin{itemize}
    \item \texttt{Ctrl + S} : Enregistrer (à faire \emph{très} souvent).
    \item \texttt{Ctrl + Z} / \texttt{Ctrl + Y} : Annuler / Rétablir.
    \item \texttt{Ctrl + Molette de la souris} : Zoom avant / arrière rapide.
    \item \texttt{Ctrl + P} : Imprimer.
\end{itemize}
\end{aretenir}

\courssub{La zone de travail, le volet des pages et les repères}

La \textbf{zone de travail} (au centre) affiche la page courante. C'est là que l'on construit la publication. Autour de la page, la \textbf{zone de fond} (grise) sert de « table de travail » : on peut y déposer temporairement des objets, les préparer, les copier-coller avant de les glisser sur la page.

Le \textbf{volet des pages} (à gauche) affiche les miniatures de toutes les pages de la publication. Il permet :
\begin{itemize}
    \item de naviguer d'une page à l'autre par un simple clic ;
    \item d'ajouter une page (bouton \texttt{+} ou clic droit → \emph{Insérer une page}) ;
    \item de dupliquer, supprimer ou réorganiser les pages (glisser-déposer) ;
    \item d'accéder aux \textbf{pages maîtresses} (modèles de fond communs à plusieurs pages).
\end{itemize}

Les \textbf{règles} (horizontale en haut, verticale à gauche) sont graduées en centimètres. Elles indiquent la position exacte du pointeur. Pour créer un repère personnalisé : \textbf{cliquer sur la règle et glisser} vers l'intérieur de la page. Une ligne pointillée apparaît ; on la relâche à la position voulue. Pour supprimer un repère : le sélectionner (clic dessus) et appuyer sur \texttt{Suppr}, ou le ramener sur la règle.

\begin{propriete}[Repères magnétiques (aimantation)]
Par défaut, les objets « accrochent » aux repères et aux marges quand on les déplace à proximité (tolérance de quelques pixels). Cela garantit un alignement parfait sans effort. On peut désactiver temporairement l'aimantation en maintenant la touche \texttt{Alt} enfoncée pendant le déplacement.
\end{propriete}

\courssub{La barre d'état : informations contextuelles}

Tout en bas de la fenêtre, la \textbf{barre d'état} donne en temps réel :
\begin{itemize}
    \item le \textbf{niveau de zoom} (ex. \texttt{100\%}) — cliquable pour ouvrir la boîte de dialogue Zoom ;
    \item le \textbf{numéro de la page courante} et le \textbf{nombre total de pages} (ex. \texttt{Page 1 sur 3}) ;
    \item la \textbf{position du pointeur} en coordonnées \texttt{X ; Y} par rapport au coin supérieur gauche de la page (ex. \texttt{X = 5.2cm ; Y = 3.1cm}) — précieux pour placer un objet à une position exacte ;
    \item le \textbf{mode d'affichage} (Normal, Arrière-plan, Aperçu avant impression).
\end{itemize}

\begin{experience}[Prise en main de l'interface — TP guidé (10 min)]
\textbf{Objectif :} Lancer Publisher, créer une publication vierge A4, identifier chaque zone de l'interface, créer et supprimer un repère personnalisé.

\textbf{Matériel / Logiciel :} Poste Windows avec Microsoft Publisher installé.

\textbf{Manipulation :}
\begin{enumerate}
    \item Lancer Publisher via le menu Démarrer (méthode en 4 étapes vue plus haut).
    \item Choisir \textbf{Page vierge} (format A4). Observer l'écran : repérer la barre de titre, le ruban, le volet des pages, la zone de travail, les règles, la barre d'état.
    \item Sur la règle verticale (à gauche), cliquer à \texttt{5 cm} et glisser vers la droite : un repère vertical apparaît. Relâcher. Vérifier qu'il est magnétique : créer un cadre de texte (\texttt{Insertion → Cadre de texte}), le déplacer près du repère → il « colle ».
    \item Sélectionner le repère (clic dessus), appuyer sur \texttt{Suppr}. Le repère disparaît.
    \item Dans le volet des pages, cliquer sur \texttt{+} pour ajouter une page 2. Cliquer sur la miniature « Page 1 » pour revenir.
    \item Enregistrer le fichier sur le Bureau sous le nom \texttt{Test\_Interface.pub} (\texttt{Ctrl + S}).
    \item Fermer Publisher.
\end{enumerate}
\textbf{Observations attendues :} L'élève distingue la page (blanche, imprimable) de la zone de fond (grise, non imprimable). Il sait créer et supprimer un repère. Il connaît le raccourci \texttt{Ctrl+S}.

\textbf{Interprétation :} La maîtrise de l'interface est la condition \emph{sine qua non} d'une production PAO efficace. Ce TP ancre les repères visuels et les gestes de base.
\end{experience}

\courssub{Démarche de production : de l'interface à la première page structurée}

Maintenant que l'interface est familière, la démarche professionnelle pour démarrer toute publication suit ce schéma logique :

\begin{methode}[Démarrer une publication structurée (check-list)]
\begin{enumerate}
    \item \textbf{Lancer} le logiciel et choisir \textbf{Page vierge} (ou un modèle adapté au besoin).
    \item \textbf{Enregistrer immédiatement} (\texttt{Ctrl+S}) sous un nom explicite (\texttt{Affiche\_Tournoi\_Foot\_Mai2026.pub}) dans un dossier dédié (ex. \texttt{Documents/PAO/3eme/}).
    \item \textbf{Configurer la page} : onglet \texttt{Mise en page} → \texttt{Taille} (A4, A5, personnalisé), \texttt{Orientation} (Portrait/Paysage), \texttt{Marges} (Normales, Étroites, Larges ou Personnalisées).
    \item \textbf{Poser les repères de structure} : onglet \texttt{Mise en page} → \texttt{Repères} → \texttt{Repères de grille et de lignes} → définir le nombre de colonnes/lignes et l'espacement (la \emph{gouttière}). Tirer au besoin des repères personnalisés depuis les règles.
    \item \textbf{Activer l'affichage des repères} : onglet \texttt{Affichage} → cocher \texttt{Repères} et \texttt{Règles} si ce n'est pas déjà fait.
    \item \textbf{Commencer l'insertion des objets} (cadres de texte, images, formes) en les \emph{ancrant} sur les repères.
\end{enumerate}
\end{methode}

\begin{exoresolu}[Configurer une page A4 paysage avec marges étroites et 3 colonnes]
\textbf{Énoncé :} Tu dois créer un dépliant 3 volets (une feuille A4 en orientation paysage, pliée en 3). Configure la page : orientation paysage, format A4, marges étroites (1 cm), 3 colonnes avec gouttière de 0,5 cm. Place ensuite un repère vertical au milieu de la colonne centrale pour centrer un titre.

\tcblower
\textbf{Solution détaillée :}
\begin{enumerate}
    \item Lancer Publisher → \texttt{Page vierge}.
    \item \texttt{Ctrl+S} → enregistrer sous \texttt{Depliant\_3volets.pub}.
    \item Onglet \textbf{Mise en page} :
    \begin{itemize}
        \item \texttt{Taille} → \texttt{A4} (21 × 29,7 cm).
        \item \texttt{Orientation} → \texttt{Paysage} (la page fait alors 29,7 cm de large sur 21 cm de haut).
        \item \texttt{Marges} → \texttt{Étroites} (environ 1 cm de chaque côté).
    \end{itemize}
    \item Toujours dans \textbf{Mise en page} → \texttt{Repères} → \texttt{Repères de grille et de lignes} :
    \begin{itemize}
        \item Colonnes : \texttt{3}, Espacement (gouttière) : \texttt{0,5 cm}.
        \item Lignes : \texttt{1}.
        \item Valider par \texttt{OK}.
    \end{itemize}
    \item Les repères de colonnes (roses) apparaissent, divisant la page en 3 bandes égales.
    \item Pour le repère au milieu de la colonne centrale : la largeur utile de la page vaut $29,7 - 2 \times 1 = 27,7$ cm. En retirant les deux gouttières ($2 \times 0,5 = 1$ cm), chaque colonne mesure environ $(27,7 - 1) / 3 \approx 8,9$ cm. La colonne centrale commence donc vers $1 + 8,9 + 0,5 = 10,4$ cm et son milieu se situe vers $10,4 + 8,9/2 \approx 14,9$ cm, c'est-à-dire quasiment au centre de la page (ce qui est logique : la colonne centrale est au milieu). Sur la règle horizontale, cliquer vers \texttt{14,9 cm} et glisser vers le bas, puis relâcher.
    \item Vérifier : onglet \texttt{Affichage} → \texttt{Repères} coché, \texttt{Règles} coché.
    \item \texttt{Ctrl+S} pour sauvegarder.
\end{enumerate}
\end{exoresolu}

\begin{attention}[Pièges classiques à l'initialisation]
\begin{itemize}
    \item \textbf{Oublier d'enregistrer au début} : une coupure de courant ou un plantage fait tout perdre. Réflexe : \texttt{Ctrl+S} \emph{avant} la première action, puis régulièrement.
    \item \textbf{Travailler sur la zone de fond} : tout objet hors de la page blanche est invisible à l'impression. Vérifier systématiquement que les cadres sont bien \emph{sur} la page.
    \item \textbf{Mélanger les unités} : les règles sont en centimètres par défaut. Si elles sont en pouces (\texttt{in}) ou en points (\texttt{pt}), changer via \texttt{Fichier → Options → Avancées → Unités d'affichage}.
    \item \textbf{Ne pas activer les repères} : sans repères visibles, l'alignement se fait « au jugé » → résultat peu soigné. Toujours : \texttt{Affichage → Repères}.
\end{itemize}
\end{attention}

\begin{savaistu}[Le format .pub et l'interopérabilité]
Le format natif de Publisher est \texttt{.pub}. Il est \textbf{propriétaire} : seul Microsoft Publisher peut l'ouvrir correctement. Pour partager une publication avec quelqu'un qui n'a pas Publisher (par exemple un imprimeur ou un correspondant sous Linux), on \textbf{exporte en PDF} (\texttt{Fichier → Exporter → Créer un document PDF/XPS}). Le PDF préserve la mise en page exacte, les polices et les images, et se lit sur tous les appareils. C'est le format standard d'échange en PAO.
\end{savaistu}

\courssub{Synthèse visuelle : les 3 zones critiques à surveiller}

\begin{center}
\begin{tabularx}{\linewidth}{|l|X|X|}
\hline
\rowcolor{popblueL}
\textbf{Zone} & \textbf{Question clé} & \textbf{Action corrective} \\
\hline
\textbf{Page (blanche)} & « Mon objet est-il \emph{entièrement} dans la page ? » & Glisser l'objet jusqu'à ce qu'il soit à l'intérieur des marges bleues. \\
\hline
\textbf{Repères (bleus/roses)} & « Mes objets sont-ils alignés sur la grille ? » & Activer \texttt{Affichage → Repères} ; utiliser l'aimantation ; ajouter des repères personnalisés. \\
\hline
\textbf{Barre d'état} & « Quel zoom ? Quelle page ? Quelle position ? » & Ajuster le zoom (100\% pour la mise en page, 50\% pour une vue d'ensemble) ; vérifier X/Y pour un placement précis. \\
\hline
\end{tabularx}
\end{center}

Cette maîtrise de l'interface — lancement, repérage, configuration, sauvegarde — est le socle sur lequel reposent les sections suivantes : insertion d'objets, mise en forme, finalisation et exportation. Dans la prochaine section, nous apprendrons à \textbf{structurer la page} (colonnes, arrière-plan, pages maîtresses) avant d'y déposer le contenu.

\coursec{Création d'une publication : page et structure}

Après avoir découvert l'interface de notre logiciel de PAO — barres d'outils, volet des pages, zone de dessin et rubans contextuels — nous sommes maintenant prêts à poser les fondations de notre document. C'est une étape décisive : \textbf{la structure de la page détermine toute la suite du travail}. Un mauvais choix de format, des marges trop justes ou l'absence de repères de colonnes transformeront une mise en page élégante en casse-tête ingérable. Dans cette section, nous apprenons à configurer la page, à exploiter les modèles, à organiser les pages et à sauvegarder correctement notre publication.

\courssub{Choix du format de page et orientation}

\paragraph{Intuition : le contenant avant le contenu.} Imagine que tu dois concevoir une affiche pour la fête de l'école (format A3, orientation paysage), une carte de visite (petit format personnalisé, orientation portrait) ou un dépliant touristique pour le parc de la Waza (format A4, orientation paysage, plié en trois). Le logiciel de PAO ne devine pas ton intention : c'est à toi de lui indiquer \emph{quelle feuille de papier virtuelle} tu utilises avant d'y poser le moindre texte ou la moindre image.

\begin{definition}[Format de page et orientation]
Le \cle{format de page} définit les dimensions physiques (largeur $\times$ hauteur) de la zone d'impression. Les formats normalisés les plus courants sont :
\begin{itemize}
    \item \textbf{A4} : 21,0 cm $\times$ 29,7 cm (standard pour lettres, rapports, dépliants).
    \item \textbf{A5} : 14,8 cm $\times$ 21,0 cm (livrets, petits flyers) — moitié d'un A4.
    \item \textbf{Personnalisé} : dimensions libres (ex. : carte de visite 8,5 cm $\times$ 5,5 cm ; affiche 40 cm $\times$ 60 cm).
\end{itemize}
L'\cle{orientation} détermine le sens de lecture par rapport au bord long de la feuille :
\begin{itemize}
    \item \textbf{Portrait} : hauteur $>$ largeur (lecture verticale, standard bureautique).
    \item \textbf{Paysage} : largeur $>$ hauteur (lecture horizontale, idéale pour dépliants, tableaux larges, présentations).
\end{itemize}
\end{definition}

\begin{savaistu}[La norme ISO 216 : l'intelligence du format A]
Le format A4 (21 $\times$ 29,7 cm) n'est pas un hasard : il respecte le rapport $\sqrt{2} \approx 1,414$. Si tu plies une feuille A4 en deux dans le sens de la largeur, tu obtiens \textbf{deux feuilles A5} aux proportions \emph{identiques}. Cette propriété permet de passer d'un format à l'autre sans déformer les mises en page. C'est la norme internationale ISO 216, utilisée dans toutes les imprimeries du Cameroun et du monde — sauf aux États-Unis et au Canada (format Letter, 8,5 $\times$ 11 pouces).
\end{savaistu}

\paragraph{Démarche : régler le format et l'orientation au démarrage.}
Dès la création d'un nouveau fichier (Fichier $\rightarrow$ Nouveau), la boîte de dialogue \texttt{Nouvelle publication} propose :
\begin{enumerate}
    \item Choisir \texttt{Publication vierge} (pour partir de zéro) ou une catégorie de modèles.
    \item Cliquer sur \texttt{Format de page} (souvent un lien ou un bouton dans le volet latéral).
    \item Dans la boîte \texttt{Format de page} : sélectionner \texttt{A4}, \texttt{A5} ou \texttt{Personnalisé}, puis cocher \texttt{Paysage} ou \texttt{Portrait}.
    \item Valider par \texttt{OK}. La zone de dessin se redimensionne instantanément.
\end{enumerate}

\begin{exemplebox}[Calcul de la largeur utile pour un dépliant 3 volets A4 paysage]
Concevons un dépliant touristique (3 volets) sur une page A4 en orientation paysage.
\begin{itemize}
    \item Largeur totale A4 paysage : $29,7$ cm.
    \item Marges gauche et droite (imprimante standard) : $2 \times 1,5 = 3,0$ cm.
    \item Largeur utile totale : $29,7 - 3,0 = 26,7$ cm.
    \item 3 colonnes $\Rightarrow$ 2 gouttières (espaces entre colonnes) de $0,5$ cm chacune : $2 \times 0,5 = 1,0$ cm.
    \item Largeur disponible pour les 3 colonnes : $26,7 - 1,0 = 25,7$ cm.
    \item \textbf{Largeur utile par colonne} : $25,7 / 3 \approx 8,57$ cm (arrondi à $8,6$ cm).
\end{itemize}
\emph{Astuce :} les repères de colonnes se placent alors à $10,1$ cm, $10,6$ cm, $19,1$ cm et $19,6$ cm depuis le bord gauche de la page (marge de $1,5$ cm comprise).
\end{exemplebox}

\begin{popfigure}
\begin{tikzpicture}[scale=0.35, every node/.style={font=\small, align=center}]
    % Page A4 paysage : 29.7 x 21 cm ; marge 1.5 ; colonnes de 8.5667 ; gouttieres 0.5
    % Bordures de colonnes : x = 1.5 | 10.0667 | 10.5667 | 19.1333 | 19.6333 | 28.2
    % Bord de page
    \draw[popdark, thick] (0,0) rectangle (29.7,21);
    \node[popdark, anchor=north west] at (0.3,20.7) {Page A4 Paysage (29,7 $\times$ 21 cm)};

    % Marges (pointilles gris)
    \draw[popmuted, dashed, thick] (1.5,1.5) rectangle (28.2,19.5);
    \node[popmuted, anchor=south west, font=\tiny] at (1.6,1.6) {Marges 1,5 cm};

    % Reperes de colonnes (bleu)
    \draw[popblue, densely dashed, thick] (10.0667,1.5) -- (10.0667,19.5);
    \draw[popblue, densely dashed, thick] (10.5667,1.5) -- (10.5667,19.5);
    \draw[popblue, densely dashed, thick] (19.1333,1.5) -- (19.1333,19.5);
    \draw[popblue, densely dashed, thick] (19.6333,1.5) -- (19.6333,19.5);

    % Etiquettes des volets
    \node[popblue, align=center] at (5.7833,10.5){Volet 1\\Largeur $\approx$ 8,6 cm};
    \node[popblue, align=center] at (14.85,10.5){Volet 2\\Largeur $\approx$ 8,6 cm};
    \node[popblue, align=center] at (23.9167,10.5){Volet 3\\Largeur $\approx$ 8,6 cm};

    % Etiquettes des gouttieres (orange)
    \node[poporange, font=\tiny, rotate=90] at (10.3167,10.5) {Gouttière 0,5 cm};
    \node[poporange, font=\tiny, rotate=90] at (19.3833,10.5) {Gouttière 0,5 cm};

    % Reperes de pliage (petits traits en haut et en bas)
    \draw[popdark, thick] (10.0667,21) -- (10.0667,21.4);
    \draw[popdark, thick] (10.5667,21) -- (10.5667,21.4);
    \draw[popdark, thick] (19.1333,21) -- (19.1333,21.4);
    \draw[popdark, thick] (19.6333,21) -- (19.6333,21.4);
    \draw[popdark, thick] (10.0667,0) -- (10.0667,-0.4);
    \draw[popdark, thick] (10.5667,0) -- (10.5667,-0.4);
    \draw[popdark, thick] (19.1333,0) -- (19.1333,-0.4);
    \draw[popdark, thick] (19.6333,0) -- (19.6333,-0.4);
    \node[popdark, anchor=north, font=\tiny] at (14.85,-0.6) {Repères de pliage (traits de coupe)};
\end{tikzpicture}
\legende{Schéma d'une page A4 paysage configurée pour un dépliant 3 volets : marges en pointillés gris, repères de colonnes (limites des volets) en bleu, gouttières en orange, repères de pliage en haut et en bas.}
\end{popfigure}

\courssub{Utilisation d'un modèle : gain de temps et professionnalisme}

\paragraph{Pourquoi réinventer la roue ?} Les logiciels de PAO (Publisher, Scribus, Canva) proposent des \cle{modèles} (templates) préconçus : dépliants, cartes de visite, bulletins, menus, certificats. Un modèle fournit une structure prête à l'emploi : zones de texte positionnées, styles de police harmonisés, jeux de couleurs cohérents, images de substitution.

\begin{propriete}[Avantages et limites du modèle]
\begin{itemize}
    \item \textbf{Gain de temps immense} : la grille, les styles et l'alignement sont déjà faits.
    \item \textbf{Cohérence graphique} : respect des règles de design (alignement, contraste, répétition, proximité).
    \item \textbf{Personnalisation obligatoire} : un modèle n'est pas un document fini. Il faut \emph{remplacer} les textes factices, \emph{substituer} les images par les tiennes, \emph{adapter} les couleurs à ta charte (ex. : vert/jaune/rouge pour le Cameroun).
    \item \textbf{Piège :} ne pas modifier la structure (marges, colonnes) du modèle sans vérifier l'impact sur l'ensemble du document.
\end{itemize}
\end{propriete}

\begin{methode}[Créer un dépliant 3 volets à partir d'un modèle (ou d'une page blanche)]
\textbf{Objectif :} produire une structure prête pour un dépliant touristique (ex. : « Découvrir le Nord-Cameroun »).
\begin{enumerate}
    \item \textbf{Nouveau document} : Fichier $\rightarrow$ Nouveau $\rightarrow$ taper « Dépliant » dans la zone de recherche, ou choisir \texttt{Publication vierge A4}.
    \item \textbf{Format} : si page blanche $\rightarrow$ Format de page $\rightarrow$ A4 $\rightarrow$ \textbf{Paysage} $\rightarrow$ OK.
    \item \textbf{Marges} : onglet \texttt{Mise en page} $\rightarrow$ \texttt{Marges} $\rightarrow$ \texttt{Marges personnalisées} $\rightarrow$ Haut/Bas/Gauche/Droite = 1.5cm (ou 2 cm pour plus de sécurité avec l'imprimante).
    \item \textbf{Repères de colonnes} : onglet \texttt{Mise en page} $\rightarrow$ \texttt{Repères} $\rightarrow$ \texttt{Repères de colonnes et de lignes} $\rightarrow$ Nombre de colonnes : \textbf{3} $\rightarrow$ Espacement (gouttière) : \textbf{0,5 cm} $\rightarrow$ OK.
    \item \textbf{Affichage des repères} : onglet \texttt{Affichage} $\rightarrow$ cocher \texttt{Repères} et \texttt{Règles} pour visualiser la grille et mesurer les positions.
    \item \textbf{Enregistrement immédiat} : Fichier $\rightarrow$ Enregistrer sous $\rightarrow$ dossier \texttt{PAO\_3eme} $\rightarrow$ Nom : \texttt{Depliant\_Nord\_Cameroun.pub} $\rightarrow$ Enregistrer.
    \item \textbf{Remplissage} : cliquer dans chaque zone de texte du modèle (ou créer des zones de texte dans les colonnes) $\rightarrow$ taper tes titres et textes, insérer tes photos.
\end{enumerate}
\end{methode}

\courssub{Réglage des marges et création de repères de colonnes}

\paragraph{Les marges : une zone de sécurité.} La \cle{marge} est l'espace blanc entre le bord physique du papier et la zone utile. Elle sert à trois choses : (1) éviter que l'imprimante ne coupe du contenu (marge technique d'au moins 0,5 cm, souvent 1 cm), (2) laisser de l'air pour la lecture (confort visuel), (3) prévoir l'espace pour la reliure ou le pliage (marge intérieure plus large pour une brochure agrafée).

\begin{attention}[Erreur fréquente : marges trop fines ou imprimante oubliée]
Beaucoup d'élèves mettent des marges à 0,5 cm « pour gagner de la place ». \textbf{Résultat :} l'imprimante du cybercafé ou du lycée coupe le bas du texte ou le logo. \textbf{Règle d'or :} marge minimum de \textbf{1,5 cm} sur les 4 côtés pour un document standard, \textbf{2 cm} côté reliure ou pliage central. Vérifie toujours les marges imprimables de l'imprimante cible (Fichier $\rightarrow$ Imprimer $\rightarrow$ Propriétés de l'imprimante).
\end{attention}

\paragraph{Les repères de colonnes (guides) : la grille invisible.} Dans un dépliant 3 volets, les repères verticaux matérialisent les bords de chaque volet. Dans un bulletin à 2 colonnes, ils séparent les articles. Ils sont \emph{non imprimables} : ils ne sortent pas sur le papier, ils guident seulement ton œil et l'alignement des objets (aimantation automatique des cadres sur les repères).

\begin{propriete}[Créer des repères de colonnes personnalisés]
\begin{enumerate}
    \item Onglet \texttt{Mise en page} $\rightarrow$ groupe \texttt{Repères} $\rightarrow$ \texttt{Repères de colonnes et de lignes...}
    \item Dans la boîte de dialogue : \texttt{Nombre de colonnes} = 3 (pour 3 volets), \texttt{Espacement} = 0.5cm (gouttière).
    \item Valider. Les lignes de repère (bleues par défaut) apparaissent sur la page.
    \item Pour déplacer un repère manuellement : maintenir \texttt{Maj} et glisser le repère avec la souris.
\end{enumerate}
\end{propriete}

\courssub{Gestion des pages : ajout, suppression, déplacement, duplication}

Une publication PAO est souvent \textbf{multipage} : recto/verso d'un dépliant (2 pages), bulletin de 4 pages, catalogue de 8 pages. Le \cle{volet des pages} (généralement à gauche de l'écran) est ton tableau de bord.

\begin{experience}[TP guidé : manipuler les pages d'un dépliant recto-verso]
\textbf{Objectif :} créer un dépliant 3 volets recto-verso (2 pages physiques = 6 volets logiques) et maîtriser la navigation.
\textbf{Logiciel :} Microsoft Publisher (ou Scribus, LibreOffice Draw).
\textbf{Manipulation :}
\begin{enumerate}
    \item Ouvre le fichier \texttt{Depliant\_Nord\_Cameroun.pub} créé précédemment. Tu as 1 page (le recto).
    \item \textbf{Ajouter le verso} : dans le volet Pages, clic droit sur la vignette de la page 1 $\rightarrow$ \texttt{Insérer une page} $\rightarrow$ \texttt{Après la page actuelle} $\rightarrow$ Nombre : 1 $\rightarrow$ OK. \emph{Observation :} une page 2 apparaît, vide, avec la même structure (marges, colonnes).
    \item \textbf{Dupliquer la structure} : au lieu d'insérer une page vierge, clic droit sur la page 1 $\rightarrow$ \texttt{Dupliquer la page}. \emph{Observation :} la page 2 contient \emph{exactement} les mêmes cadres, repères et objets que la page 1. C'est l'idéal pour le verso d'un dépliant : les plis tombent au même endroit.
    \item \textbf{Déplacer} : glisse-dépose la vignette de la page 2 avant la page 1 dans le volet. L'ordre change.
    \item \textbf{Supprimer} : clic droit sur une page vide $\rightarrow$ \texttt{Supprimer la page}.
\end{enumerate}
\textbf{Interprétation :} la duplication préserve la \textbf{cohérence géométrique} (plis alignés recto/verso). L'insertion d'une page vierge est utile pour ajouter une page différente (ex. : page de garde).
\end{experience}

\courssub{Enregistrement de la publication : formats et organisation}

\paragraph{Le réflexe vital : Ctrl+S.} En PAO, les fichiers grossissent vite (images haute résolution intégrées). Une coupure de courant ou un plantage du logiciel, et tout est perdu si rien n'est sauvegardé. \textbf{Enregistre après chaque étape majeure} (création de la page, insertion d'une image, fin d'une colonne).

\begin{definition}[Fichier source et fichier d'échange]
\begin{itemize}
    \item \textbf{Format natif} (\texttt{.pub} pour Publisher, \texttt{.sla} pour Scribus) : conserve \emph{toutes} les fonctionnalités (styles, repères, objets modifiables). C'est ton \textbf{fichier de travail}. \textbf{Ne l'envoie jamais à l'imprimeur ni au destinataire final.}
    \item \textbf{Format d'échange / impression (PDF)} : figé, lisible sur tout ordinateur, universel. C'est ce qu'on fournit pour l'impression ou la diffusion.
    \item \textbf{Format image (JPG, PNG)} : pour un aperçu écran ou une insertion dans une page web, pas pour l'impression professionnelle.
\end{itemize}
\end{definition}

\begin{popfigure}
\begin{tikzpicture}[
    node distance=1.1cm and 1.5cm,
    every node/.style={font=\small, align=center},
    menu/.style={rectangle, rounded corners=3pt, draw=popblue, thick, fill=popblueL, minimum width=3cm, minimum height=0.8cm},
    action/.style={rectangle, rounded corners=3pt, draw=popgreen, thick, fill=popgreenL, minimum width=3.5cm, minimum height=0.8cm},
    file/.style={rectangle, rounded corners=3pt, draw=poporange, thick, fill=poporangeL, minimum width=4cm, minimum height=0.8cm},
    arrow/.style={-Stealth, thick, popdark}
]
    % Noeuds
    \node[menu] (fichier) {Fichier};
    \node[menu, below of=fichier] (enrsous) {Enregistrer sous...};
    \node[file, below of=enrsous, yshift=-0.4cm, align=center] (boite){Boîte de dialogue\\« Enregistrer sous »};

    \node[action, left of=boite, xshift=-2.2cm, align=center] (dossier){Choisir le dossier\\(ex. PAO\_3eme)};
    \node[action, right of=boite, xshift=2.2cm, align=center] (nom){Saisir le nom\\(ex. Depliant\_Nord)};
    \node[action, below of=boite, yshift=-0.8cm, align=center] (format){Choisir le format\\(.pub / .pdf / .jpg)};
    \node[menu, below of=format, yshift=-0.4cm] (btn) {Bouton [Enregistrer]};

    % Fleches
    \draw[arrow] (fichier) -- (enrsous);
    \draw[arrow] (enrsous) -- (boite);
    \draw[arrow] (boite) -- (dossier);
    \draw[arrow] (boite) -- (nom);
    \draw[arrow] (boite) -- (format);
    \draw[arrow] (format) -- (btn);
    \draw[arrow] (dossier) |- (btn);
    \draw[arrow] (nom) |- (btn);
\end{tikzpicture}
\legende{Arborescence du menu Fichier $\rightarrow$ Enregistrer sous : navigation vers le dossier de travail, nommage explicite (sans espaces ni accents), choix du format natif (.pub) pour la sauvegarde de travail.}
\end{popfigure}

\begin{methode}[Enregistrer correctement sa publication PAO]
\begin{enumerate}
    \item \textbf{Raccourci} : \texttt{Ctrl + S} (Enregistrer) si le fichier a déjà un nom ; Fichier $\rightarrow$ \texttt{Enregistrer sous} pour la première fois ou pour créer une nouvelle version.
    \item \textbf{Dossier} : crée un dossier dédié par projet (ex. \texttt{Documents\textbackslash PAO\_3eme\textbackslash Depliant\_Nord\_Cameroun\textbackslash}). Range-y le fichier \texttt{.pub} \textbf{et} les images utilisées.
    \item \textbf{Nom de fichier} : sans espaces, sans accents, explicite : \texttt{Depliant\_Nord\_Cameroun\_v1.pub}. Utilise les suffixes \texttt{\_v1}, \texttt{\_v2} pour gérer les versions.
    \item \textbf{Format} : sélectionne \textbf{Publication Publisher (*.pub)} (ou le format natif de ton logiciel) pour le travail ; exporte en \textbf{PDF} (Fichier $\rightarrow$ Exporter $\rightarrow$ Créer un document PDF) pour la diffusion ou l'impression.
    \item \textbf{Copie de sécurité} : avant de remettre ton travail sur clé USB, copie le fichier \texttt{.pub} \textbf{et toutes les images} dans un même dossier, afin que rien ne manque à l'ouverture sur un autre poste.
\end{enumerate}
\end{methode}

\courssub{Démarche de production et justification}

\paragraph{Synthèse : de l'intention à la structure verrouillée.} Produire une publication PAO, ce n'est pas « faire joli » au feeling. C'est une démarche méthodique :
\begin{enumerate}
    \item \textbf{Analyse du besoin} : quel objet ? (dépliant, affiche, carte) $\rightarrow$ détermine le format et l'orientation.
    \item \textbf{Choix du support} : modèle (gain de temps) ou page blanche (liberté totale) ?
    \item \textbf{Architecture de la page} : marges (sécurité), colonnes et repères (structure), règles (alignement).
    \item \textbf{Gestion du document} : pages multiples (recto/verso), duplication (cohérence), nommage (traçabilité).
    \item \textbf{Sauvegarde continue} : \texttt{Ctrl+S} (source \texttt{.pub}) puis export PDF (livraison).
\end{enumerate}

\begin{exoresolu}[Structurer une invitation pour la Journée Portes Ouvertes du lycée]
\textbf{Énoncé :} tu dois créer une invitation format carte postale (A6 : 10,5 $\times$ 14,8 cm), recto seul, avec marge de sécurité de 1 cm, une colonne unique, image de fond sur toute la page. Donne la procédure complète de création du fichier source.
\tcblower
\textbf{Solution détaillée :}
\begin{enumerate}
    \item \textbf{Analyse} : objet = invitation. Format = A6 (quart d'un A4). Orientation = \textbf{Paysage} (14,8 cm de large $\times$ 10,5 cm de haut), fréquent pour une carte postale. Une seule colonne. Image de fond = cadre couvrant toute la page.
    \item \textbf{Création} : Fichier $\rightarrow$ Nouveau $\rightarrow$ Page blanche $\rightarrow$ Format de page $\rightarrow$ Personnalisé : largeur 14,8 cm, hauteur 10,5 cm $\rightarrow$ Orientation Paysage $\rightarrow$ OK.
    \item \textbf{Marges} : Mise en page $\rightarrow$ Marges $\rightarrow$ Personnalisées $\rightarrow$ 1 cm sur les 4 côtés. (Zone utile : 12,8 $\times$ 8,5 cm.)
    \item \textbf{Repères} : pas de colonnes nécessaires (un seul bloc). Activer l'affichage des repères de marge (Affichage $\rightarrow$ Repères).
    \item \textbf{Image de fond} : Insertion $\rightarrow$ Image $\rightarrow$ choisir le fichier $\rightarrow$ étirer l'image pour couvrir toute la page $\rightarrow$ clic droit $\rightarrow$ \texttt{Mettre à l'arrière-plan}.
    \item \textbf{Zones de texte} : dessiner des zones de texte \emph{à l'intérieur} des marges (zone utile) pour le titre, la date et le lieu.
    \item \textbf{Enregistrement} : Fichier $\rightarrow$ Enregistrer sous $\rightarrow$ dossier \texttt{PAO\_3eme\textbackslash Invitation\_JPO} $\rightarrow$ Nom : \texttt{Invitation\_JPO\_2026\_v1.pub} $\rightarrow$ Type : \texttt{.pub} $\rightarrow$ Enregistrer.
    \item \textbf{Justification} : le format A6 est standard en imprimerie (découpe facile depuis une feuille A4). La marge de 1 cm garantit la sécurité à la coupe et la lisibilité. L'image en arrière-plan ne gêne pas l'édition du texte. Le fichier \texttt{.pub} reste modifiable pour l'année suivante (il suffira de changer la date).
\end{enumerate}
\end{exoresolu}

\begin{exercice}[Mise en situation : bulletin d'information du club informatique]
Le club informatique du lycée veut un bulletin de 4 pages (feuille A4 pliée en 2 = 4 pages A5) : page 1 (couverture), pages 2-3 (articles sur 2 colonnes), page 4 (contacts).
\begin{enumerate}
    \item Quels sont le format de page et l'orientation à choisir dans le logiciel ? Justifie.
    \item Combien de pages physiques faut-il créer dans le volet Pages ? Explique la différence entre page physique et page logique (volet).
    \item Quelle méthode utiliser pour que les pages 2 et 3 aient exactement la même structure de 2 colonnes ? Détaille les clics.
    \item Donne le nom de fichier recommandé (convention v1) et le format de sauvegarde de travail.
\end{enumerate}
\end{exercice}

\begin{corrige}[Exercice : bulletin d'information du club informatique]
\begin{enumerate}
    \item \textbf{Format : A4, orientation Portrait.} Justification : le produit final est une feuille A4 pliée en deux (format fini A5). En PAO, on travaille \emph{à plat} sur la feuille d'impression (A4) ; le pli central sépare les pages logiques. L'orientation portrait convient à un livret standard plié verticalement.
    \item \textbf{2 pages physiques.} La page physique 1 (recto de la feuille) porte les pages logiques 4 (à gauche) et 1 (à droite) ; la page physique 2 (verso) porte les pages logiques 2 (à gauche) et 3 (à droite). La \textbf{page physique} est la feuille réelle qui sort de l'imprimante ; la \textbf{page logique} est la page du bulletin telle que le lecteur la lit après pliage.
    \item \textbf{Duplication.} Créer la page 1 (couverture). Insérer une page 2 vierge. Sur la page 2 : Mise en page $\rightarrow$ Repères $\rightarrow$ Repères de colonnes $\rightarrow$ 2 colonnes, gouttière 0,5 cm $\rightarrow$ OK. Puis \textbf{clic droit sur la vignette de la page 2 $\rightarrow$ Dupliquer la page} : on obtient une page 3 avec exactement la même structure de colonnes.
    \item \textbf{Nom :} \texttt{Bulletin\_Club\_Info\_Mars2026\_v1.pub}. \textbf{Format :} le format natif du logiciel (\texttt{.pub} pour Publisher) pour conserver styles, repères et objets modifiables ; l'export PDF sera réservé à la diffusion finale.
\end{enumerate}
\end{corrige}

\begin{aretenir}[Check-list avant de passer au contenu (section 4)]
Avant d'insérer le moindre titre ou la moindre photo, vérifie :
\begin{itemize}
    \item [$\square$] Le format de page (A4, A5, personnalisé) et l'orientation (Portrait/Paysage) correspondent au produit final.
    \item [$\square$] Les marges sont d'au moins 1,5 cm (2 cm côté pliure/reliure).
    \item [$\square$] Les repères de colonnes sont posés (nombre, gouttière) pour structurer l'espace.
    \item [$\square$] Les pages nécessaires sont créées (ajout/duplication) et ordonnées.
    \item [$\square$] Le fichier source est enregistré (\texttt{.pub}) dans un dossier de projet nommé, avec version (\texttt{\_v1}).
    \item [$\square$] \texttt{Ctrl+S} devient un réflexe.
\end{itemize}
\end{aretenir}

\coursec{Insertion et mise en forme des objets}

Dans la section précédente, nous avons créé la \cle{page} de notre publication : choix du format, des marges, de l'orientation et de la structure (colonnes, repères). Mais une page vide, même bien structurée, ne communique rien. Il faut maintenant la \emph{remplir} : y placer du texte, des images, des formes décoratives, puis les organiser harmonieusement. C'est exactement ce que fait l'élève chargé de concevoir l'affiche de la semaine culturelle du collège : une fois la page préparée, il insère le titre, le texte d'annonce, la photo du groupe de danse et le logo de l'établissement.

La grande différence avec un logiciel de traitement de texte est ici fondamentale : en \cle{PAO}, rien ne se place « au fil du texte ». Chaque contenu est un \cle{objet} indépendant (zone de texte, image, forme) que l'on pose librement sur la page, que l'on déplace, redimensionne, superpose et aligne. Cette section t'apprend à insérer ces objets et à les mettre en forme avec goût et rigueur.

\courssub{Insertion d'une zone de texte}

\textbf{Pourquoi une « zone » de texte ?} Dans un traitement de texte, le texte coule automatiquement de la première à la dernière page. En PAO, le concepteur veut décider \emph{précisément} où chaque bloc de texte apparaît : un titre en haut, un encadré à droite, une légende sous une photo. La solution est la \cle{zone de texte} : un cadre rectangulaire, positionné librement, qui contient du texte.

\begin{definition}[Zone de texte]
Une \cle{zone de texte} est un objet rectangulaire placé librement sur la page d'une publication et destiné à contenir du texte. Elle possède des \cle{poignées} (petits carrés aux coins et au milieu des côtés) qui apparaissent quand elle est sélectionnée et permettent de la redimensionner ou de la déplacer.
\end{definition}

\begin{methode}[Insérer et mettre en forme une zone de texte en 5 étapes]
\begin{enumerate}
\item \textbf{Insérer} : ouvrir l'onglet \texttt{Insertion}, cliquer sur \texttt{Zone de texte} (ou \texttt{Zone de texte dessinée}).
\item \textbf{Tracer} : sur la page, cliquer au coin supérieur gauche souhaité, maintenir le bouton de la souris enfoncé et faire glisser jusqu'au coin inférieur droit, puis relâcher. Le cadre apparaît avec un curseur clignotant à l'intérieur.
\item \textbf{Saisir} : taper le texte directement dans la zone (ou le coller depuis un autre document).
\item \textbf{Mettre en forme} : sélectionner le texte, puis dans l'onglet \texttt{Accueil} choisir la police, la taille, la couleur et l'alignement.
\item \textbf{Ajuster} : cliquer sur le bord du cadre pour sélectionner la zone entière, puis la déplacer (glisser-déposer) ou la redimensionner avec les poignées jusqu'à ce que tout le texte soit visible.
\end{enumerate}
\end{methode}

\begin{attention}[Texte qui déborde de la zone]
Si le texte est plus long que la zone, il devient \emph{invisible} : la zone affiche parfois un petit symbole d'avertissement. Deux solutions : agrandir la zone avec les poignées, ou réduire la taille de la police. Ne laisse jamais une partie du texte cachée dans une publication finale !
\end{attention}

\courssub{Mise en forme du texte : les choix raisonnés}

Une fois le texte saisi, on le met en forme avec les outils classiques déjà connus du traitement de texte :

\begin{itemize}
\item la \cle{police} (Arial, Times New Roman, Calibri\ldots) : une police \emph{sans empattements} (Arial) convient aux titres, une police \emph{avec empattements} (Times) au corps du texte ;
\item la \cle{taille} : un titre entre 24 et 48 points, un texte courant entre 10 et 12 points ;
\item la \cle{couleur} : à choisir lisible sur le fond de la page (jamais de jaune clair sur fond blanc !) ;
\item l'\cle{alignement} : gauche, centré, droit ou justifié ;
\item l'\cle{interligne} : l'espace vertical entre les lignes (simple, 1,15, 1,5\ldots).
\end{itemize}

Mais la PAO est un art de la \emph{sobriété}. Une affiche qui mélange six polices et huit couleurs fatigue l'œil et paraît désordonnée : le lecteur ne sait plus où regarder. La pratique professionnelle a dégagé une règle simple.

\begin{propriete}[Règle de sobriété typographique]
Un document de publication lisible et professionnel utilise \textbf{au plus 3 polices} et \textbf{au plus 3 couleurs dominantes}. En pratique : une police pour les titres, une pour le texte courant, éventuellement une troisième pour un élément spécial (date, slogan) ; et une palette de deux à trois couleurs en harmonie avec les images du document.
\end{propriete}

Cette règle n'est pas un théorème mathématique : elle est \emph{admise} car elle est justifiée par l'expérience. L'œil humain repère la structure d'un document grâce à la répétition : si tous les titres ont la même police et la même couleur, le lecteur les identifie instantanément. Multiplier les styles détruit ce repère.

\begin{exemplebox}[Corriger une affiche surchargée]
Une élève a conçu l'affiche de la kermesse avec 6 polices différentes (Comic Sans, Algerian, Arial, Times, Impact, Broadway) et 8 couleurs (rouge, vert, bleu, jaune, violet, orange, rose, marron). Résultat : illisible. On applique la règle de sobriété :
\begin{itemize}
\item \textbf{Titres} : police Impact, taille 36, couleur bleu foncé ;
\item \textbf{Texte courant} : police Arial, taille 12, couleur noire ;
\item \textbf{Date et lieu} : Arial, taille 14, gras, couleur rouge (seule touche vive, pour attirer l'œil sur l'information essentielle).
\end{itemize}
On passe ainsi de 6 polices et 8 couleurs à \textbf{2 polices et 3 couleurs} : l'affiche devient claire et professionnelle.
\end{exemplebox}

\courssub{Insertion et redimensionnement d'une image}

Une image illustre, attire l'attention, explique mieux qu'un long paragraphe. Pour l'insérer : onglet \texttt{Insertion} $\rightarrow$ \texttt{Images} (pour un fichier enregistré sur l'ordinateur, par exemple une photo du club informatique) ou \texttt{Images en ligne} (pour rechercher une illustration sur Internet, en respectant les droits d'auteur).

Une fois insérée, l'image est un objet comme la zone de texte : elle possède des poignées. C'est ici que se cache le piège le plus fréquent des débutants.

\begin{attention}[Ne jamais déformer une image]
Les poignées situées \emph{au milieu des côtés} étirent l'image dans un seul sens : tirer la poignée latérale droite élargit la photo et rend les personnages « écrasés » ou « allongés ». Pour agrandir ou réduire une image \textbf{en conservant ses proportions}, il faut \textbf{toujours utiliser les poignées d'angle} (les quatre coins). On parle de \cle{redimensionnement proportionnel}. Astuce : maintenir la touche \texttt{Maj} enfoncée pendant le glisser garantit la proportionnalité dans de nombreux logiciels.
\end{attention}

\begin{exemplebox}[Redimensionnement proportionnel : exemple chiffré]
Une photo mesure 10 cm de large et 6 cm de haut. Son rapport largeur/hauteur vaut $\dfrac{10}{6} \approx 1,67$. On veut la réduire à 5 cm de large. Pour conserver les proportions, la hauteur doit vérifier :
\[ \frac{\text{largeur}}{\text{hauteur}} = \frac{10}{6} \quad \Longrightarrow \quad \text{hauteur} = \frac{5 \times 6}{10} = 3 \text{ cm}. \]
En tirant une poignée d'angle, le logiciel effectue ce calcul automatiquement : l'image passe de $10 \times 6$ cm à $5 \times 3$ cm sans déformation.
\end{exemplebox}

\courssub{Habillage du texte et ordre de superposition}

Quand une image et une zone de texte se chevauchent, deux questions se posent : comment le texte contourne-t-il l'image ? Et quel objet apparaît \emph{devant} ?

\begin{definition}[Habillage]
L'\cle{habillage} est le mode de répartition du texte autour d'une image. Les modes courants sont : \emph{habillé} (le texte épouse le contour de l'image), \emph{encadré} (le texte passe au-dessus et en dessous), \emph{derrière le texte} (l'image sert de fond) et \emph{devant le texte} (l'image recouvre le texte).
\end{definition}

Chaque objet posé sur la page occupe un \emph{plan}, comme des feuilles transparentes empilées. On modifie cet empilement par clic droit sur l'objet : \texttt{Mettre au premier plan}, \texttt{Mettre à l'arrière-plan}, \texttt{Avancer}, \texttt{Reculer}. Ainsi, un bandeau de couleur en arrière-plan peut servir de décor à un titre placé au premier plan.

\begin{popfigure}
\begin{tikzpicture}[scale=0.95, every node/.style={font=\small}]
% Zone de texte avec poignées
\draw[thick, popblue, fill=popblueL] (0,0) rectangle (5.4,3.2);
\node[align=left, text width=4.6cm, text=popdark] at (2.7,1.6)
 {Semaine culturelle du coll\`ege : danses traditionnelles, th\'e\^atre et concerts vous attendent du 12 au 16 mai\ldots};
% Poignées
\foreach \x/\y in {0/0, 2.7/0, 5.4/0, 0/1.6, 5.4/1.6, 0/3.2, 2.7/3.2, 5.4/3.2}
  \draw[fill=white, draw=popdark] (\x,\y) circle (0.09);
\draw[stealth-stealth, poporange, thick] (5.4,3.2) -- (6.3,4.1);
\node[poporange, font=\footnotesize, align=center] at (7.3,4.35) {poign\'ee d'angle :\\ redimensionner};
\node[popblue, font=\footnotesize] at (2.7,-0.45) {zone de texte s\'electionn\'ee (8 poign\'ees)};
% Image avec habillage
\begin{scope}[xshift=8.2cm]
\draw[thick, popgreen, fill=popgreenL] (0,0) rectangle (4.6,3.2);
\draw[fill=poppurpleL, draw=poppurple, thick] (0.35,0.9) rectangle (2.0,2.9);
\draw[fill=poppinkL, draw=poppink] (1.17,2.1) circle (0.4);
\draw[fill=popgoldL, draw=popgold] (0.35,0.9) -- (1.0,1.9) -- (1.5,1.3) -- (2.0,2.2) -- (2.0,0.9) -- cycle;
\node[font=\footnotesize, popdark] at (1.17,0.55) {image};
\foreach \x/\y in {0.35/0.9, 2.0/0.9, 0.35/2.9, 2.0/2.9}
  \draw[fill=white, draw=popdark] (\x,\y) circle (0.08);
\node[align=left, text width=2.1cm, font=\footnotesize, text=popdark] at (3.35,2.35) {Le texte contourne l'image : c'est l'habillage.};
\node[align=left, text width=4.0cm, font=\footnotesize, text=popdark] at (2.3,0.35) {Le texte continue sous l'image.};
\end{scope}
\end{tikzpicture}
\legende{À gauche : une zone de texte sélectionnée et ses huit poignées ; la poignée d'angle permet de redimensionner. À droite : une image insérée avec habillage du texte autour d'elle.}
\end{popfigure}

\begin{popfigure}
\begin{tikzpicture}[scale=1, every node/.style={font=\small}]
% Vue éclatée : trois plans
\draw[thick, popgold, fill=popgoldL, opacity=0.9] (0,0) rectangle (4.4,2.8);
\node[popdark] at (2.2,1.4) {bandeau de couleur};
\node[popgold, font=\footnotesize, align=center] at (2.2,-0.5) {arri\`ere-plan\\ (pos\'e en premier)};
\begin{scope}[xshift=1.6cm, yshift=1.5cm]
\draw[thick, poppurple, fill=poppurpleL, opacity=0.92] (0,0) rectangle (4.4,2.8);
\node[popdark] at (2.2,1.4) {photo du groupe de danse};
\node[poppurple, font=\footnotesize, align=center] at (2.2,-0.5) {plan interm\'ediaire};
\end{scope}
\begin{scope}[xshift=3.2cm, yshift=3.0cm]
\draw[thick, popblue, fill=popblueL, opacity=0.92] (0,0) rectangle (4.4,2.8);
\node[popdark, align=center] at (2.2,1.4) {titre : \og Semaine culturelle \fg};
\node[popblue, font=\footnotesize, align=center] at (2.2,-0.5) {premier plan\\ (pos\'e en dernier)};
\end{scope}
% Flèches d'empilement
\draw[-stealth, very thick, poporange] (9.6,1.2) -- (9.6,3.4);
\node[poporange, font=\footnotesize, align=left] at (11.3,2.3) {Mettre au\\ premier plan};
\draw[-stealth, very thick, popmuted] (10.6,3.4) -- (10.6,1.2);
\node[popmuted, font=\footnotesize, align=left] at (12.6,1.6) {Mettre \`a l'arri\`ere-plan};
\end{tikzpicture}
\legende{Vue éclatée de la superposition des objets : chaque objet occupe un plan. Les commandes « premier plan » et « arrière-plan » modifient l'ordre d'empilement.}
\end{popfigure}

\courssub{Formes, bordures et effets : la sobriété avant tout}

Pour structurer la page, on peut insérer des \cle{formes} (rectangles, flèches, étoiles, bulles de bande dessinée) via l'onglet \texttt{Insertion} $\rightarrow$ \texttt{Formes}. Chaque forme possède :

\begin{itemize}
\item un \cle{remplissage} : couleur unie, \cle{dégradé} (transition douce entre deux couleurs), motif ou image ;
\item une \cle{bordure} : couleur, épaisseur et style du contour (plein, pointillés) ;
\item des \cle{effets} : \cle{ombre} portée, lueur, biseautage.
\end{itemize}

\begin{attention}[L'excès d'effets tue l'effet]
Une ombre légère sous un titre peut lui donner du relief ; mais multiplier ombres, dégradés criards et bordures épaisses sur tous les objets rend la publication illisible et « amateur ». Retiens la règle : \textbf{un effet doit souligner une information, jamais la décorer gratuitement}. En cas de doute, n'en mets pas.
\end{attention}

\courssub{Alignement, distribution et groupement des objets}

Une page professionnelle se reconnaît à ses alignements : les bords gauches des zones de texte suivent la même verticale, les images sont régulièrement espacées. Deux outils nous y aident.

\textbf{Les repères et guides.} Des lignes vertes (ou des lignes magnétiques) apparaissent automatiquement quand un objet s'aligne avec un autre pendant le déplacement : le logiciel « accroche » l'objet à la bonne position.

\textbf{Les outils d'alignement.} Après avoir sélectionné plusieurs objets (cliquer sur chacun en maintenant la touche \texttt{Maj}), l'onglet \texttt{Format} propose : \texttt{Aligner à gauche}, \texttt{Aligner au centre}, \texttt{Aligner en haut}\ldots{} ainsi que \texttt{Distribuer horizontalement} et \texttt{Distribuer verticalement}, qui répartissent les objets à \emph{égale distance} les uns des autres.

\begin{definition}[Groupement]
Le \cle{groupement} réunit plusieurs objets en un seul objet composite : sélectionner les objets, puis clic droit $\rightarrow$ \texttt{Grouper}. Le groupe se déplace, se redimensionne et s'aligne comme un tout. La commande \texttt{Dissocier} (ou \texttt{Dégrouper}) inverse l'opération.
\end{definition}

C'est indispensable pour un logo composé d'une forme, d'une image et d'un texte : groupés, ces trois éléments restent ensemble quand on les déplace sur la page.

\begin{experience}[TP guidé : composer l'encart « Journée portes ouvertes »]
\textbf{Objectif} : produire un encart d'annonce combinant texte, image et forme, correctement alignés et groupés.\par
\textbf{Matériel} : un ordinateur avec un logiciel de PAO (Microsoft Publisher ou Scribus), une photo de l'établissement.\par
\textbf{Manipulation} :
\begin{enumerate}
\item Insérer une zone de texte en haut de page, saisir « Journée portes ouvertes — Samedi 21 juin », police Arial, 28 pt, bleu foncé, centré.
\item Insérer la photo (onglet \texttt{Insertion} $\rightarrow$ \texttt{Images}), la réduire à 8 cm de large \emph{par une poignée d'angle}, vérifier que les proportions sont conservées.
\item Régler l'habillage de l'image sur « encadré » et la placer sous le titre.
\item Insérer une forme « rectangle arrondi » sous l'image, remplissage dégradé bleu clair, sans ombre ; y taper le texte « Entrée libre à partir de 9 h ».
\item Sélectionner les trois objets (touche \texttt{Maj}), appliquer \texttt{Aligner au centre}, puis \texttt{Grouper}.
\item Déplacer le groupe au centre de la page à l'aide des repères.
\end{enumerate}
\textbf{Observations} : l'image n'est pas déformée ; les trois objets partagent le même axe vertical ; le groupe se déplace d'un seul bloc.\par
\textbf{Interprétation} : l'alignement et le groupement garantissent une mise en page nette et facile à réorganiser, qualité essentielle d'une publication professionnelle.
\end{experience}

\begin{exoresolu}[Organiser les objets d'une affiche]
\textbf{Énoncé.} Sur une affiche, une photo recouvre la moitié du titre, et trois photos d'artistes sont placées à des hauteurs différentes, ce qui donne un rendu désordonné. (a) Comment rendre le titre entièrement visible ? (b) Comment aligner les trois photos sur la même ligne avec des espacements égaux ? (c) Pourquoi faut-il éviter d'élargir la photo principale en tirant sa poignée latérale droite ?
\tcblower
\textbf{Solution.} (a) Cliquer sur la photo, puis clic droit $\rightarrow$ \texttt{Mettre à l'arrière-plan} (ou sélectionner le titre et choisir \texttt{Mettre au premier plan}) : le titre passe devant et redevient lisible.\par
(b) Sélectionner les trois photos avec la touche \texttt{Maj}, puis dans l'onglet \texttt{Format} : \texttt{Aligner en haut} pour les mettre sur la même ligne horizontale, puis \texttt{Distribuer horizontalement} pour obtenir des espacements égaux.\par
(c) La poignée latérale modifie la largeur sans toucher la hauteur : le rapport largeur/hauteur change et l'image est déformée (personnages écrasés). Il faut utiliser une poignée d'angle, qui conserve les proportions.
\end{exoresolu}

\begin{aretenir}[L'essentiel de la section]
\begin{itemize}
\item En PAO, tout contenu est un \cle{objet} libre : zone de texte, image, forme.
\item Une zone de texte s'insère par \texttt{Insertion} $\rightarrow$ \texttt{Zone de texte}, puis se trace à la souris.
\item Une image se redimensionne \textbf{toujours par les poignées d'angle} pour conserver ses proportions.
\item L'\cle{habillage} gère le texte autour de l'image ; les commandes \emph{premier plan / arrière-plan} gèrent la superposition.
\item Sobriété : au plus 3 polices, 3 couleurs dominantes, et des effets avec parcimonie.
\item Les outils \emph{Aligner} et \emph{Distribuer} ordonnent la page ; le \cle{groupement} fait de plusieurs objets un seul bloc.
\end{itemize}
\end{aretenir}

\begin{exercice}[À toi de jouer]
Tu conçois le bulletin d'information du club informatique. (1) Décris les étapes pour insérer le titre dans une zone de texte et le mettre en forme (police, taille, couleur, alignement) en respectant la règle de sobriété. (2) Tu insères deux photos côte à côte : explique comment les aligner parfaitement et les espacer régulièrement. (3) Le logo du club est composé d'une étoile et d'un texte : quelle commande utilises-tu pour les déplacer ensemble, et pourquoi ?
\end{exercice}

\coursec{Finalisation, exportation et impression}

Dans la section précédente, tu as appris à insérer et à mettre en forme les objets de ta publication : zones de texte, images, formes, couleurs. Ton affiche ou ton dépliant est maintenant terminé à l'écran. Mais attention : un document qui semble parfait à l'écran peut réserver de mauvaises surprises une fois imprimé — un texte coupé, une image floue, une faute d'orthographe en gros titre. Cette dernière étape, souvent négligée, est pourtant celle qui fait la différence entre un travail d'amateur et un travail de professionnel. Nous allons apprendre à \cle{vérifier}, \cle{exporter} et \cle{imprimer} correctement une publication.

\courssub{Vérification avant impression}

\textbf{Pourquoi vérifier ?}\par
Imprimer coûte cher : du papier, de l'encre, du temps. Au cybercafé de ton quartier, chaque page imprimée par erreur est de l'argent perdu. Avant d'imprimer, il faut donc s'assurer que la publication est parfaite. Trois vérifications sont indispensables.

\textbf{1. L'aperçu avant impression.}\par
L'\cle{aperçu avant impression} montre à l'écran exactement ce qui va sortir de l'imprimante : les marges, la position des objets, le nombre de pages. On y accède par le menu \texttt{Fichier} $\rightarrow$ \texttt{Imprimer} : la fenêtre qui s'ouvre affiche, à droite, l'aperçu de chaque page. C'est le moment de vérifier que rien ne dépasse, que rien n'est coupé, et que le dépliant tient bien sur le nombre de pages prévu.

\textbf{2. Le contrôle de l'orthographe.}\par
Une faute d'orthographe sur une affiche publicitaire ou sur la couverture du journal du lycée est très mal perçue. Les logiciels de PAO disposent d'un \cle{correcteur orthographique} qui souligne en rouge les mots inconnus du dictionnaire. Il faut relire chaque zone de texte et corriger les mots signalés — sans oublier qu'une relecture attentive reste nécessaire, car le correcteur ne détecte pas tout (par exemple « \emph{sa} » au lieu de « \emph{ça} » est un mot correctement orthographié, mais mal employé).

\textbf{3. La vérification des objets.}\par
Il faut enfin s'assurer que \textbf{tous les objets sont bien dans la page} : une image ou une zone de texte placée en dehors des limites de la page ne sera pas imprimée, ou sera coupée. On zoome sur les bords de la page et on vérifie chaque élément.

\begin{attention}[L'erreur qui coûte cher]
L'erreur la plus fréquente est d'\textbf{imprimer directement sans passer par l'aperçu avant impression}. Résultat : on découvre après coup qu'un texte est coupé ou qu'une image dépasse, et on a gaspillé papier et encre. \textbf{Règle d'or : on ne clique jamais sur « Imprimer » sans avoir examiné l'aperçu page par page.}
\end{attention}

\courssub{Le vérificateur de conception (concepteur de contrôle)}

Les logiciels de PAO comme Microsoft Publisher intègrent un assistant automatique appelé \cle{vérificateur de conception} (aussi nommé \emph{concepteur de contrôle}). Son rôle : examiner toute la publication et signaler les erreurs courantes avant la diffusion.

\begin{definition}[Vérificateur de conception]
Le \cle{vérificateur de conception} est un outil du logiciel de PAO qui analyse automatiquement la publication et détecte les problèmes de mise en page, par exemple :
\begin{itemize}
\item un \textbf{texte en débordement} : la zone de texte est trop petite, une partie du texte est cachée et ne sera pas imprimée ;
\item un \textbf{objet hors page} : un élément placé en dehors des limites de la feuille ;
\item une image de qualité insuffisante (faible résolution) ou un objet mal aligné.
\end{itemize}
Chaque problème détecté est signalé dans un volet, et il suffit de cliquer dessus pour être conduit à l'objet concerné et le corriger.
\end{definition}

La démarche est simple : on lance le vérificateur (dans Publisher : \texttt{Fichier} $\rightarrow$ \texttt{Informations} $\rightarrow$ \texttt{Vérificateur de conception}), on lit la liste des avertissements, puis on corrige un par un — par exemple en agrandissant une zone de texte trop petite ou en ramenant un objet dans la page.

\courssub{Exportation au format PDF}

\textbf{Le problème.}\par
Tu veux envoyer ton affiche à l'imprimeur de la ville ou la partager avec ton professeur. Si tu envoies le fichier Publisher, le destinataire doit posséder \emph{le même logiciel} pour l'ouvrir — et il pourrait le modifier par accident. Il existe un format universel qui résout ce problème : le PDF.

\begin{definition}[Format PDF]
Le format \cle{PDF} (\emph{Portable Document Format}, « format de document portable ») est un format de fichier qui conserve exactement la mise en page d'un document (textes, images, couleurs, polices), quel que soit l'ordinateur ou le téléphone utilisé pour le lire. Ses avantages :
\begin{itemize}
\item le document est \textbf{lisible partout}, avec un simple lecteur PDF gratuit ;
\item la mise en page est \textbf{identique} sur tous les appareils ;
\item le document est \textbf{difficilement modifiable}, ce qui protège le travail de l'auteur.
\end{itemize}
C'est le format idéal pour diffuser une publication : affiche, dépliant, bulletin d'information, certificat.
\end{definition}

\begin{methode}[Exporter une publication en PDF en 4 étapes]
\begin{enumerate}
\item Ouvrir le menu \texttt{Fichier} et choisir \texttt{Exporter}.
\item Cliquer sur \texttt{Créer un document PDF/XPS}, puis sur le bouton \texttt{Créer PDF/XPS}.
\item Dans la fenêtre d'enregistrement, choisir le dossier de destination, donner un nom clair au fichier (par exemple \texttt{Affiche\_Fete\_Lycee.pdf}) et vérifier que le type de fichier est bien \texttt{PDF}.
\item Cliquer sur \texttt{Publier} (ou \texttt{Enregistrer}), puis ouvrir le fichier PDF obtenu pour vérifier qu'il est conforme à l'aperçu.
\end{enumerate}
\end{methode}

\begin{savaistu}[Le PDF, langue des imprimeurs]
Les imprimeurs professionnels — ceux qui produisent les affiches des concerts à Yaoundé ou les bulletins des établissements scolaires — exigent des fichiers \textbf{PDF haute résolution}. Pourquoi ? Parce que le PDF garantit que les polices, les couleurs et les images seront reproduites exactement comme le concepteur les a préparées, et la haute résolution garantit que l'affiche, même imprimée en très grand format, restera nette. Un simple fichier Word ou une image de faible qualité est refusé !
\end{savaistu}

\courssub{Impression}

Une fois l'aperçu validé (et le PDF exporté pour la diffusion), on peut imprimer. La fenêtre \texttt{Fichier} $\rightarrow$ \texttt{Imprimer} propose plusieurs réglages qu'il faut contrôler attentivement :

\begin{itemize}
\item \textbf{Le choix de l'imprimante} : si plusieurs imprimantes sont disponibles (au cybercafé ou à la salle multimédia du lycée), on sélectionne celle qui fonctionne et qui convient (couleur ou noir et blanc) ;
\item \textbf{Le nombre de copies} : pour imprimer 50 dépliants de la fête du lycée, on saisit 50 dans la zone « Copies » — inutile de relancer 50 fois l'impression ;
\item \textbf{Les pages à imprimer} : on peut imprimer tout le document, la page active, ou une plage de pages (par exemple les pages 2 à 4) ;
\item \textbf{L'impression recto-verso} : pour un \cle{dépliant}, l'impression se fait sur les \textbf{deux faces} de la feuille. On choisit l'option « recto-verso » (ou « impression des deux côtés ») ; si l'imprimante ne le fait pas automatiquement, elle imprime d'abord les pages impaires, puis demande de retourner les feuilles pour imprimer les pages paires.
\end{itemize}

\begin{propriete}[Règle de finalisation d'une publication]
Toute publication destinée à être diffusée suit la même démarche de finalisation, dans cet ordre : \textbf{Vérifier} (orthographe, objets, vérificateur de conception) $\rightarrow$ \textbf{Aperçu} avant impression $\rightarrow$ \textbf{Corriger} les problèmes détectés $\rightarrow$ \textbf{Exporter} en PDF pour la diffusion $\rightarrow$ \textbf{Imprimer} avec les bons réglages. On ne saute aucune étape.
\end{propriete}

\begin{popfigure}
\centering
\begin{tikzpicture}[
  etape/.style={rectangle, rounded corners=4pt, draw=popblue, fill=popblueL, thick, minimum width=3cm, minimum height=0.9cm, align=center, font=\small\bfseries},
  fleche/.style={-{Stealth[length=2.5mm]}, thick, popdark},
  note/.style={font=\scriptsize\itshape, text=popmuted, align=center}
]
\node[etape, fill=popgreenL, draw=popgreen] (verif) {1. Vérifier};
\node[etape, right=1.1cm of verif] (apercu) {2. Aperçu};
\node[etape, right=1.1cm of apercu, fill=poporangeL, draw=poporange] (corriger) {3. Corriger};
\node[etape, right=1.1cm of corriger, fill=poppurpleL, draw=poppurple] (pdf) {4. Exporter PDF};
\node[etape, right=1.1cm of pdf, fill=popgoldL, draw=popgold] (impr) {5. Imprimer};
\draw[fleche] (verif) -- (apercu);
\draw[fleche] (apercu) -- (corriger);
\draw[fleche] (corriger) -- (pdf);
\draw[fleche] (pdf) -- (impr);
\draw[fleche, poporange, dashed] (corriger.north) to[bend left=45] node[above, font=\scriptsize\itshape, text=poporange]{si un problème est détecté} (verif.north);
\node[note, below=0.15cm of verif, align=center]{orthographe, objets,\\vérificateur};
\node[note, below=0.15cm of apercu, align=center]{page par page\\avant d'imprimer};
\node[note, below=0.15cm of corriger, align=center]{texte coupé,\\objet hors page};
\node[note, below=0.15cm of pdf, align=center]{pour diffuser\\sans modification};
\node[note, below=0.15cm of impr, align=center]{imprimante, copies,\\recto-verso};
\end{tikzpicture}
\legende{La démarche de finalisation d'une publication : cinq étapes, avec retour en arrière si l'aperçu révèle un problème.}
\end{popfigure}

\courssub{Complément utile : la résolution d'une image}

\textbf{Observation.}\par
As-tu déjà imprimé une photo prise sur Internet et découvert qu'elle était toute \emph{floue} ou \emph{pixellisée} sur le papier, alors qu'elle semblait nette à l'écran ? Ce phénomène s'explique par la notion de \cle{résolution}.

\begin{definition}[Résolution d'une image]
La \cle{résolution} d'une image numérique est le nombre de \cle{pixels} (petits points de couleur) qui composent l'image, souvent exprimé en pixels par pouce (ppp, ou \emph{dpi} en anglais). Plus la résolution est \textbf{élevée}, plus l'image contient de détails et plus elle restera \textbf{nette à l'impression}. Une image de \textbf{faible résolution} (par exemple une petite image récupérée sur un site web, souvent 72 ppp) apparaît \textbf{floue} ou avec des carrés visibles une fois imprimée en grand (l'impression exige au minimum 150 à 300 ppp).
\end{definition}

\begin{popfigure}
\centering
\begin{tikzpicture}[scale=1]
\begin{scope}
\foreach \i/\j/\c in {0/0/popblue,1/0/popblue,2/0/popblueL,3/0/popblueL,0/1/popblue,1/1/popblueL,2/1/popblueL,3/1/poporange,0/2/popblueL,1/2/popblueL,2/2/poporange,3/2/poporange,0/3/popblueL,1/3/popblueL,2/3/poporange,3/3/popgold}
  \fill[\c] (\i*0.55,\j*0.55) rectangle ++(0.55,0.55);
\draw[popdark, thick] (0,0) rectangle (2.2,2.2);
\draw[popline, thin] (0.55,0)--(0.55,2.2) (1.1,0)--(1.1,2.2) (1.65,0)--(1.65,2.2) (0,0.55)--(2.2,0.55) (0,1.1)--(2.2,1.1) (0,1.65)--(2.2,1.65);
\node[font=\small\bfseries, text=popdark, below=0.2cm] at (1.1,0) {Basse résolution};
\node[font=\scriptsize\itshape, text=popmuted, below=0.55cm] at (1.1,0) {pixels visibles, image floue};
\end{scope}
\begin{scope}[xshift=6.5cm]
\shade[left color=popblue, right color=poporange] (0,0) rectangle (2.2,2.2);
\fill[popgold, opacity=0.85] (1.6,1.5) circle (0.35);
\draw[popdark, thick] (0,0) rectangle (2.2,2.2);
\node[font=\small\bfseries, text=popdark, below=0.2cm] at (1.1,0) {Haute résolution};
\node[font=\scriptsize\itshape, text=popmuted, below=0.55cm] at (1.1,0) {image nette à l'impression};
\end{scope}
\end{tikzpicture}
\legende{À gauche, une image de basse résolution : les pixels (carrés) sont visibles et l'impression sera floue. À droite, une image de haute résolution : les détails sont fins et l'impression nette.}
\end{popfigure}

\begin{aretenir}[Résolution et impression]
Pour une publication destinée à l'impression, on choisit des images de \textbf{haute résolution} (300 ppp recommandés). Une image qui paraît acceptable à l'écran (72 ppp) peut être de qualité insuffisante pour le papier, car l'impression exige beaucoup plus de détails que l'affichage. C'est pourquoi le vérificateur de conception signale parfois les images de faible résolution.
\end{aretenir}

\courssub{Mise en pratique}

\begin{exoresolu}[Finaliser le dépliant de la fête du lycée]
Énoncé. Le club informatique du lycée a créé avec Publisher un dépliant recto-verso annonçant la fête de fin d'année. Avant de le faire imprimer en 100 exemplaires au cybercafé, le président du club te demande de lui indiquer la démarche complète à suivre, dans l'ordre, en justifiant chaque étape.
\tcblower
Solution détaillée. On applique la démarche de finalisation en cinq étapes :
\begin{enumerate}
\item \textbf{Vérifier} : on lance le correcteur orthographique sur toutes les zones de texte, puis le vérificateur de conception (\texttt{Fichier} $\rightarrow$ \texttt{Informations}) pour détecter un éventuel texte en débordement ou un objet hors page. Justification : une faute ou un texte coupé sur 100 exemplaires serait catastrophique.
\item \textbf{Aperçu} : on ouvre \texttt{Fichier} $\rightarrow$ \texttt{Imprimer} et on examine l'aperçu des deux pages (recto et verso). Justification : l'aperçu montre exactement ce qui sortira de l'imprimante.
\item \textbf{Corriger} : si l'aperçu révèle un problème (marge trop grande, image coupée), on retourne dans la publication, on corrige, puis on refait l'aperçu.
\item \textbf{Exporter en PDF} : \texttt{Fichier} $\rightarrow$ \texttt{Exporter} $\rightarrow$ \texttt{Créer un document PDF/XPS}. Justification : le PDF s'ouvrira sans problème sur l'ordinateur du cybercafé, avec une mise en page strictement identique, et personne ne pourra le modifier par erreur.
\item \textbf{Imprimer} : on choisit l'imprimante couleur, on règle le nombre de copies à 100, on sélectionne toutes les pages et on active l'option \textbf{recto-verso} pour que le dépliant soit imprimé sur les deux faces. On imprime d'abord \textbf{un seul exemplaire d'essai} pour vérifier le résultat avant de lancer les 100 copies.
\end{enumerate}
Conclusion : en suivant cet ordre, on évite le gaspillage de papier et d'encre et l'on garantit un dépliant de qualité professionnelle.
\end{exoresolu}

\begin{experience}[TP guidé : exporter et imprimer une affiche]
\textbf{Objectif} : finaliser l'affiche créée dans les séances précédentes et la préparer pour la diffusion.\par
\textbf{Matériel et logiciels} : un ordinateur, un logiciel de PAO (Microsoft Publisher ou équivalent), un lecteur PDF, une imprimante (si disponible).\par
\textbf{Manipulations} :
\begin{enumerate}
\item Ouvre ta publication et lance le correcteur orthographique ; corrige toutes les fautes signalées.
\item Lance le vérificateur de conception et note les problèmes détectés ; corrige-les (agrandis les zones de texte en débordement, ramène les objets dans la page).
\item Ouvre \texttt{Fichier} $\rightarrow$ \texttt{Imprimer} et observe l'aperçu page par page sans imprimer.
\item Exporte la publication en PDF (\texttt{Fichier} $\rightarrow$ \texttt{Exporter} $\rightarrow$ \texttt{Créer un document PDF/XPS}) sous le nom \texttt{Mon\_Affiche.pdf}.
\item Ouvre le fichier PDF obtenu et compare-le à l'aperçu : la mise en page est-elle identique ?
\item Si une imprimante est disponible, imprime \textbf{un seul} exemplaire en réglant l'imprimante, le nombre de copies (1) et les pages à imprimer.
\end{enumerate}
\textbf{Observations et interprétation} : le fichier PDF reproduit exactement la publication, y compris sur un ordinateur ne possédant pas le logiciel de PAO. L'aperçu avant impression a permis de détecter d'éventuels défauts sans consommer ni papier ni encre. Tu as ainsi appliqué la démarche complète d'un professionnel de la publication.
\end{experience}

\begin{exercice}[À toi de jouer]
Ton association veut imprimer 30 affiches annonçant une campagne de propreté au quartier. La publication contient une photo téléchargée sur Internet qui apparaît petite à l'écran.
\begin{enumerate}
\item Cite, dans l'ordre, les cinq étapes de la démarche de finalisation.
\item Pourquoi est-il risqué d'utiliser cette photo en grand format sur l'affiche ? Quel outil du logiciel peut te signaler ce problème ?
\item Pourquoi faut-il exporter l'affiche en PDF avant de l'apporter chez l'imprimeur ?
\end{enumerate}
\end{exercice}

\begin{corrige}[Exercice 1]
\begin{enumerate}
\item Vérifier (orthographe, objets, vérificateur de conception) $\rightarrow$ Aperçu avant impression $\rightarrow$ Corriger $\rightarrow$ Exporter en PDF $\rightarrow$ Imprimer.
\item Une image de faible résolution, agrandie, apparaîtra floue ou pixellisée à l'impression. Le vérificateur de conception peut signaler cette image de qualité insuffisante ; il faut alors la remplacer par une image de haute résolution.
\item Le PDF conserve exactement la mise en page, s'ouvre sur n'importe quel ordinateur avec un lecteur gratuit, et ne peut pas être modifié par accident : l'imprimeur reproduira l'affiche exactement telle qu'elle a été conçue.
\end{enumerate}
\end{corrige}

\begin{aretenir}[L'essentiel de la section]
\begin{itemize}
\item Avant d'imprimer : corriger l'orthographe, vérifier que tous les objets sont dans la page, lancer le \cle{vérificateur de conception} (texte en débordement, objets hors page, images de faible résolution).
\item Toujours passer par l'\cle{aperçu avant impression} (\texttt{Fichier} $\rightarrow$ \texttt{Imprimer}) : jamais d'impression à l'aveugle.
\item Exporter en \cle{PDF} (\texttt{Fichier} $\rightarrow$ \texttt{Exporter} $\rightarrow$ \texttt{Créer un document PDF/XPS}) pour diffuser un document lisible partout et non modifiable.
\item À l'impression : choisir l'imprimante, le nombre de copies, les pages, et le \cle{recto-verso} pour un dépliant.
\item Une image de \textbf{faible résolution} (72 ppp) est floue à l'impression : pour le papier, on utilise des images de haute résolution (300 ppp recommandés).
\end{itemize}
\end{aretenir}

\coursec{Démarche de production et justification}

Dans la section précédente, nous avons appris à finaliser une publication : vérification, exportation en PDF et impression. Tu sais donc désormais manipuler un logiciel de PAO de bout en bout. Mais un vrai professionnel ne se contente pas de « cliquer au bon endroit » : il sait \emph{pourquoi} il fait chaque choix, et il sait l'expliquer. C'est exactement ce que l'on te demandera au BEPC : décrire une démarche de production et \cle{justifier} tes choix. Cette dernière section te donne la méthode complète.

\courssub{Pourquoi une démarche de production ?}

Imagine que le proviseur de ton collège te demande de concevoir une affiche contre le paludisme pour la semaine de

\newpage

\coursec{Activité d'intégration}

\coursec{Leçon N08 : Utilisation d'un logiciel de publication}

\courssub{Objectifs de l'activité}
Cette activité d'intégration vise à valider les compétences suivantes issues du programme officiel de 3ème :
\begin{itemize}
    \item \textbf{Analyser un besoin de communication réel} : identifier le public cible, le message clé, les contraintes techniques (format, support) et les informations obligatoires.
    \item \textbf{Maîtriser l'interface d'un logiciel de PAO} (ici Microsoft Publisher, référence en milieu scolaire camerounais) : espace de travail, ruban, volet \texttt{Pages}, outils \texttt{Zone de texte}, \texttt{Image}, \texttt{Forme}, repères et gabarits.
    \item \textbf{Produire deux objets de PAO distincts et complémentaires} :
    \begin{enumerate}
        \item Une \cle{affiche} A4 portrait (communication événementielle, impact visuel fort, lecture rapide).
        \item Un \cle{dépliant} A4 paysage 3 volets (information détaillée, structuration hiérarchique, lecture séquentielle).
    \end{enumerate}
    \item \textbf{Appliquer les règles fondamentales de la mise en page} : règle des 3 polices maximum, alignements, contraste, équilibre blanc/plein, lisibilité.
    \item \textbf{Valider techniquement la production} : aperçu avant impression, vérificateur de conception, exportation en \cle{PDF} (format universel d'échange et d'impression).
    \item \textbf{Justifier par écrit les choix techniques et esthétiques} : rédiger une argumentation structurée reliant les décisions prises (format, couleurs, disposition, polices) aux objectifs de communication et au public visé.
\end{itemize}

\courssub{Situation}
Le \textbf{Collège Bilingue « La Réussite »} de \textbf{Bafoussam} organise sa \emph{Semaine Culturelle Annuelle} du \textbf{lundi 12 au vendredi 16 mai 2025}. Le thème retenu cette année est : \og \textbf{Culture, Sciences et Numérique : le trio gagnant pour l'excellence} \fg{}.

Le \textbf{Club Informatique}, encadré par M. \textbf{Kamga}, professeur d'informatique, a reçu la commande officielle de la direction pour assurer la \textbf{communication visuelle} de l'événement. Le club compte 18 élèves répartis en 6 groupes de 3. Chaque groupe doit livrer un \textbf{kit de communication complet} composé de deux supports :

\begin{enumerate}
    \item \textbf{Une affiche grand format (A4 portrait)} destinée à être placardée sur les panneaux d'affichage du collège, aux abords de la mairie de Bafoussam, chez les partenaires (CAMTEL Bafoussam, Librairie du Centre) et dans les cybercafés du quartier \emph{Commercial}. Elle doit capter l'attention en moins de 3 secondes.
    \item \textbf{Un dépliant 3 volets (A4 paysage plié en 3)} destiné à être distribué aux élèves, aux parents lors de la réunion de l'APE (Association des Parents d'Élèves) du 5 mai, et déposé à l'accueil du collège. Il détaille le programme jour par jour.
\end{enumerate}

\textbf{Données imposées (cahier des charges fourni par la direction) :}
\begin{itemize}
    \item \textbf{Titre principal} : \og SEMAINE CULTURELLE 2025 \fg{}.
    \item \textbf{Sous-titre / Thème} : \og Culture, Sciences et Numérique \fg{}.
    \item \textbf{Période} : Du lundi 12 au vendredi 16 mai 2025.
    \item \textbf{Lieux principaux} : Cour d'honneur (cérémonies), Salle polyvalente (conférences, théâtre), Terrain de sport (finale de football), Laboratoire informatique (concours de codage / robotique).
    \item \textbf{Programme synthétique} :
    \begin{itemize}
        \item Lundi 12 : Cérémonie d'ouverture (8h), Conférence \og IA et Patrimoine \fg{} (10h), Soirée talents (15h).
        \item Mardi 13 : Concours de dictée/poésie (9h), Atelier Robotique (14h), Match de gala Filles/Garçons (16h).
        \item Mercredi 14 : Journée porte ouverte / Exposition travaux manuels (8h-14h), Théâtre forum \og Harcèlement \fg{} (15h).
        \item Jeudi 15 : Concours de codage \cle{Scratch}/\cle{Python} (9h-12h), Finale Football (15h), Soirée culturelle danses traditionnelles (19h).
        \item Vendredi 16 : Cérémonie de clôture \& Remise des prix (9h), Repas communautaire (13h).
    \end{itemize}
    \item \textbf{Partenaires à citer} : Mairie de Bafoussam, CAMTEL (Fibre Optique), ONG \og Jeunesse Connectée \fg{}, Librairie du Centre.
    \item \textbf{Contacts} : Tel : 6XX XX XX XX / Email : culture@college-lareussite.cm.
\end{itemize}

\textbf{Contraintes techniques et graphiques} :
\begin{itemize}
    \item Logiciel imposé : \textbf{Microsoft Publisher 2016/2019/365} (installé sur les postes de la salle informatique).
    \item Règle des \textbf{3 polices maximum} par document (titre, corps, accent/eventuel).
    \item Charte couleurs suggérée : Vert (\#007A33), Jaune (\#FCD116), Rouge (\#CE1126) — couleurs du drapeau national — + Blanc et Gris foncé (\#333333) pour le texte.
    \item Images : Utiliser la banque d'images libres de droits du logiciel ou les photos de l'édition 2024 fournies sur le disque partagé \texttt{Z:\textbackslash SemaineCulturelle2025\textbackslash Photos}.
    \item Livrables numériques : 1 fichier \texttt{.pub} (source) + 1 fichier \texttt{.pdf} (diffusion) par support, nommés \texttt{Groupe\_N\_Affiche} et \texttt{Groupe\_N\_Depliant}.
\end{itemize}

\courssub{Consignes}
Travaillant en groupe de 3, vous réaliserez les étapes suivantes. Chaque élève doit pouvoir expliquer la démarche du groupe.

\begin{enumerate}
    \item \textbf{Analyse du besoin (15 min)} : Sur une feuille de brouillon, répondez collectivement :
    \begin{itemize}
        \item Qui est le public cible principal ? Secondaire ?
        \item Quel est le message unique à retenir pour l'affiche ? Pour le dépliant ?
        \item Quelles informations sont \emph{obligatoires} (juridiques, pratiques) vs \emph{secondaires} ?
        \item Quelle structure de navigation pour le dépliant 3 volets (6 pages) ?
    \end{itemize}
    \item \textbf{Production de l'affiche A4 Portrait (35 min)} :
    \begin{enumerate}
        \item Créer une publication \texttt{A4 (21 $\times$ 29,7 cm)} orientation \texttt{Portrait}.
        \item Définir des \cle{marges} de sécurité (1,5 cm minimum) et des \cle{repères} pour structurer la page (tiers, centre).
        \item Insérer une \cle{zone de texte} pour le titre principal (police d'affichage, grande taille, couleur forte).
        \item Insérer une \cle{image} illustrative centrale (occupant ~40-50\% de la hauteur).
        \item Insérer les blocs d'informations pratiques (Date, Lieu, Thème) en bas de page, alignés.
        \item Appliquer la règle des 3 polices. Vérifier l'alignement (gauche, centre ou justifié selon l'équilibre).
        \item Ajouter les logos partenaires (petits, en pied de page) et le QR code (généré via un site gratuit) vers le site du collège.
    \end{enumerate}
    \item \textbf{Production du dépliant A4 Paysage 3 volets (45 min)} :
    \begin{enumerate}
        \item Créer une publication \texttt{A4 (29,7 $\times$ 21 cm)} orientation \texttt{Paysage} \textbf{avec gabarit « 3 volets »} (Publisher propose ce modèle au démarrage ou via \texttt{Fichier > Nouveau > Modèles intégrés > Dépliants}).
        \item Comprendre la correspondance pages/volets : \textbf{Page 1 = Volet 1 (Couverture recto)}, \textbf{Pages 2-3-4 = Intérieur (Volets 2, 3, 4)}, \textbf{Page 5 = Volet 5 (Dos)}, \textbf{Page 6 = Volet 6 (Couverture verso / Rabat)}. \emph{Note : en impression recto-verso, l'ordre physique diffère de l'ordre logique à l'écran.}
        \item \textbf{Volet 1 (Couverture)} : Reprendre l'identité visuelle de l'affiche (cohérence), titre, dates, image accrocheuse, mention \og Programme détaillé à l'intérieur \fg{}.
        \item \textbf{Volets 2 à 5 (Intérieur)} : Structurer le programme par jour. Utiliser des \cle{colonnes} (2 colonnes par volet) ou un tableau inséré pour aligner Heure / Événement / Lieu. Une police lisible (corps 10-11 pt), interligne 1,15.
        \item \textbf{Volet 6 (Dos/Rabat)} : Partenaires (logos alignés), Contacts, Mentions légales (ex: \og Publication Club Informatique - Collège La Réussite \fg{}), QR Code inscription ateliers.
        \item Utiliser les \cle{gabarits} (menu \texttt{Disposition > Gabarits}) pour poser les éléments fixes (fond coloré, numérotation de page, logo filigrane) sur toutes les pages intérieures d'un coup.
    \end{enumerate}
    \item \textbf{Vérification technique et Exportation (15 min)} :
    \begin{enumerate}
        \item Basculer en mode \texttt{Aperçu avant impression} (\texttt{Fichier > Imprimer} ou \texttt{Ctrl+P}). Vérifier : aucun texte coupé au bord (marges), images non pixelisées (résolution > 150 dpi), orthographe (correcteur \texttt{F7}), liens hypertextes (QR codes).
        \item Lancer le \texttt{Vérificateur de conception} (\texttt{Fichier > Informations > Vérificateur de conception}) : corriger les erreurs « Texte dans la zone de rognage », « Image en basse résolution ».
        \item Exporter en \textbf{PDF} : \texttt{Fichier > Exporter > Créer un document PDF/XPS}. Options : \texttt{Qualité standard (publishing en ligne et impression)}. Cocher \texttt{Optimiser pour : Impression} si destiné à l'imprimeur professionnel.
        \item Enregistrer les fichiers \texttt{.pub} sources sur la clé USB du groupe / disque réseau.
    \end{enumerate}
    \item \textbf{Justification écrite (15 min)} : Rédiger collectivement un texte d'une demi-page (environ 150 mots) structuré ainsi :
    \begin{itemize}
        \item \textbf{Choix du format} : Pourquoi A4 Portrait pour l'affiche ? Pourquoi A4 Paysage 3 volets pour le dépliant ?
        \item \textbf{Choix des polices} : Citer les 3 polices utilisées et justifier (ex: \og Police à empattement pour le titre = autorité/tradition ; police sans empattement pour le corps = lisibilité écran/impression \fg{}).
        \item \textbf{Choix des couleurs} : Lien avec le thème / identité nationale / contraste accessibilité.
        \item \textbf{Disposition / Hiérarchie} : Comment l'œil est guidé (point focal, flux de lecture Z ou F).
    \end{itemize}
\end{enumerate}

\courssub{Ressources à mobiliser}

\begin{propriete}[Règle des 3 polices maximum]
Dans une publication professionnelle, on limite le nombre de familles de polices distinctes à \textbf{3} :
\begin{itemize}
    \item \textbf{Police 1 (Titres / Display)} : Forte personnalité, grande taille, graisse \texttt{Black} ou \texttt{Bold}. Ex: \texttt{Montserrat ExtraBold}, \texttt{Impact}, \texttt{Playfair Display}, \texttt{Arial Black}.
    \item \textbf{Police 2 (Corps de texte / Body)} : Haute lisibilité en petit corps (9-12 pt), sans empattement (\texttt{Sans-serif}) pour l'écran ou avec empattement (\texttt{Serif}) pour l'impression longue. Ex: \texttt{Open Sans}, \texttt{Roboto}, \texttt{Calibri}, \texttt{Garamond}, \texttt{Arial}.
    \item \textbf{Police 3 (Accent / Légende / Citation)} : Optionnelle, pour différencier un type d'info (citations, légendes photos, boutons d'action). Ex: \texttt{Italique de la Police 2}, ou une police manuscrite \texttt{Script} utilisée avec parcimonie.
\end{itemize}
\emph{Variantes autorisées :} Gras, Italique, Petites capitales, Tailles différentes \textbf{ne comptent pas} comme de nouvelles polices.
\emph{Conseil examen :} Privilégier les polices système standard (Arial, Calibri, Times New Roman, Georgia) pour éviter les problèmes d'incorporation à l'export PDF.
\end{propriete}

\begin{propriete}[Structure logique d'un dépliant 3 volets (A4 Paysage) - 6 pages]
\begin{center}
\begin{tabularx}{\linewidth}{|c|X|c|}\hline
\textbf{Page Publisher} & \textbf{Contenu éditorial} & \textbf{Position physique (pliée)} \\\hline
1 & Couverture (Titre, Image, Date) & Recto - Volet 1 (Visible en premier) \\\hline
2 & Intérieur Gauche (Jour 1-2) & Intérieur - Volet 2 \\\hline
3 & Intérieur Milieu (Jour 3) & Intérieur - Volet 3 (Central) \\\hline
4 & Intérieur Droit (Jour 4-5) & Intérieur - Volet 4 \\\hline
5 & Dos (Partenaires, Mentions) & Verso - Volet 5 (Dos de la couv.) \\\hline
6 & Rabat / Inscription (QR Code) & Verso - Volet 6 (Rabat rentrant) \\\hline
\end{tabularx}
\end{center}
\emph{Astuce :} En mode \texttt{Affichage > Volet Pages}, double-cliquer sur une miniature pour naviguer. Utiliser \texttt{Disposition > Gabarits > Modifier le gabarit} pour les éléments récurrents.
\end{propriete}

\begin{methode}[Exportation PDF prête pour impression professionnelle]
\begin{enumerate}
    \item \texttt{Fichier > Exporter > Créer un document PDF/XPS}.
    \item Cliquer \texttt{Options} pour affiner les réglages.
    \item Dans \texttt{Options} : Sélectionner \textbf{Qualité d'impression (300 ppp)} pour l'imprimeur ou \textbf{ISO 19005-1 (PDF/A)} pour l'archivage.
    \item Cocher \texttt{Inclure les repères de rognage} si l'imprimeur le demande (traits de coupe).
    \item Cocher \texttt{Incorporer les polices} (ou vérifier \texttt{Fichier > Informations > Gérer les polices incorporées}).
    \item Vérifier le PDF ouvert dans Adobe Reader / Edge : zoomer à 100\%, vérifier la netteté des images et la vectorisation du texte.
\end{enumerate}
\end{methode}

\begin{attention}[Pièges fréquents en PAO (Publisher)]
\begin{itemize}
    \item \textbf{Texte masqué / débordement} : Un petit carré rouge avec une croix en bas à droite d'une zone de texte = texte invisible. \emph{Solution :} Agrandir la zone, réduire la police, ou lier à une zone suivante (\texttt{Format > Liaison}).
    \item \textbf{Images basse résolution} : Une image trouvée sur Google (souvent 72 dpi) sera floue en A4. \emph{Solution :} Utiliser \texttt{Insérer > Images > Images en ligne (Creative Commons)} ou photos locales haute définition (> 1500 px de large).
    \item \textbf{Polices non embarquées} : Si vous utilisez une police exotique installée seulement sur votre poste, le PDF l'affichera mal chez l'imprimeur. \emph{Solution :} Privilégier les polices système standard (Arial, Calibri, Times New Roman, Georgia) ou \texttt{Fichier > Informations > Gérer les polices incorporées}.
    \item \textbf{Ordre des pages dépliant} : Ne pas confondre l'ordre à l'écran (1,2,3,4,5,6) et l'ordre d'impression recto-verso. Publisher gère l'imposition si on choisit \texttt{Imprimer > Paramètres > Imprimer sur les deux faces - Retourner les feuilles sur le bord court}.
    \item \textbf{Marge de pli (gouttière)} : Le texte trop près du pli central sera illisible une fois plié. \emph{Solution :} Augmenter la marge intérieure via \texttt{Disposition > Marges > Personnalisées} ou décaler les objets manuellement.
\end{itemize}
\end{attention}

\courssub{Démarche et corrigé commenté}

\begin{exoresolu}[Production du Groupe Modèle « Les Codeurs » — Corrigé commenté étape par étape]
\textbf{Étape 1 — Analyse du besoin (Brouillon)}

\begin{itemize}
    \item \textbf{Public Cible Principal} : Élèves du collège (11-18 ans) $\rightarrow$ Visuel dynamique, couleurs vives, langage accessible, QR Code.
    \item \textbf{Public Secondaire} : Parents, Administration, Partenaires $\rightarrow$ Infos pratiques précises (horaires, lieux, contacts), ton institutionnel sur le dépliant.
    \item \textbf{Message Affiche} : \og Grande Fête Culturelle : 12-16 Mai ! Venez nombreux ! \fg{} (Accroche + Date + Lieu).
    \item \textbf{Message Dépliant} : \og Programme complet jour par jour, inscriptions ateliers, remerciements partenaires. \fg{} (Exhaustivité).
    \item \textbf{Structure Dépliant (6 pages logiques)} :
    \begin{align*}
        \text{P1 (Couv.)} &: \text{Titre + Image + Dates + Thème} \\
        \text{P2 (Volet 2)} &: \text{Lundi 12 \& Mardi 13 (2 colonnes)} \\
        \text{P3 (Volet 3)} &: \text{Mercredi 14 (Focus Expo/Théâtre)} \\
        \text{P4 (Volet 4)} &: \text{Jeudi 15 (Coding + Foot + Soirée)} \\
        \text{P5 (Volet 5 Dos)} &: \text{Vendredi 16 (Clôture) + Partenaires} \\
        \text{P6 (Volet 6 Rabat)} &: \text{Contacts + QR Code Inscription + Mentions}
    \end{align*}
\end{itemize}

\textbf{Étape 2 — Création de l'Affiche A4 Portrait}

\begin{lstlisting}[language=text, title={Journal d'actions Publisher - Affiche}]
1. Fichier > Nouveau > Plus de modèles vierges > A4 (Portrait) > Créer.
2. Affichage > Cocher "Repères" et "Règles". Définir marges 1,5 cm (Disposition > Marges > Personnalisées).
3. Insérer > Zone de texte > Dessiner en haut (hauteur 4 cm, largeur 18 cm centré).
   - Texte : "SEMAINE CULTURELLE 2025"
   - Police : Montserrat ExtraBold (Police 1 - Titre), Taille 48 pt, Couleur Vert #007A33.
   - Alignement : Centré. Interligne : 1,0. Espacement après : 6 pt.
4. Insérer > Zone de texte (sous titre).
   - Texte : "Culture, Sciences et Numérique : le trio gagnant"
   - Police : Playfair Display Italic (Police 3 - Accent), Taille 24 pt, Couleur Rouge #CE1126.
   - Alignement : Centré.
5. Insérer > Images > Ce périphérique > Z:\SemaineCulturelle2025\Photos\Danse_Traditionnelle.jpg.
   - Redimensionner : Largeur 17 cm (marges respectées), Hauteur ~10 cm. Centrer horizontalement.
   - Format > Recadrer > Remplir (si ratio diffère).
   - Effets artistiques > "Légèreté" (optionnel) ou Cadre simple "Carré arrondi" 2 pt Gris.
6. Bas de page : 3 zones de texte alignées gauche (marge 1,5 cm) ou tableau 1 ligne 3 colonnes invisibles.
   - Col 1 (Gauche) : " 12 au 16 Mai 2025" | Police Open Sans Regular (Police 2 - Corps), 14 pt, Gras, Gris #333.
   - Col 2 (Centre) : " Collège La Réussite, Bafoussam" | Même police.
   - Col 3 (Droite) : " Cour d'Honneur / Salle Polyvalente / Labo Info" | Même police.
7. Pied de page (marge bas 1 cm) : Logos partenaires (Mairie, CAMTEL, ONG) insérés comme images, hauteur 1 cm, alignés horizontalement, espacés équitablement (Format > Aligner > Distribuer horizontalement).
8. QR Code : Généré sur qr-code-generator.com (URL: college-lareussite.cm/culture2025). Inséré en bas à droite, 2,5 cm côté. Légende "Scannez-moi !" Police 2, 9 pt, Italique.
9. Vérification : Aucune zone de texte en débordement. Contraste texte/fond OK (Fond blanc, texte Vert). Règle 3 polices respectée (Montserrat, Open Sans, Playfair).
\end{lstlisting}

\begin{popfigure}
\begin{tikzpicture}[scale=0.7, every node/.style={font=\small}]
    % Page A4 Portrait Wireframe
    \draw[popdark, thick] (0,0) rectangle (14.8, 21); % A4 ratio approx 1:1.414
    % Marges
    \draw[popmuted, dashed] (1.5,1.5) rectangle (13.3, 19.5);
    % Titre Zone
    \draw[popblue, fill=popblueL] (1.5, 17.5) rectangle (13.3, 19.5);
    \node[popblue, align=center] at (7.4, 18.5) {TITRE PRINCIPAL\\ormalsize (Montserrat ExtraBold, 48pt, Vert)};
    % Sous-titre
    \draw[poppink, fill=poppinkL] (1.5, 16.2) rectangle (13.3, 17.3);
    \node[poppink, align=center] at (7.4, 16.75) {SOUS-TITRE THÈME (Playfair Italic, 24pt, Rouge)};
    % Image Centrale
    \draw[poporange, fill=poporangeL] (2, 6.5) rectangle (12.8, 15.8);
    \node[poporange, align=center] at (7.4, 11.15) {IMAGE ILLUSTRATIVE\\ormalsize (Danse / Sciences / Numérique)\\~40-50\% Hauteur};
    % Infos Pratiques (3 colonnes)
    \draw[popgreen, fill=popgreenL] (1.5, 3.5) rectangle (5.5, 6.0);
    \node[popgreen, align=center] at (3.5, 4.75) {DATE\\ 12-16 Mai 2025};
    \draw[popgreen, fill=popgreenL] (6.0, 3.5) rectangle (10.0, 6.0);
    \node[popgreen, align=center] at (8.0, 4.75) {LIEU\\ Collège La Réussite};
    \draw[popgreen, fill=popgreenL] (10.5, 3.5) rectangle (13.3, 6.0);
    \node[popgreen, align=center] at (11.9, 4.75) {ESPACES\\ Cour / Salle / Labo};
    % Pied : Logos + QR
    \draw[poppurple, fill=poppurpleL] (1.5, 1.5) rectangle (10, 3.0);
    \node[poppurple, align=center] at (5.75, 2.25) {LOGOS PARTENAIRES (Mairie, CAMTEL, ONG, Librairie)};
    \draw[popgold, fill=popgoldL] (10.5, 1.5) rectangle (13.3, 3.5);
    \node[popgold, align=center] at (11.9, 2.5) {QR CODE\\Inscription / Infos};
\end{tikzpicture}
\legende{Maquette structurelle (Wireframe) de l'Affiche A4 Portrait — Zones fonctionnelles et styles typographiques}
\end{popfigure}

\textbf{Justification choix Affiche} :
\begin{itemize}
    \item \textbf{Format Portrait} : Standard affichage mural, lecture verticale naturelle (Haut $\rightarrow$ Bas), compatible panneaux d'affichage verticaux.
    \item \textbf{Polices} : \texttt{Montserrat} (géométrique, moderne, lisible de loin) pour l'impact ; \texttt{Playfair Display} (élégance, culture) pour le thème ; \texttt{Open Sans} (neutre, très lisible en petit) pour les infos pratiques. \textbf{3 familles exactes.}
    \item \textbf{Couleurs} : Vert/Rouge/Jaune (Drapeau) pour l'ancrage national/camerounais. Fond blanc = économie encre, lisibilité max, propreté.
    \item \textbf{Disposition} : Structure en \textbf{Z} (Titre $\rightarrow$ Image $\rightarrow$ Infos) ou \textbf{F} centrée. Point focal = Image centrale (émotion). Infos pratiques groupées en bas (zone de lecture finale).
\end{itemize}

\textbf{Étape 3 — Création du Dépliant A4 Paysage 3 Volets}

\begin{lstlisting}[language=text, title={Journal d'actions Publisher - Dépliant}]
1. Fichier > Nouveau > Taper "Dépliant" dans recherche > Sélectionner "Dépliant 3 volets A4" > Créer.
   -> Publisher crée 6 pages logiques avec gabarits "Volet 1" à "Volet 6" préconfigurés (marges, colonnes).
2. GABARIT (Disposition > Gabarit > Modifier le gabarit "Dépliant 3 volets") :
   - Fond : Rectangle plein page, remplissage Dégradé Linéaire (Gauche Vert #007A33 -> Droite Jaune #FCD116) Transparence 90% (très léger).
   - Pied de page gabarit : Zone de texte fine, centrée, Police 2 (Open Sans) 8 pt, Gris : "Semaine Culturelle 2025 - Collège La Réussite - culture@college-lareussite.cm".
   - Fermer le mode gabarit.
3. PAGE 1 (VOLET 1 - COUVERTURE) :
   - Reprise visuelle Affiche : Titre "SEMAINE CULTURELLE 2025" (Montserrat, 36pt, Vert), Sous-titre (Playfair, 18pt, Rouge).
   - Image : Photo groupe élèves 2024 (recadrée format paysage), largeur 100% volet.
   - Bas : "Programme détaillé à l'intérieur ►" (Open Sans Italic, 12pt, Gras, Rouge).
4. PAGE 2 (VOLET 2 - LUNDI/MARDI) :
   - Insérer > Tableau > 3 colonnes x 6 lignes (Heure, Événement, Lieu). Largeur 100%.
   - En-tête tableau : "LUNDI 12 MAI" (fusion 3 cellules, fond Vert, texte Blanc, Montserrat Bold).
   - Remplissage lignes. Ligne paire : fond Gris 10% (zébrage).
   - Sous tableau : "MARDI 13 MAI" (même style). Remplissage.
   - Police corps : Open Sans 10 pt.
5. PAGE 3 (VOLET 3 - MERCREDI) : Structure identique. Focus "Journée Portes Ouvertes" : encadré coloré (Remplissage Jaune 20%, bordure Orange).
6. PAGE 4 (VOLET 4 - JEUDI) : Structure identique. Focus "Concours Codage" : Icône </> (Insérer > Icônes > Technologie) devant le titre.
7. PAGE 5 (VOLET 5 - VENDREDI / DOS) :
   - Haut : "VENDREDI 16 MAI - CLÔTURE" (Tableau).
   - Bas : "NOS PARTENAIRES". Insérer logos (Mairie, CAMTEL, ONG, Librairie). Format > Aligner > Aligner au centre / Distribuer horizontalement. Hauteur uniforme 1.5 cm.
8. PAGE 6 (VOLET 6 - RABAT) :
   - "INSCRIPTIONS ATELIERS (Robotique, Codage, Théâtre)".
   - QR Code géant (4 cm) centré. Légende : "Places limitées ! Inscrivez-vous avant le 10 Mai."
   - Mentions légales : "Conception : Club Informatique - Encadreur M. Kamga - Impression : Imprimerie Centrale Bafoussam".
9. VÉRIFICATION GLOBALE (Fichier > Informations > Vérificateur de conception) :
   - Corriger "Texte près du bord" sur P2 (marge intérieure pli) : décaler tableau de 0.5 cm vers l'extérieur.
   - Corriger "Image Basse Résolution" sur P1 : Remplacer par version HD du disque Z.
10. EXPORT PDF : Fichier > Exporter > PDF/XPS > Options > "Qualité d'impression (300 ppp)" > Cocher "Repères de rognage" > Publier sous "Groupe3_Depliant.pdf".
\end{lstlisting}

\begin{popfigure}
\begin{tikzpicture}[scale=0.55, every node/.style={font=\tiny}]
    % Dépliant A4 Paysage - 6 volets pliés (Vue à plat)
    % Dimensions A4 Paysage : 29.7 x 21 cm. 3 volets = 9.9 cm large chacun.
    \def\voletW{9.9}
    \def\voletH{21}
    \def\gap{0.3}
    % Fonction pour dessiner un volet : #1=x, #2=y, #3=pageNum, #4=title, #5=color, #6=content
    \newcommand{\drawVolet}[6]{
        \draw[popdark, thick] (#1, #2) rectangle (#1+\voletW, #2+\voletH);
        \node[#5, align=center, anchor=north west] at (#1+0.2, #2+\voletH-0.2) {\textbf{Page #3}\\\small #4};
        \node[popmuted, align=center] at (#1+\voletW/2, #2+\voletH/2) {#6};
    }
    % RECTO (Pages 1, 2, 3) - Haut
    \drawVolet{0}{0}{1}{Couverture}{popblue}{Titre + Image + Dates\\« Programme à l'intérieur »};
    \drawVolet{\voletW+\gap}{0}{2}{Intérieur Gauche}{popgreen}{Lundi 12 \& Mardi 13\\Tableaux Heure/Événement/Lieu};
    \drawVolet{\voletW+\gap+\voletW+\gap}{0}{3}{Intérieur Centre}{poporange}{Mercredi 14\\Expo / Théâtre Forum\\(Focus Visuel)};
    % VERSO (Pages 4, 5, 6) - Bas (imprimé au dos)
    \drawVolet{0}{-\voletH-\gap}{4}{Intérieur Droit}{poppurple}{Jeudi 15\\Codage / Foot / Soirée};
    \drawVolet{\voletW+\gap}{-\voletH-\gap}{5}{Dos (Volet 5)}{poppink}{Vendredi 16 Clôture\\PARTENAIRES (Logos)};
    \drawVolet{\voletW+\gap+\voletW+\gap}{-\voletH-\gap}{6}{Rabat (Volet 6)}{popgold}{INSCRIPTIONS\\QR CODE Géant\\Mentions Légales};
    % Flèches pliage
    \draw[popdark, <->, thick] (\voletW, -0.5) -- (\voletW+\gap, -0.5) node[midway, above, font=\tiny] {Pli 1};
    \draw[popdark, <->, thick] (2*\voletW+\gap, -0.5) -- (2*\voletW+2*\gap, -0.5) node[midway, above, font=\tiny] {Pli 2};
\end{tikzpicture}
\legende{Structure logique du Dépliant 3 Volets (6 pages) dans Publisher — Correspondance Page/Volet et Contenu éditorial}
\end{popfigure}

\textbf{Justification choix Dépliant} :
\begin{itemize}
    \item \textbf{Format Paysage 3 volets} : Standard professionnel pour documentation pliée (format DL une fois plié = 99$\times$210 mm, rentre dans enveloppe standard). Orientation paysage = lecture en « accordéon » confortable, permet 2 colonnes de texte par volet (gain de place, lisibilité).
    \item \textbf{Polices} : Identiques à l'affiche (\textbf{Cohérence identitaire}). \texttt{Montserrat} pour les en-têtes de jour (navigation), \texttt{Open Sans} pour le corps des tableaux (densité d'info), \texttt{Playfair} uniquement pour la citation d'accroche en P3.
    \item \textbf{Couleurs} : Reprise du dégradé Vert/Jaune en fond de gabarit (subtil, 10\% opacité) pour l'unité visuelle. Zébrage tableaux (Gris 10\%) pour guider l'œil sur les lignes. Couleurs fortes (Vert/Rouge) réservées aux en-têtes de jour (repérage rapide).
    \item \textbf{Disposition} : \textbf{Gabarit} utilisé pour éléments récurrents (pied de page, fond) $\rightarrow$ Gain de temps, cohérence garantie. \textbf{Tableaux} pour le programme $\rightarrow$ Alignement parfait Heure/Lieu, modification facile. \textbf{Hiérarchie} : Jour (Gras, Couleur) > Heure (Gras) > Événement (Normal) > Lieu (Italique, Gris).
\end{itemize}

\textbf{Étape 4 — Vérification \& Export (Check-list finale)}

\begin{center}
\begin{tabularx}{\linewidth}{|l|X|c|}\hline
\textbf{Point de contrôle} & \textbf{Action / Outil Publisher} & \textbf{Statut} \\\hline
Orthographe / Grammaire & \texttt{Révision > Orthographe (F7)} + Relecture croisée binôme & \textbf{OK} \\\hline
Débordement texte & Icône \texttt{Redimensionnement automatique} / Liaison zones & \textbf{OK} \\\hline
Résolution images & Vérificateur de conception > « Images en basse résolution » & \textbf{OK} (Remplacées) \\\hline
Marges / Zone de rognage & Aperçu avant impression + Vérificateur > « Texte dans zone rognage » & \textbf{OK} (Décalé 0,5 cm) \\\hline
Polices incorporées & Fichier > Informations > Gérer les polices > « Incorporer toutes les polices » & \textbf{OK} \\\hline
Hyperliens / QR Codes & Test scan smartphone (QR) + Clic Ctrl+Clique (liens mailto:) & \textbf{OK} \\\hline
Cohérence visuelle (Affiche vs Dépliant) & Comparaison côte à côte : Polices, Couleurs, Logos, Ton & \textbf{OK} \\\hline
Nommage fichiers & \texttt{Groupe3\_Affiche.pub/.pdf}, \texttt{Groupe3\_Depliant.pub/.pdf} & \textbf{OK} \\\hline
\end{tabularx}
\end{center}

\textbf{Étape 5 — Texte de justification rédigé par le groupe (Exemple)}

\begin{quote}
{\itshape « Pour l'affiche, nous avons choisi le format A4 Portrait car c'est le standard d'affichage mural dans les couloirs et les panneaux extérieurs du collège ; la lecture verticale (Z-pattern) guide naturellement l'œil du titre accrocheur (Montserrat, vert drapeau, 48 pt) vers l'image centrale émotionnelle (danse traditionnelle), pour finir sur les infos pratiques (Open Sans, 14 pt) en bas. La troisième police, Playfair Display Italic (rouge), signe le thème "Culture" avec élégance sans surcharger. 

Pour le dépliant, le format A4 Paysage 3 volets (pli roulé) offre 6 panneaux logiques idéaux pour séquencer le programme jour par jour. L'orientation paysage permet une mise en page en 2 colonnes par volet (tableaux Heure/Événement/Lieu), optimisant la densité d'information tout conservant une lisibilité (Open Sans 10 pt). Nous avons utilisé les gabarits pour imposer le fond dégradé vert/jaune (identité nationale) et le pied de page contact sur les 4 pages intérieures, garantissant une cohérence graphique parfaite avec l'affiche. L'export PDF "Qualité Impression 300 ppp" avec repères de rognage assure une reproduction fidèle chez l'imprimeur de la Mairie. »}
\end{quote}
\tcblower
\textbf{Commentaires pédagogiques du professeur (M. Kamga)} :
\begin{itemize}
    \item \textbf{Analyse} : Le groupe a bien distingué la fonction \emph{accroche} (affiche) de la fonction \emph{information complète} (dépliant). C'est la clé de la PAO : \textbf{le format suit la fonction}.
    \item \textbf{Technique} : L'utilisation du \textbf{Gabarit} pour le dépliant est une compétence niveau 3ème attendue (gain de temps, cohérence). La correction proactive des alertes du \textbf{Vérificateur de conception} (marge de pli, résolution images) montre une maturité technique « pré-professionnelle ».
    \item \textbf{Typographie} : Respect strict de la règle des 3 polices. Le choix \texttt{Montserrat / Open Sans / Playfair} est pertinent : contraste Serif/Sans-Serif maîtrisé, lisibilité écran/impression assurée. Attention : \texttt{Playfair} en italique corps 24 pt reste lisible, mais en corps 10 pt elle le serait moins (empattements fins).
    \item \textbf{Couleur} : L'usage du drapeau (Vert/Jaune/Rouge) est pertinent pour un événement scolaire officiel au Cameroun. Le fond blanc de l'affiche assure le contraste maximal (accessibilité). Le fond dégradé 10\% du dépliant est une touche pro subtile.
    \item \textbf{Justification} : Le texte écrit est structuré, utilise le vocabulaire technique (Z-pattern, gabarit, zébrage, repères de rognage, incorporation polices) et relie \textbf{choix technique $\rightarrow$ finalité communication}. C'est exactement ce qui est attendu à l'oral du BEPC ou en épreuve pratique.
\end{itemize}
\end{exoresolu}

\courssub{Bilan de l'activité}
Cette activité d'intégration a permis de mettre en évidence les acquis majeurs de l'unité « Utilisation d'un logiciel de publication » :

\begin{aretenir}[Compétences validées]
\begin{itemize}
    \item \textbf{Savoir-faire technique} : Lancer Publisher, choisir le bon modèle (Affiche vs Dépliant 3 volets), naviguer dans l'interface (Ruban, Volet Pages, Gabarits), manipuler les objets (Zones de texte, Images, Tableaux, Formes, QR Codes), gérer les pages/volets.
    \item \textbf{Savoir-faire méthodologique} : Analyser un cahier des charges réel (public, message, contraintes), planifier la structure (wireframe mental ou papier), appliquer les règles de mise en page (marges, repères, alignement, grille, zébrage, hiérarchie visuelle), utiliser le \textbf{Vérificateur de conception} comme outil qualité.
    \item \textbf{Savoir-faire de production} : Respecter la \textbf{règle des 3 polices}, assurer la \textbf{cohérence graphique} entre deux supports distincts (charte couleurs, polices, logo), produire des fichiers \textbf{PDF normés} (qualité impression, polices incorporées, repères de coupe) prêts pour la chaîne graphique professionnelle.
    \item \textbf{Savoir-être / Communication} : Travailler en équipe (répartition : chef de projet/maquettiste, rédacteur/relecteur, iconographe/technique export), argumenter par écrit ses choix techniques avec un vocabulaire précis, présenter son travail face à la classe/client (direction).
\end{itemize}
\end{aretenir}

\begin{savaistu}[Le saviez-vous ? — La PAO au Cameroun]
Au Cameroun, la \textbf{PAO (Publication Assistée par Ordinateur)} est un pilier de l'économie numérique locale. Les \emph{cybercafés} et \emph{bureaux de photocopie} (très présents à Yaoundé, Douala, Bafoussam, Garoua) réalisent 80\% de leur chiffre d'affaires sur la conception/impression de : \textbf{faire-part} (mariages, décès), \textbf{affiches} (événements culturels, politiques, religieux), \textbf{dépliants} (promotions boutiques, programmes scolaires), \textbf{cartes de visite} et \textbf{en-têtes de correspondance} pour PME/ONG. Maîtriser Publisher (ou Scribus, alternative libre) est une \textbf{compétence directement employable} dès la sortie du BEPC pour un job d'été ou un stage en imprimerie. Le format \textbf{PDF/X-1a} est le standard exigé par les grosses imprimeries industrielles (ex: \emph{Imprimerie Nationale}, \emph{SOPECAM}) pour garantir la fidélité des couleurs (CMJN) et des polices.
\end{savaistu}

\begin{exercice}[Entraînement individuel - Évaluation formative]
\begin{enumerate}
    \item Vous devez créer une affiche A5 (14,8 $\times$ 21 cm) pour la vente de gâteaux du club. Quel paramètre modifier \textbf{en premier} dans \texttt{Fichier > Nouveau} ?
    \item Dans un dépliant 3 volets A4 paysage, sur quelle \textbf{page logique Publisher} place-t-on l'adresse de l'école et le logo pour qu'ils se retrouvent au dos de la couverture une fois plié ?
    \item Citez 3 erreurs détectées par le \texttt{Vérificateur de conception} qui bloqueraient une impression professionnelle.
    \item Pourquoi est-il risqué d'utiliser la police \texttt{« DJB Chalk It Up »} (téléchargée sur DaFont) sans l'incorporer dans le PDF ?
    \item Proposez une combinaison de 3 polices (Titre / Corps / Accent) élégante et lisible pour un bulletin de notes officiel.
\end{enumerate}
\end{exercice}

\begin{corrige}[Exercice Entraînement PAO]
\begin{enumerate}
    \item \textbf{Format de page / Taille} : Choisir « A5 » dans les modèles ou \texttt{Disposition > Taille > A5}. C'est la base géométrique de tout le reste.
    \item \textbf{Page 5 (Volet 5 / Dos)}. Dans la structure 3 volets pli roulé : Page 1=Couv. Recto, Pages 2-3-4=Intérieur, Page 5=Dos Couv. (Verso P1), Page 6=Rabat.
    \item Exemples : « Texte dans la zone de rognage » (trop près du bord/du pli), « Image en basse résolution » (< 150 ppp pour l'impression), « Police non incorporée » (risque de substitution), « Objet en dehors de la page », « Lien hypertexte brisé ».
    \item Si le PDF est ouvert sur un poste sans cette police (imprimeur, jury BEPC, parent), le lecteur PDF la substituera par une police par défaut (ex: Arial), cassant la mise en page (débordements, alignements ruinés) et l'identité visuelle. \textbf{Solution} : Incorporer la police (Fichier > Informations > Gérer les polices) ou utiliser une police système standard.
    \item Proposition : \textbf{Titre} : \texttt{Garamond Bold} (Serif, autorité, tradition scolaire) — \textbf{Corps} : \texttt{Calibri Regular} (Sans-serif, lisible petit corps, standard Windows) — \textbf{Accent} : \texttt{Garamond Italic} (Cohérence famille titre, pour mentions « Mention Bien / Très Bien »). \emph{Alternative moderne} : \texttt{Montserrat SemiBold} / \texttt{Source Sans Pro} / \texttt{Montserrat Italic}.
\end{enumerate}
\end{corrige}

\newpage

\coursec{Méthodes et exercices résolus}

\coursec{Leçon 08 : Utilisation d'un logiciel de publication assistée par ordinateur (PAO)}

\noindent
Bienvenue dans cette leçon consacrée à la \cle{Publication Assistée par Ordinateur (PAO)}. Contrairement au traitement de texte (Word, Writer) qui gère un flux de texte continu, la PAO (Scribus, InDesign, Publisher) repose sur une approche \textbf{objet} et \textbf{page} : chaque élément (texte, image, forme) vit dans un \textbf{cadre} (frame) positionné avec une précision millimétrique. C'est la compétence clé pour produire des documents « prêts à imprimer » : affiches, flyers, journaux scolaires, cartes de visite, bulletins. Au BEPC, l'épreuve pratique attend de vous la maîtrise du flux complet : \emph{Gabarit $\rightarrow$ Structure (Cadres/Maîtres) $\rightarrow$ Contenu (Styles/Images) $\rightarrow$ Contrôle qualité (Preflight) $\rightarrow$ Export normé (PDF/X)}.

\begin{savaistu}[De la composition manuelle au « Desktop Publishing »]
    Historiquement, la composition d'une page impliquait du plomb fondu (Linotype) ou du photocomposition coûteuse, réservée aux imprimeries. En 1985, \textbf{Aldus PageMaker} sur Macintosh (avec l'imprimante LaserWriter d'Apple et le langage PostScript d'Adobe) démocratise la mise en page : on voit la page à l'écran « tel quel » (\emph{WYSIWYG}). Au Cameroun, l'arrivée de la PAO dans les années 90 a transformé les cybercafés et les secrétariats en mini-studios graphiques. Aujourd'hui, \textbf{Scribus} (logiciel libre, multiplateforme) est le standard pédagogique au MINESEC car il offre des fonctions professionnelles (gestion CMJN, PDF/X, pages maîtres, styles) sans licence coûteuse, tout en tournant sur les configurations modestes de nos salles informatiques.
\end{savaistu}

\courssub{Concepts fondamentaux et vocabulaire technique}

\begin{definition}[Publication Assistée par Ordinateur (PAO)]
    Ensemble de techniques informatiques permettant de composer des documents destinés à l'impression professionnelle ou à la diffusion numérique haute qualité, en assemblant du texte, des images et des graphiques vectoriels dans des \textbf{cadres} positionnés librement sur une \textbf{page} virtuelle aux dimensions réelles.
\end{definition}

\begin{definition}[Bords perdus (Bleeds / Fonds perdus)]
    Zone d'extension du document (\textbf{généralement 3 à 5 mm}) au-delà du format fini, qui sera rognée par le massicot de l'imprimeur. Tout élément graphique touchant le bord de la page (fond coloré, photo pleine page) \textbf{doit} déborder dans cette zone pour éviter un filet blanc inesthétique après coupe.
\end{definition}

\begin{definition}[Marge de sécurité (Zone tranquille / Safe zone)]
    Espace intérieur (généralement \textbf{5 à 10 mm} depuis le bord fini) où l'on place les éléments vitaux (textes, logos, visages). Rien d'important ne doit se trouver entre le bord fini et cette marge, sous peine d'être coupé ou trop près du bord.
\end{definition}

\begin{propriete}[Différences clés : Traitement de texte vs PAO]
    \begin{center}
    \begin{tabularx}{\linewidth}{|l|X|X|}
    \hline
    \rowcolor{popblueL}
    \textbf{Critère} & \textbf{Traitement de texte (Word, Writer)} & \textbf{PAO (Scribus, InDesign, Publisher)} \\
    \hline
    Modèle & Flux continu (texte qui coule) & Objets positionnés (Cadres/Frames) \\
    \hline
    Précision & Caractère / Paragraphe & Millimétrique (position X, Y, Largeur, Hauteur) \\
    \hline
    Pages & Séquentielles, ajout auto & Maîtrisées, nombre fixe, \texttt{Facing pages} \\
    \hline
    Couleurs & RVB (écran) & \textbf{CMJN (Quadrichromie)} + Pantone + RVB \\
    \hline
    Images & Embarquées (dans le .docx) & \textbf{Liées (lien vers fichier source)} \\
    \hline
    Sortie pro & PDF basique & \textbf{PDF/X-1a, PDF/X-3, PDF/X-4} (normes ISO) \\
    \hline
    Usage type & Rapports, lettres, thèses & Affiches, journaux, packaging, magazines, flyers \\
    \hline
    \end{tabularx}
    \end{center}
\end{propriete}

\courssub{Lancement, interface et configuration d'un nouveau document}

\begin{methode}[Découvrir l'interface de Scribus (Fenêtre principale)]
    \textbf{Objectif :} Identifier les zones de travail pour gagner en efficacité.
    \begin{enumerate}
        \item \textbf{Barre de menus} (\texttt{Fichier, Édition, Affichage, Page, Objet, Texte, Fichier, Fenêtre, Aide}) : Accès à toutes les fonctions.
        \item \textbf{Barre d'outils principale} (sous les menus) : Nouveau, Ouvrir, Enregistrer, PDF, Annuler/Rétablir, Zoom, Mode Aperçu (\texttt{W}).
        \item \textbf{Barre d'outils Outils} (gauche, verticale) : \texttt{Sélection} (flèche noire, \texttt{V}), \texttt{Sélection contenu} (flèche blanche), \texttt{Texte} (\texttt{T}), \texttt{Image} (\texttt{I}), \texttt{Forme} (Rectangle, Ellipse, Polygone), \texttt{Ligne}, \texttt{Bézier}, \texttt{Tableau}, \texttt{Chaîner cadres} (icône chaîne).
        \item \textbf{Propriétés (F2)} : \textbf{Le panneau le plus important}. Onglets : \texttt{X,Y,Z} (Position/Taille/Rotation), \texttt{Texte} (Police, Taille, Alignement, Style, Habillage), \texttt{Image} (Échelle, Résolution, Aperçu), \texttt{Couleurs} (Remplissage, Contour), \texttt{Ligne} (Type, Largeur), \texttt{Forme} (Coins arrondis, Conversion), \texttt{GPI} (Gestionnaire d'images).
        \item \textbf{Panneau Pages / Pages Maîtres} (souvent à droite) : Vignettes des pages, double-clic pour naviguer, icône « Éditer le maître » (page avec coin plié).
        \item \textbf{Guides et Règles} : Clics sur la règle pour créer un guide manuel. \texttt{Affichage > Guides} pour verrouiller/cacher.
        \item \textbf{Barre d'état} (bas) : Infos curseur, zoom, mode grille, préflight (icône avion/feuille verte/rouge).
    \end{enumerate}
    \textbf{Reflexe :} Travaillez en \textbf{Mode Aperçu (\texttt{W})} régulièrement pour masquer les contours de cadres, guides et marges, et voir le rendu final.
\end{methode}

\begin{methode}[Préparer et configurer un nouveau document de publication]
    \textbf{Objectif :} Créer un fichier PAO aux dimensions adaptées au support final (affiche, flyer, carte de visite, bulletin) et définir les repères de mise en page.
    \begin{enumerate}
        \item \textbf{Analyser le besoin :} Identifier le format final (ex. A5 pour un flyer, A4 pour une lettre, dimensions personnalisées pour une carte de visite $85 \times 55$ mm), l'orientation (Portrait/Paysage) et les marges de sécurité (généralement $5$ à $10$ mm bords perdus).
        \item \textbf{Lancer le logiciel} (Scribus) et choisir \texttt{Fichier > Nouveau} (ou \texttt{Ctrl+N}).
        \item \textbf{Configurer la page (Onglet \texttt{Document}) :}
              \begin{itemize}
                  \item \texttt{Taille} : Prédéfinie (A4, A5, Lettre, Légal) ou \texttt{Personnalisée} (Largeur $\times$ Hauteur en mm/cm).
                  \item \texttt{Orientation} : Portrait (hauteur $>$ largeur) ou Paysage.
                  \item \texttt{Unités} : \textbf{Millimètres (mm)} impératif pour l'impression.
                  \item \texttt{Marges} : Définir les marges intérieures (zone de texte sûre) et extérieures. Cocher \emph{Facing pages} (Pages face à face) \textbf{uniquement} pour les documents multipages reliés ou pliés (livre, magazine, flyer plié 4 pages).
              \end{itemize}
        \item \textbf{Configurer les bords perdus (Onglet \texttt{Bords perdus}) :} Mettre \textbf{3 mm} (standard imprimeur) sur les 4 côtés (Haut, Bas, Gauche, Droite). \emph{Note :} Si \texttt{Facing pages} coché, le bord perdu « Intérieur » (côté pli) est souvent mis à 0 mm, mais 3 mm partout simplifie l'export.
        \item \textbf{Activer les aides visuelles :} \texttt{Affichage > Grille}, \texttt{Guides de marge}, \texttt{Guides de colonne}. Configurer le nombre de colonnes (ex. $3$ pour un journal) et la gouttière (espace entre colonnes, ex. $4$ mm) via \texttt{Fichier > Configurer le document > Guides}.
        \item \textbf{Enregistrer immédiatement :} \texttt{Fichier > Enregistrer sous} $\rightarrow$ Nom explicite sans accents ni espaces (ex. \texttt{Flyer\_Kermesse\_College\_2025.sla}) $\rightarrow$ Dossier projet dédié (contenant le \texttt{.sla} et un sous-dossier \texttt{/Images}).
    \end{enumerate}
    \textbf{Repère :} Un document bien configuré dès le départ évite les décalages à l'impression et garantit le respect des \emph{bords perdus}.
\end{methode}

\begin{exoresolu}[Configuration d'un flyer plié pour la kermesse du collège]
    \textbf{Énoncé :} Vous êtes chargé de concevoir un flyer A5 (format fini $148 \times 210$ mm) pour la kermesse de fin d'année du Collège Bilingue de Ngaoundéré. L'imprimeur exige un fond perdu de $3$ mm de chaque côté et une marge de sécurité intérieure de $5$ mm. Le flyer sera plié en deux (format fermé A5, ouvert A4 $\rightarrow$ 4 pages). Configurez le document dans Scribus.
    \tcblower
    \textbf{Solution détaillée :}
    \begin{enumerate}
        \item \textbf{Analyse du format :} Format fini = A5 ($148 \times 210$ mm). Plié en 2 $\rightarrow$ Feuillet ouvert = A4 ($210 \times 297$ mm). Un flyer plié en 2 (1 feuillet) = \textbf{4 pages A5}. Dans Scribus, on travaille en « Pages d'imprimeur » (Printer Spreads) ou « Pages lecteurs » (Reader Spreads). \textbf{Approche pédagogique standard :} On crée un document de \textbf{4 pages} au format \textbf{A5}, \textbf{Facing pages} coché (pour voir les doubles pages), marges $10$ mm (sécurité), bords perdus $3$ mm. L'imposition (réordonnancement pour l'impression) sera gérée par l'imprimeur ou via \texttt{Fichier > Imprimer > Imposition} plus tard.
        \item \textbf{Configuration dans Scribus :}
              \begin{itemize}
                  \item \texttt{Fichier > Nouveau Document}.
                  \item \texttt{Taille} : A5 ($148 \times 210$ mm).
                  \item \texttt{Orientation} : Portrait.
                  \item \texttt{Pages} : \textbf{4}. Cocher \textbf{Pages face à face (Facing pages)}.
                  \item \texttt{Marges} : Gauche/Droite/Haut/Bas = \textbf{10 mm} (marge de sécurité confortable).
                  \item \textbf{Bords perdus :} Onglet \texttt{Bords perdus} $\rightarrow$ Haut, Bas, Gauche, Droite = \textbf{3 mm}.
              \end{itemize}
        \item \textbf{Guides de pli (Optionnel mais pro) :} Sur le Maître (Master Page), ajouter un guide horizontal manuel à $Y = 105$ mm (milieu de l'A5) si pli horizontal, ou guide vertical à $X = 74$ mm si pli vertical (format italien). Ici pli standard (portrait $\rightarrow$ pli horizontal au milieu de l'A5 ? Non, A5 portrait plié en 2 donne A6. L'énoncé dit "Flyer A5 plié en deux, ouvert A4". Donc le document *ouvert* est A4. Le format *fermé* est A5. Cela signifie qu'on conçoit sur un format \textbf{A4} (2 pages face à face = 1 feuillet recto/verso = 4 pages A5).
              \textbf{Correction stratégie :} On configure en \textbf{A4 Portrait}, \textbf{2 pages}, \textbf{Facing pages}. Le pli est au milieu ($148,5$ mm). Marge intérieure (côté pli) = $10$ mm, Extérieure = $10$ mm. Bords perdus $3$ mm. C'est la méthode la plus simple pour un élève de 3ème.
        \item \textbf{Enregistrement :} \texttt{Fichier > Enregistrer sous} $\rightarrow$ \texttt{/Projets\_PAO/Kermesse\_Ngaoundere/Flyer\_Kermesse\_2025.sla}.
    \end{enumerate}
    \textbf{Conclusion :} Le document est prêt. La zone de travail utile (hors bords perdus) est de $204 \times 291$ mm par page A4. Le pli central est repéré à $148,5$ mm. Les éléments graphiques de fond doivent déborder sur la zone de bords perdus ($3$ mm), les textes importants doivent rester dans les marges de $10$ mm.
\end{exoresolu}

\courssub{Structurer la mise en page : Pages Maîtres, Cadres et Chaînage}

\begin{definition}[Page Maître (Master Page / Gabarit)]
    Modèle de page contenant les éléments récurrents (en-têtes, pieds de page, numérotation, logos, fonds de page, guides de colonnes). Toute page à laquelle ce maître est appliqué hérite de ces éléments. Modification du maître $\rightarrow$ mise à jour automatique sur toutes les pages concernées.
\end{definition}

\begin{definition}[Chaînage des cadres texte (Text Flow / Linking)]
    Mécanisme qui relie plusieurs cadres texte entre eux pour former un seul flux de texte continu. Lorsque le premier cadre est plein, le texte « coule » automatiquement dans le suivant. Indispensable pour les articles sur plusieurs colonnes ou pages.
\end{definition}

\begin{methode}[Structurer la mise en page avec les Pages Maîtres et les Cadres]
    \textbf{Objectif :} Organiser l'espace par des conteneurs (cadres texte, image, tableau) et automatiser les éléments répétitifs via les Pages Maîtres.
    \begin{enumerate}
        \item \textbf{Passer en mode "Page Maître" :} Menu \texttt{Fenêtre > Pages Maîtres} (ou icône « Éditer le maître » en bas du panneau Pages). Par défaut : \texttt{Normal Left} (page gauche) et \texttt{Normal Right} (page droite).
        \item \textbf{Créer les cadres récurrents sur le Maître :}
              \begin{itemize}
                  \item \textbf{Cadre Texte (Pied de page) :} Outil \texttt{Insérer un cadre texte} (\texttt{Shift+T}). Dessiner la zone en bas de page. Insérer \texttt{Insertion > Marqueurs > Numéro de page} (raccourci \texttt{Ctrl+Alt+Maj+P} ou clic droit \texttt{Insérer un marqueur}). Appliquer un style « PiedPage ».
                  \item \textbf{Cadre Image (Logo/Fond fixe) :} Outil \texttt{Insérer un cadre image} (\texttt{I}). Positionner le logo. \textbf{Verrouiller} (\texttt{Clic droit > Verrouiller} ou icône cadenas dans Propriétés X,Y,Z) pour éviter déplacement accidentel.
                  \item \textbf{Lignes / Formes décoratives :} Outil \texttt{Insérer une ligne} ou \texttt{Forme} pour filets de séparation.
              \end{itemize}
        \item \textbf{Configurer la grille de colonnes sur le Maître :} \texttt{Fichier > Configurer le document > Guides} : Colonnes = $3$ (ex.), Gouttière = $5$ mm. Les guides violets apparaissent sur le maître et toutes les pages associées.
        \item \textbf{Appliquer le Maître aux pages :} Revenir en mode \texttt{Pages normales} (icône feuille blanche). Panneau \texttt{Pages} : glisser-déposer l'icône du Maître sur les pages cibles, ou clic droit sur page $\rightarrow$ \texttt{Appliquer le maître}.
        \item \textbf{Créer les cadres de structure sur les pages courantes (Contenu) :} Sur la page 1, dessiner les blocs : \emph{Zone Titre}, \emph{Zone Image principale}, \emph{Zone Infos}. Utiliser \texttt{Fenêtre > Aligner et Distribuer} pour centrer parfaitement par rapport à la page ou aux marges.
        \item \textbf{Lier les cadres texte (Chaînage) :} Pour un article sur plusieurs colonnes/pages :
              \begin{enumerate}
                  \item Créer Cadre 1, Cadre 2, Cadre 3 (alignés sur guides colonnes).
                  \item Sélectionner Cadre 1 (outil Flèche noire).
                  \item Cliquer icône \texttt{Chaîner les cadres texte} (chaîne) dans barre d'outils (ou \texttt{Édition > Chaîner les cadres}).
                  \item Cliquer sur Cadre 2 $\rightarrow$ Lien créé (flèche visible). Répéter Cadre 2 vers Cadre 3.
              \end{enumerate}
    \end{enumerate}
    \textbf{Reflexe Pro :} Toujours créer les \textbf{Styles de paragraphe} (\texttt{Édition > Styles} ou \texttt{F3}) \textbf{avant} de saisir le texte (Titre1, CorpsDeTexte, Legende, PiedPage). Appliquer le style au cadre ou au texte sélectionné.
\end{methode}

\begin{exoresolu}[Mise en place d'un gabarit pour un journal scolaire "L'Écho du Lycée"]
    \textbf{Énoncé :} Le club presse veut un journal de 4 pages (format A3 plié en 2 = 4 pages A4). Page 1 : Une (Titre journal, Édito, Une photo). Pages 2-3 : Articles sur 3 colonnes (double page centrale). Page 4 : Vie du club, contacts. Mettez en place les Pages Maîtres et la structure de cadres pour les pages intérieures (2-3).
    \tcblower
    \textbf{Solution détaillée :}
    \begin{enumerate}
        \item \textbf{Configuration doc :} \texttt{Nouveau} > Taille A3 ($297 \times 420$ mm), Orientation Portrait, \textbf{4 pages}, \textbf{Facing pages} coché. Marges $15$ mm. Bords perdus $3$ mm.
        \item \textbf{Pages Maîtres (Double page centrale = Pages 2-3) :}
              \begin{itemize}
                  \item Ouvrir \texttt{Fenêtre > Pages Maîtres}. Sélectionner \texttt{Normal Left} (Page 2) et \texttt{Normal Right} (Page 3).
                  \item \textbf{Grille de colonnes :} \texttt{Fichier > Configurer le document > Guides} : Colonnes = $3$, Gouttière = $5$ mm.
                  \item \textbf{Pied de page commun :} Sur \texttt{Normal Left} ET \texttt{Normal Right} : Créer cadre texte bas (Hauteur $10$ mm, Largeur page moins marges). Insérer : \texttt{« L'Écho du Lycée — N°12 — Juin 2025 — Page »} + \texttt{Insertion > Marqueur > Numéro de page}. Aligner à gauche (page gauche) / à droite (page droite). Police \textit{Times New Roman}, Corps $8$ pt, Gris $60\%$. Style \texttt{PiedPage}.
                  \item \textbf{En-tête (Nom journal petit) :} Cadre texte haut, aligné extérieur (gauche p.2, droite p.3). Texte : \texttt{L'Écho du Lycée}. Style \texttt{EnteteJournal}.
              \end{itemize}
        \item \textbf{Structure Pages 2 \& 3 (Mode normal) :}
              \begin{itemize}
                  \item Sur Page 2 : Créer $3$ cadres texte alignés sur les guides de colonnes (Largeur cadre = largeur colonne). Les sélectionner dans l'ordre de lecture (Col 1, Col 2, Col 3) $\rightarrow$ \texttt{Édition > Chaîner les cadres}. Répéter sur Page 3.
                  \item Créer un grand \textbf{Cadre Image} en haut de la Page 2 (sur 3 colonnes, hauteur $\approx 80$ mm) pour la photo d'illustration. Option \texttt{Propriétés > Forme > Le texte contourne le cadre} (Habillage) activé plus tard.
              \end{itemize}
        \item \textbf{Vérification :} Cliquer dans le premier cadre de la page 2, taper du "Lorem Ipsum" long. Le texte doit couler vers cadre 2, puis 3, puis page 3 cadres 1, 2, 3. Les numéros de page s'affichent automatiquement en bas.
    \end{enumerate}
    \textbf{Conclusion :} La structure est industrielle. Les rédacteurs n'ont plus qu'à importer le texte (\texttt{Fichier > Importer > Texte}) dans le premier cadre. La cohérence visuelle (colonnes, pagination, en-têtes) est garantie sur tout le journal.
\end{exoresolu}

\courssub{Insertion, habillage et ajustement des images (Gestion des ressources)}

\begin{propriete}[Règles d'or pour les images en PAO professionnelle]
    \begin{enumerate}
        \item \textbf{Liaison, pas incorporation :} Scribus ne copie pas l'image dans le \texttt{.sla}. Il garde un \textbf{lien} (chemin relatif ou absolu). \textbf{Conséquence :} Le dossier \texttt{/Images} doit accompagner le \texttt{.sla} pour toute modification future.
        \item \textbf{Résolution effective (Effective PPI) :} C'est la résolution \textbf{après mise à l'échelle} dans le cadre. Cible : \textbf{$\ge 300$ DPI} (impression), $72/96$ DPI (écran). Vérifiable dans \texttt{Propriétés > Image} ou \texttt{Gestionnaire d'images (F12)}.
        \item \textbf{Mode couleur :} \textbf{CMJN} (Cyan, Magenta, Jaune, Noir) pour l'impression offset. \textbf{RVB} (Rouge, Vert, Bleu) pour écran/PDF web. Une image RVB dans un PDF/X-1a sera convertie (souvent mal) à l'export. \textbf{Idéal :} Convertir en CMJN (Profil \texttt{ISO Coated v2 300\%}) dans GIMP/Photoshop \textbf{avant} import.
        \item \textbf{Formats :} \texttt{TIFF} (sans perte, CMJN), \texttt{PSD} (calques), \texttt{JPEG} (photos, qualité max), \texttt{PNG} (transparence, logos, RVB), \texttt{SVG/PDF} (vectoriel, logos, graphiques).
        \item \textbf{Nommage :} Sans accents, sans espaces, minuscules (ex. \texttt{logo\_ecole\_cmjn.tif}, \texttt{photo\_kermesse\_01.jpg}).
    \end{enumerate}
\end{propriete}

\begin{methode}[Insérer, habiller et ajuster les images]
    \textbf{Objectif :} Placer des images sans les déformer, gérer le lien source et faire contourner le texte (habillage).
    \begin{enumerate}
        \item \textbf{Préparer les fichiers sources :} Rassembler dans dossier \texttt{/Images} (à côté du \texttt{.sla}). Nommage propre. Résolution $300$ DPI. Mode CMJN pour impression.
        \item \textbf{Insérer l'image :} Outil \texttt{Cadre Image} (\texttt{I} ou icône montagne). Dessiner le cadre à la taille souhaitée. \texttt{Clic droit > Obtenir une image} (\texttt{Ctrl+I}) $\rightarrow$ Sélectionner fichier.
        \item \textbf{Ajuster l'image au cadre (ou inversement) :} \texttt{Clic droit} (ou \texttt{Propriétés > Image}) :
              \begin{itemize}
                  \item \texttt{Ajuster l'image au cadre} : Échelle proportionnelle pour remplir le cadre (rognage possible).
                  \item \texttt{Ajuster le cadre à l'image} : Redimensionne le cadre aux dimensions natives de l'image (pixels $\rightarrow$ mm selon DPI).
                  \item \texttt{Échelle libre} : Déformer (interdit pour photos !).
              \end{itemize}
              \emph{Astuce :} Maintenir \texttt{Ctrl} + molette souris pour zoomer/déplacer l'image \emph{à l'intérieur} du cadre sans bouger le cadre.
        \item \textbf{Gérer l'habillage (Contour de texte) :} Sélectionner cadre image $\rightarrow$ \texttt{Propriétés (F2) > Onglet Forme (ou Contour)} $\rightarrow$ Cocher \texttt{Le texte contourne le cadre}. Choisir type :
              \begin{itemize}
                  \item \texttt{Contour du cadre} : Rectangle autour du cadre.
                  \item \texttt{Contour de l'image (Canal Alpha)} : Suit la transparence (PNG) ou chemin de détourage (Clipping Path). \textbf{Idéal pour logos/png sans fond.}
                  \item \texttt{Distance du texte} : Espace (ex. $3$ mm) entre image et texte.
              \end{itemize}
        \item \textbf{Vérifier le Gestionnaire d'images :} \texttt{Fichier > Gestionnaire d'images} (\texttt{F12}). Colonnes : \texttt{État} (OK/Manquant/Modifié), \texttt{Résolution effective}, \texttt{Espace couleur}. Si \texttt{Manquant} $\rightarrow$ \texttt{Rechercher} le fichier déplacé. Si \texttt{Modifié} $\rightarrow$ \texttt{Mettre à jour}.
    \end{enumerate}
    \textbf{Attention :} Dans Scribus, l'image n'est \textbf{pas intégrée} dans le fichier \texttt{.sla}. Si vous déplacez le dossier \texttt{/Images}, le lien casse. Toujours envoyer le dossier \texttt{Images} + fichier \texttt{.sla} (ou exporter en PDF pour diffusion finale).
\end{methode}

\begin{exoresolu}[Insertion du logo et d'une photo avec habillage pour l'affiche de la kermesse]
    \textbf{Énoncé :} Sur la page de garde du flyer (Page 1, format A5), vous devez placer :
    \begin{itemize}
        \item Le logo de l'école (\texttt{logo\_college\_ngaoundere.png}, fond transparent, $500 \times 500$ px) en haut à gauche, largeur $30$ mm.
        \item Une photo d'ambiance (\texttt{kermesse\_2024.jpg}, $3000 \times 2000$ px, RVB) en fond de page entière (avec bords perdus).
        \item Un cadre texte pour le titre "KERMESSE 2025" qui doit se superposer à la photo en bas à droite, le texte devant contourner un cercle imaginaire (simulé par un cadre image vide avec contour alpha).
    \end{itemize}
    Réalisez les opérations dans l'ordre et justifiez le choix de résolution pour la photo.
    \tcblower
    \textbf{Solution détaillée :}
    \begin{enumerate}
        \item \textbf{Photo de fond (Page 1) :}
              \begin{itemize}
                  \item Créer cadre image couvrant \textbf{toute la page y compris bords perdus} (Largeur $154$ mm, Hauteur $216$ mm, position $X=-3$ mm, $Y=-3$ mm par rapport au coin page physique).
                  \item \texttt{Clic droit > Obtenir une image} $\rightarrow$ \texttt{kermesse\_2024.jpg}.
                  \item \texttt{Clic droit > Ajuster l'image au cadre} (échelle proportionnelle, l'image remplit le cadre, parties rognées si ratio différent).
                  \item \textbf{Ordre des calques :} \texttt{Clic droit > Niveau > Envoyer au fond} (icône page avec flèche bas). Ainsi le texte passera au-dessus.
              \end{itemize}
        \item \textbf{Logo (Haut gauche) :}
              \begin{itemize}
                  \item Créer cadre image ($30 \times 30$ mm) positionné à $X=10$ mm, $Y=10$ mm (marge sécurité).
                  \item Insérer \texttt{logo\_college\_ngaoundere.png}.
                  \item \texttt{Ajuster l'image au cadre} (proportionnel). PNG transparent $\rightarrow$ pas de fond blanc.
                  \item \textbf{Résolution effective :} $500$ px / $30$ mm $\approx 500 / 1,18 \approx 423$ DPI. \textbf{OK} ($> 300$ DPI).
              \end{itemize}
        \item \textbf{Simulation habillage circulaire (Titre) :}
              \begin{itemize}
                  \item Créer un \textbf{cadre image vide} (sans image) de $60 \times 60$ mm, positionné bas droite (zone titre).
                  \item \texttt{Propriétés (F2) > Forme} : Choisir \texttt{Forme > Ellipse} (le cadre devient rond).
                  \item \texttt{Propriétés > Forme > Contour de texte} : Cocher \texttt{Le texte contourne le cadre} > Type \texttt{Contour du cadre}. Distance $3$ mm.
                  \item Créer le \textbf{Cadre Texte} du titre (large, recouvrant la zone). Le texte évitera automatiquement le cercle.
                  \item \emph{Astuce :} Mettre le cadre cercle \textbf{Invisible} (Couleur de fond : Aucune, Contour : Aucun) mais garder l'habillage actif.
              \end{itemize}
        \item \textbf{Justification résolution photo :} $3000$ px sur largeur utile $\approx 148$ mm ($14,8$ cm). Résolution = $3000 / (14,8 / 2,54) \approx 3000 / 5,83 \approx 514$ DPI. \textbf{Excellent} pour l'impression (qualité photo). Le fichier est en RVB $\rightarrow$ À l'export PDF (Impression), Scribus/Imprimeur convertira en CMJN (perte possible de vifs). Idéal : convertir en CMJN dans GIMP/Photoshop avant import.
    \end{enumerate}
    \textbf{Conclusion :} Les éléments graphiques sont en place. La structure de calques (Fond photo > Logo > Cercle habillage > Texte Titre) est respectée. Le gestionnaire d'images signale "OK" pour les deux fichiers.
\end{exoresolu}

\courssub{Mise en forme du texte : Styles de paragraphe et de caractère}

\begin{definition}[Style de paragraphe]
    Ensemble de paramètres de formatage (police, corps, interlignage, alignement, alinéa, espacements avant/après, couleur, tabulations, césure) appliqué à un paragraphe entier. Permet la cohérence globale et la modification en cascade.
\end{definition}

\begin{definition}[Style de caractère]
    Ensemble de paramètres (gras, italique, couleur, police, taille relative) appliqué à une \textbf{sélection de caractères} à l'intérieur d'un paragraphe (ex. mot en gras, référence en italique). Ne doit pas contenir de paramètres de paragraphe (alignement, interlignage).
\end{definition}

\begin{methode}[Mettre en forme le texte avec les Styles (Cohérence et gain de temps)]
    \textbf{Objectif :} Appliquer une charte graphique rigoureuse via des \textbf{Styles de paragraphe} et de \textbf{caractère}, plutôt que du formatage manuel.
    \begin{enumerate}
        \item \textbf{Ouvrir l'éditeur de styles :} \texttt{Édition > Styles} (ou \texttt{F3}).
        \item \textbf{Créer les styles de base (Hiérarchie par héritage) :}
              \begin{itemize}
                  \item \texttt{Style\_Base} (Base) : Police \textit{Source Serif Pro} / \textit{DejaVu Serif} / \textit{Times New Roman}, Corps $10$ pt, Interlignage $1,3$ (ou $13$ pt fixe), Alignement \texttt{Justifié}, Alinéa premier $4$ mm, Espace après paragraphe $2$ mm. Couleur \texttt{Noir} (ou \texttt{TextePrincipal}).
                  \item \texttt{Titre\_Principal} (basé sur \texttt{Style\_Base}) : Police \textit{Montserrat Bold} / \textit{Oswald Bold} / \textit{Arial Black}, Corps $24$ pt, Interlignage $1,1$ (ou $26$ pt), Alignement \texttt{Centré}, Couleur \texttt{BleuEcole} (ex. \texttt{C=100 M=60 J=0 N=10}), Espace avant $10$ mm, Après $6$ mm. \textbf{Pas d'alinéa}.
                  \item \texttt{Sous\_Titre} (basé sur \texttt{Titre\_Principal}) : Corps $14$ pt, Gras, Alignement \texttt{Gauche}, Couleur \texttt{GrisFonce}.
                  \item \texttt{Corps\_Journal} (basé sur \texttt{Style\_Base}) : Corps $9,5$ pt, Interlignage $12,5$ pt fixe, Alinéa $3$ mm.
                  \item \texttt{Corps\_Premier} (basé sur \texttt{Corps\_Journal}) : \textbf{Alinéa $0$ mm} (pour le 1er paragraphe après titre).
                  \item \texttt{Liste\_Puces} (basé sur \texttt{Corps\_Journal}) : Retrait gauche $8$ mm, Puce « • », Espace après $1$ mm.
                  \item \texttt{Legende\_Photo} (basé sur \texttt{Style\_Base}) : Corps $8$ pt, Italique, Alignement \texttt{Centré}, Couleur \texttt{GrisMoyen}, Espace avant $1$ mm.
              \end{itemize}
        \item \textbf{Créer des styles de caractère (pour l'intérieur du paragraphe) :}
              \begin{itemize}
                  \item \texttt{Emphase} : Italique.
                  \item \texttt{GrasImportant} : Gras, Couleur \texttt{RougeAccent}.
                  \item \texttt{LienWeb} : Souligné, Bleu.
              \end{itemize}
        \item \textbf{Appliquer :} Sélectionner le cadre texte (applique à tout le cadre) ou placer curseur dans paragraphe / sélectionner paragraphes $\rightarrow$ Double-clic sur le nom du style dans le panneau \texttt{Propriétés > Texte > Style de paragraphe} ou panneau \texttt{Styles (F3)}.
        \item \textbf{Modifier globalement :} Double-clic sur le style \texttt{Style\_Base} $\rightarrow$ Changer Police en \textit{Libre Baskerville} $\rightarrow$ \textbf{Tout le document met à jour instantanément}.
    \end{enumerate}
    \textbf{Règle d'or :} \textbf{Zéro formatage manuel} (pas de \texttt{Ctrl+B}, \texttt{Ctrl+I}, changement taille police via barre d'outils). Tout passe par les Styles. C'est la condition \emph{sine qua non} pour la note "Qualité de la mise en page" au BEPC.
\end{methode}

\begin{exoresolu}[Application d'une charte graphique au journal "L'Écho du Lycée"]
    \textbf{Énoncé :} La charte graphique du journal impose :
    \begin{itemize}
        \item Police de titre : \textbf{Oswald Bold}, Couleur \textbf{Bleu Marine (C=100, M=80, J=0, N=50)}.
        \item Police de corps : \textbf{Libre Baskerville}, Corps $9,5$ pt, Interlignage $12,5$ pt, Justifié, Alinéa $3$ mm.
        \item Citation (encadré) : \textit{Libre Baskerville Italique}, Corps $10$ pt, Aligné centre, Marge G/D $15$ mm, Fond \textbf{Jaune Pâle (C=0, M=0, J=15, N=0)}.
    \end{itemize}
    Créez les styles nécessaires dans Scribus et appliquez-les à un article type (Titre, Chapeau, 3 paragraphes, Citation, Légende photo).
    \tcblower
    \textbf{Solution détaillée :}
    \begin{enumerate}
        \item \textbf{Création Styles (Édition > Styles > Nouveau) :}
              \begin{enumerate}
                  \item \texttt{Style\_Corps\_Journal} : \texttt{Police: Libre Baskerville Regular}, \texttt{Taille: 9.5 pt}, \texttt{Interlignage: Fixe 12.5 pt}, \texttt{Alignement: Justifié}, \texttt{Alinéa premier: 3 mm}, \texttt{Couleur: Noir (C=0 M=0 J=0 N=100)}. \textbf{Valider}.
                  \item \texttt{Style\_Chapeau} (basé sur \texttt{Style\_Corps\_Journal}) : \texttt{Police: Libre Baskerville Bold Italic}, \texttt{Taille: 11 pt}, \texttt{Interlignage: 14 pt}, \texttt{Alinéa: 0 mm}, \texttt{Espace après: 4 mm}.
                  \item \texttt{Style\_Titre\_Article} : \texttt{Police: Oswald Bold}, \texttt{Taille: 18 pt}, \texttt{Interlignage: 22 pt}, \texttt{Alignement: Gauche}, \texttt{Couleur: BleuMarine (C=100 M=80 J=0 N=50)}, \texttt{Espace avant: 6 mm, Après: 3 mm}.
                  \item \texttt{Style\_Citation} : \texttt{Police: Libre Baskerville Italic}, \texttt{Taille: 10 pt}, \texttt{Interlignage: 13 pt}, \texttt{Alignement: Centré}, \texttt{Marge Gauche: 15 mm, Marge Droite: 15 mm}, \texttt{Fond: JaunePale (C=0 M=0 J=15 N=0)}, \texttt{Espace avant/après: 4 mm}.
                  \item \texttt{Style\_Legende} : \texttt{Police: Libre Baskerville Regular}, \texttt{Taille: 8 pt}, \texttt{Italique: Oui}, \texttt{Alignement: Centré}, \texttt{Couleur: Gris 60\%}.
              \end{enumerate}
        \item \textbf{Application sur Page 2 (Cadre texte chaîné Col 1) :}
              \begin{itemize}
                  \item Ligne 1 : "La rentrée sportive bat son plein" $\rightarrow$ Appliquer \texttt{Style\_Titre\_Article}.
                  \item Ligne 2 : "Les équipes de football et de handball..." $\rightarrow$ Appliquer \texttt{Style\_Chapeau}.
                  \item Paragraphe 1 : Texte long $\rightarrow$ Appliquer \texttt{Style\_Corps\_Premier} (sans alinéa).
                  \item Paragraphes 2 et 3 : Appliquer \texttt{Style\_Corps\_Journal} (avec alinéa $3$ mm).
                  \item Citation au milieu : "Le sport unit les cœurs..." $\rightarrow$ Appliquer \texttt{Style\_Citation}. Vérifier fond jaune et marges $15$ mm.
                  \item Photo en bas colonne : Insérer image $\rightarrow$ Cadre légende dessous $\rightarrow$ Appliquer \texttt{Style\_Legende}.
              \end{itemize}
        \item \textbf{Test modification globale :} Double-clic \texttt{Style\_Corps\_Journal} $\rightarrow$ Taille $10$ pt $\rightarrow$ Valider. Tous les paragraphes corps grossissent, la pagination peut décaler (reflux texte), l'habillage s'adapte.
    \end{enumerate}
    \textbf{Conclusion :} La charte est respectée techniquement. La modification de la taille de police de corps a propagé le changement sur les 3 colonnes et 2 pages instantanément, prouvant l'efficacité des styles.
\end{exoresolu}

\courssub{Finalisation, vérification (Preflight), exportation PDF et impression}

\begin{definition}[Preflight (Contrôle amont / Vérification pré-presse)]
    Processus automatisé de vérification technique d'un document PAO avant export : détection des polices manquantes/non intégrables, images basse résolution (< 300 DPI) ou RVB (pour PDF/X), textes masqués (overset), bords perdus manquants, transparences non aplaties. \textbf{Obligatoire avant tout export professionnel.}
\end{definition}

\begin{definition}[PDF/X-1a:2001 (ISO 15931-1)]
    Standard ISO pour l'échange de fichiers prêts pour l'impression (Blind Exchange). Contraintes : \textbf{Toutes les polices intégrées}, \textbf{Toutes les images en CMJN (ou Niveaux de gris)}, \textbf{Pas de transparence vive}, \textbf{Pas de RGB}, \textbf{Pas de multimédia/JS}, \textbf{Bords perdus et marques de coupe inclus}. C'est le format le plus demandé par les imprimeurs au Cameroun (Imprim'Plus, Sodicam, etc.).
\end{definition}

\begin{methode}[Finaliser, vérifier (Preflight), exporter en PDF et imprimer]
    \textbf{Objectif :} Garantir un fichier final sans erreur technique et produire le PDF standard pour l'imprimeur ou le web.
    \begin{enumerate}
        \item \textbf{Relecture "Preflight" (Vérification pré-presse) :} \texttt{Fichier > Vérifier le document} (ou icône avion/feuille dans barre d'état).
              \begin{itemize}
                  \item \textbf{Polices :} Toutes \texttt{Intégrées} (Embedded) ou \texttt{Vectorisées} (Outline). Aucune "Manquante" ou "Non intégrable" (licence restrictive).
                  \item \textbf{Images :} Résolution effective $\ge 300$ DPI (impression). Mode couleur \texttt{CMJN} (ou \texttt{Niveaux de gris}). Pas d'images \texttt{RVB} si PDF/X-1a. Pas d'images \texttt{Manquantes}.
                  \item \textbf{Débordements (Overset text) :} Icône \texttt{œil barré} ou croix rouge sur cadre texte $\rightarrow$ Texte tronqué. \textbf{Résoudre :} Agrandir cadre, réduire police, lier à cadre suivant, ou couper texte.
                  \item \textbf{Bords perdus :} Vérifier visuellement (Mode \texttt{Aperçu} / \texttt{W}) que les fonds/images touchent le trait rouge (bord perdu).
              \end{itemize}
        \item \textbf{Corriger les anomalies :} Traiter chaque erreur (rouge) et avertissement (orange). Ne pas exporter tant qu'il reste des erreurs rouges.
        \item \textbf{Exporter en PDF :} \texttt{Fichier > Exporter > Enregistrer sous PDF}.
              \begin{itemize}
                  \item \textbf{Préréglage :} \texttt{PDF/X-1a:2001} (Standard imprimerie, tout en CMJN, polices embarquées, pas de transparence vive, pas de JS/Video) \textbf{OU} \texttt{PDF 1.4/1.5} (Généraliste, garde transparences, RVB possible pour diffusion web/email).
                  \item \textbf{Onglet Général :} Pages : \texttt{Toutes} / \texttt{Étendue}. \textbf{Décocher \texttt{Pages face à face}} pour PDF simple page (imprimeur préfère souvent pages simples + bords perdus). \textbf{Cocher :} \texttt{Inclure les bords perdus du document (Use Document Bleeds)}. \texttt{Marques de coupe (Crop marks)} : Oui, Décalage \textbf{3 mm} (doit être $\ge$ bord perdu).
                  \item \textbf{Onglet Polices :} \texttt{Intégrer toutes : OUI}. \texttt{Sous-ensemble (Subset) : OUI} (optimise taille fichier).
                  \item \textbf{Onglet Couleur :} \texttt{Sortie destinée à : Imprimante (CMJN)}. \texttt{Convertir les couleurs en : CMJN}. Profil : \texttt{ISO Coated v2 300\% (ECI)} (Standard Europe/Afrique francophone).
                  \item \textbf{Onglet Compression :} Images Couleur/Niveaux de gris : \texttt{Bicubique} / $300$ DPI / \texttt{Seuil 450 DPI}. Monochrome : $1200$ DPI. Qualité \texttt{Élevée / Maximum}.
              \end{itemize}
        \item \textbf{Enregistrer et Vérifier le PDF :} Ouvrir le PDF dans \texttt{Adobe Acrobat Reader} / \texttt{Okular} / \texttt{Foxit}. Vérifier : Pages complètes, marques de coupe visibles, texte sélectionnable (pas vectorisé en courbes sauf logos), couleurs correctes.
        \item \textbf{Impression locale (Test / Épreuvage) :} \texttt{Fichier > Imprimer} $\rightarrow$ Imprimante \texttt{Microsoft Print to PDF} ou physique. \textbf{Échelle :} \texttt{Taille réelle (100\%)}. Imprimer une "épreuve" sur papier ordinaire pour valider mise en page, pli, lisibilité.
    \end{enumerate}
    \textbf{Livrable final :} Fichier \texttt{.sla} (source modifiable) + Dossier \texttt{/Images} + Fichier \texttt{.pdf} (PDF/X-1a pour imprimeur, PDF 1.4 pour diffusion email/site web).
\end{methode}

\begin{exoresolu}[Export PDF/X-1a pour l'imprimeur "Imprim'Plus Yaoundé"]
    \textbf{Énoncé :} Votre flyer kermesse (2 pages A4 face à face, bords perdus 3mm) est finalisé. L'imprimeur "Imprim'Plus" à Yaoundé demande un fichier \textbf{PDF/X-1a:2001}, pages simples (non face-à-face), avec marques de coupe décalées de 3mm, polices intégrées, images en CMJN à 300 DPI. Le document utilise la police \textit{Montserrat} (Google Font, libre) et des photos RVB. Détaillez la procédure d'export exacte dans Scribus et les vérifications post-export.
    \tcblower
    \textbf{Solution détaillée :}
    \begin{enumerate}
        \item \textbf{Preflight (Contrôle amont) :} \texttt{Fichier > Vérifier le document}.
              \begin{itemize}
                  \item \textbf{Police Montserrat :} Statut "OK - Intégrable". (Licence SIL Open Font License $\rightarrow$ autorise intégration PDF).
                  \item \textbf{Photos RVB :} Avertissement "Image RVB dans document CMJN". \textbf{Action :} Onglet \texttt{Couleur} de l'export gérera la conversion. Ou convertir manuellement dans GIMP avant import (mieux maîtrisé). On note l'avertissement, on continue.
                  \item \textbf{Résolution :} Vérifier $> 300$ DPI effective. OK.
                  \item \textbf{Texte masqué :} Aucun cadre avec croix rouge. OK.
              \end{itemize}
        \item \textbf{Export PDF :} \texttt{Fichier > Exporter > Enregistrer sous PDF}.
              \item \textbf{Préréglage :} Sélectionner \textbf{PDF/X-1a:2001} (Charge les bonnes options de base : Compatibilité 1.3, Pas de transparence, Pas de RGB, Polices embarquées).
              \item \textbf{Onglet Général :}
                    \begin{itemize}
                        \item \texttt{Pages} : \textbf{Toutes}. \textbf{Important :} Décocher \texttt{Pages face à face} (L'imprimeur veut pages simples : Page 1, Page 2 séparées).
                        \item \texttt{Bords perdus} : Cocher \textbf{Utiliser les bords perdus du document}. Marques de coupe : \textbf{Cocher}. Décalage : \textbf{3,00 mm} (doit être $\ge$ bord perdu).
                        \item \texttt{Options} : Décocher \texttt{Inclure les annotations}, \texttt{Inclure les formulaires}.
                    \end{itemize}
              \item \textbf{Onglet Polices :} \texttt{Intégrer toutes : OUI}. \texttt{Sous-ensemble (Subset) : OUI} (Montserrat a beaucoup de glyphes, subset réduit le poids). \texttt{Vectoriser : NON} (garde texte sélectionnable/recherchable).
              \item \textbf{Onglet Couleur :} \texttt{Sortie destinée à : Imprimante (CMJN)}. \texttt{Convertir les couleurs en : CMJN}. \texttt{Profil de sortie : ISO Coated v2 300\% (ECI)}. \texttt{Simuler l'impression sur écran : Non}.
              \item \textbf{Onglet Compression :} Méthode \texttt{Échantillonnage bicubique vers le bas}. Résolution \texttt{300 dpi}. Seuil \texttt{450 dpi}. Compression \texttt{Automatique (JPEG/ZIP)}. Qualité \texttt{Élevée / Maximum}.
              \item \textbf{Enregistrer :} \texttt{Flyer\_Kermesse\_IMPRIM\_PLUS\_20250615.pdf}.
        \item \textbf{Vérification Post-Export (Acrobat Reader / Preflight Acrobat) :}
              \begin{itemize}
                  \item Ouvrir PDF. \texttt{Fichier > Propriétés > Description} : Vérifier "Conforme PDF/X-1a:2001".
                  \item \texttt{Outils > Protection et normalisation > Preflight (Contrôle amont)} > Profil \texttt{PDF/X-1a}. Lancer. \textbf{Résultat :} "Aucun problème trouvé" (Vert).
                  \item Visuel : Marques de coupe aux 4 coins, décalées de 3mm du bord page. Fond photo va jusqu'au trait de coupe. Texte "KERMESSE 2025" bien dans marge sécurité ($>5$mm du trait de coupe).
                  \item Texte sélectionnable (Copier/Coller fonctionne).
              \end{itemize}
    \end{enumerate}
    \textbf{Conclusion :} Le fichier est techniquement conforme aux exigences de l'imprimeur. Il est prêt à être envoyé par WeTransfer/Email (taille $\approx 5-10$ Mo). On archive le \texttt{.sla} + \texttt{/Images} pour modification future (édition 2026).
\end{exoresolu}

\courssub{Démarche de production et justification (Compétence BEPC : Rédiger et justifier)}

\begin{methode}[Justifier sa démarche de production PAO (Écriture réflexive / Oral)]
    \textbf{Objectif :} Structurer l'explication orale ou écrite de vos choix techniques et graphiques (Compétence "Rédiger et justifier une démarche"). Utiliser la structure : \textbf{Contexte $\rightarrow$ Contraintes $\rightarrow$ Choix techniques $\rightarrow$ Validation}.
    \begin{enumerate}
        \item \textbf{Contextualiser (La situation) :} "Dans le cadre de la kermesse du collège, il fallait informer les parents et élèves (cible) sur la date, le lieu, le programme. Support choisi : Flyer A5 recto-verso, imprimé 500 ex. chez Imprim'Plus."
        \item \textbf{Lister les contraintes techniques (Le cahier des charges) :}
              \begin{itemize}
                  \item Format imposé (A5), Bords perdus 3mm, CMJN, 300 DPI, PDF/X-1a.
                  \item Charte graphique école : Bleu Marine \#003366 (C100 M80 N50), Police Oswald/Libre Baskerville.
                  \item Contenu : Logo, Photo ambiance, Titre, Programme (liste), Contacts, Partenaires (logos bas de page).
              \end{itemize}
        \item \textbf{Expliquer les choix d'interface et structure (Le "Comment") :}
              \begin{itemize}
                  \item \texttt{Pages Maîtres} : "J'ai utilisé un maître unique pour le pied de page (contacts + numérotation) afin d'assurer la cohérence sur les 2 pages et gagner du temps."
                  \item \texttt{Cadres chaînés} : "Pour le programme (texte long), j'ai chainé 3 cadres en colonnes pour une lecture journalistique fluide."
                  \item \texttt{Styles} : "Création de 5 styles (Titre, Sous-titre, Corps, Liste, Pied) pour appliquer la charte en 1 clic et modifier globalement si besoin."
              \end{itemize}
        \item \textbf{Justifier les choix graphiques (Le "Pourquoi" esthétique/ergonomique) :}
              \begin{itemize}
                  \item \textbf{Hiérarchie visuelle :} "Titre 24pt Bleu Marine en haut $\rightarrow$ Accroche regard. Photo fond traitée en transparence 30\% (ou voile blanc sur photo) $\rightarrow$ Lisibilité texte préservée."
                  \item \textbf{Habillage :} "Cercle d'habillage invisible pour le titre $\rightarrow$ Dynamisme, casse la rigidité des colonnes."
                  \item \textbf{Couleur :} "Passage RVB $\rightarrow$ CMJN profil ISO Coated v2 $\rightarrow$ Fidélité couleurs impression (le bleu marine ne vire pas au violet)."
              \end{itemize}
        \item \textbf{Valider la qualité finale (Le "Preuve") :} "Preflight Scribus OK (0 erreur). Export PDF/X-1a validé par Preflight Acrobat. Épreuve papier imprimée au collège : pli central net, texte lisible, couleurs conformes à l'écran (épreuvage)."
    \end{enumerate}
    \textbf{Phrase type de conclusion :} "Cette démarche, de la configuration du gabarit à l'export normé, garantit un document professionnel, fidèle à la charte de l'établissement, techniquement irreprochable pour l'imprimeur et modifiable pour l'an prochain."
\end{methode}

\begin{exoresolu}[Rédaction de la justification pour l'examen pratique (BEPC)]
    \textbf{Énoncé :} À l'oral de l'examen, le jury vous demande : \textit{"Présentez la démarche que vous avez suivie pour réaliser le flyer de la kermesse, en justifiant vos choix techniques majeurs (gestion des pages, styles, images, export)."} Rédigez la réponse structurée que vous donnerez (environ 2 minutes de parole).
    \tcblower
    \textbf{Solution détaillée (Réponse type élève) :}
    \begin{quote}
        \textbf{1. Contexte et Cahier des charges :} « J'ai réalisé un flyer A5 recto-verso pour la kermesse du Collège Bilingue de Ngaoundéré. L'imprimeur imposait un PDF/X-1a, bords perdus 3 mm, CMJN, 300 DPI. La charte graphique impose le Bleu Marine (C100 M80 N50) et les polices Oswald / Libre Baskerville. »
        \textbf{2. Structure du document (Pages Maîtres) :} « J'ai configuré un document Scribus 2 pages A4 face-à-face (format ouvert) avec bords perdus. J'ai créé un \textbf{Maître unique} pour automatiser le pied de page (contacts école, numéro de page) sur les deux pages. Cela assure la cohérence et évite la ressaisie. »
        \textbf{3. Mise en page et Styles :} « J'ai structuré la page de garde avec des cadres distincts : fond photo, logo, titre. Pour la page 2 (programme), j'ai utilisé \textbf{3 cadres texte chaînés en colonnes} pour une lecture fluide. J'ai défini \textbf{5 styles de paragraphe} (Titre, Chapeau, Corps, Liste, Pied) basés sur un style "Base". Cela m'a permis d'appliquer la charte instantanément et de tester la taille de police du corps (9,5 $\rightarrow$ 10 pt) sur tout le document en une modification. »
        \textbf{4. Traitement des images :} « La photo de fond (RVB, 514 DPI effective) a été envoyée au fond. Le logo (PNG transparent, 423 DPI) placé en haut à gauche. J'ai simulé un habillage circulaire pour le titre via un cadre ellipse invisible avec contour de texte actif, ce qui dynamise la une. \textbf{Point de vigilance :} J'ai converti les images en CMJN (Profil ISO Coated v2) à l'export pour garantir le rendu "Bleu Marine" attendu, car le RVB donne un bleu plus vif à l'écran mais décevant en impression. »
        \textbf{5. Finalisation et Export :} « J'ai lancé le \textbf{Preflight (Vérifier le document)} : 0 erreur (polices intégrables, images OK, pas de texte masqué). Export \textbf{PDF/X-1a:2001}, pages simples, marques de coupe 3mm, polices en sous-ensemble. Validation finale par Preflight Acrobat : "Conforme". Épreuve papier validée au collège. »
        \textbf{Conclusion :} « Ce flux de travail (Gabarit $\rightarrow$ Styles $\rightarrow$ Chaînage $\rightarrow$ Preflight $\rightarrow$ PDF/X) est reproductible, professionnel et respecte scrupuleusement les contraintes de l'imprimeur et de la charte graphique. »
    \end{quote}
    \textbf{Évaluation :} Réponse complète, vocabulaire technique précis (Preflight, PDF/X-1a, Chaînage, Style basé sur, Profil ICC, Bords perdus, Sous-ensemble), structure logique (Contexte $\rightarrow$ Technique $\rightarrow$ Validation). Note maximale attendue.
\end{exoresolu}

\begin{attention}[Les 5 erreurs qui coûtent des points au BEPC]
    \begin{enumerate}
        \item \textbf{Oublier les bords perdus (Bleeds) à l'export :} Configurer les bords perdus dans le document (3 mm) mais \textbf{ne pas cocher "Utiliser les bords perdus du document" dans l'onglet Général de l'export PDF}. Résultat : Le PDF est rogné au format fini, le fond blanc apparaît sur les bords après coupe. \textbf{Perte :} 2 à 4 pts (Qualité technique / Conformité imprimeur).
        \item \textbf{Texte masqué (Overset) non détecté :} Un cadre texte contient plus de texte qu'il n'en peut afficher (croix rouge / œil barré). L'élève exporte sans vérifier le Preflight. Résultat : La fin de l'article manque sur le PDF imprimé. \textbf{Perte :} 3 à 5 pts (Fonctionnalité / Lisibilité). \emph{Reflexe :} Toujours lancer \texttt{Fichier > Vérifier le document} avant export.
        \item \textbf{Images en RVB dans un PDF/X-1a (Impression) :} Laisser des photos RVB (profil sRGB) alors que la norme PDF/X-1a impose le CMJN. L'imprimeur convertira "à la volée" (souvent mal : verts ternes, bleus violacés). \textbf{Perte :} 2 pts (Gestion couleur). \emph{Solution :} Onglet Couleur export $\rightarrow$ Convertir en CMJN (Profil ISO Coated v2) OU convertir les images dans GIMP avant import.
        \item \textbf{Formatage manuel au lieu des Styles :} Mettre en gras, changer police, taille, couleur via la barre d'outils paragraphe par paragraphe. Au moindre changement de consigne (ex. "Police finale : Arial 11pt"), l'élève doit tout refaire. À l'oral, incapacité à expliquer la gestion de la charte graphique. \textbf{Perte :} 3 pts (Méthodologie / Compétence "Justifier démarche") + temps perdu en examen.
        \item \textbf{Polices non intégrées (ou remplacées) dans le PDF :} Utiliser une police exotique (téléchargée sur DaFont) sans vérifier "Intégrable" dans le Preflight, ou cocher "Vectoriser les polices" (Outline) ce qui rend le texte non sélectionnable/non accessible. Résultat : Chez l'imprimeur, police substituée par Courier New $\rightarrow$ Mise en page cassée. \textbf{Perte :} 3 pts (Interopérabilité / Qualité finale). \emph{Règle :} Polices Google Fonts (OFL) ou système standard + \texttt{Intégrer toutes / Sous-ensemble} à l'export.
    \end{enumerate}
\end{attention}

\courssub{Exercices d'entraînement (Type BEPC Pratique)}

\begin{exercice}[Création d'une carte de visite - Niveau 1]
    \textbf{Consigne :} Réaliser une carte de visite standard ($85 \times 55$ mm) pour le Proviseur du Lycée de Bafoussam.
    \begin{enumerate}
        \item Nouveau document : Format personnalisé $85 \times 55$ mm, Orientation Paysage, Marges $5$ mm, Bords perdus $3$ mm.
        \item Page Maître : Pied de page avec adresse école (style \texttt{PiedPage} : Police Arial 7 pt, Gris 60\%, Centré).
        \item Page 1 (Recto) : Logo école (haut gauche, 20 mm), Nom "PROVISEUR" (Style \texttt{Titre} : Oswald Bold 18 pt, Bleu Marine), Nom prénom (Style \texttt{Corps} : Libre Baskerville 11 pt), Titre/Grade (Style \texttt{SousTitre} : Italique 10 pt), Tél/Email (Style \texttt{Corps} 9 pt).
        \item Page 2 (Verso) : Plan d'accès simplifié (forme vectorielle) + QR Code (image) centré.
        \item Export PDF/X-1a avec marques de coupe 3 mm. Vérifier Preflight.
    \end{enumerate}
\end{exercice}

\begin{exercice}[Mise en page d'une affiche A3 - Niveau 2]
    \textbf{Consigne :} Concevoir l'affiche A3 (Portrait) de la "Journée de l'Excellence" du Collège.
    \begin{itemize}
        \item Contraintes : Fond perdu 3 mm. Photo fond pleine page (fichier fourni \texttt{excellence.jpg}, 4000 px larg). Titre principal "JOURNÉE DE L'EXCELLENCE" (Oswald Bold 48 pt, Blanc, Ombre portée). Date/Lieu (Montserrat 24 pt, Jaune Or). Liste des lauréats en 2 colonnes (Libre Baskerville 14 pt, Blanc, interlignage 18 pt). Logos partenaires bas de page (alignés, distribués horizontalement).
        \item Gestion calques : Calque "Fond" (verrouillé), Calque "Contenu", Calque "Logos".
        \item Habillage : Texte contourne un cadre ellipse invisible au centre pour aérer.
        \item Export PDF/X-1a + PDF 1.5 (pour écran).
    \end{itemize}
\end{exercice}

\begin{corrige}[Exercice Carte de visite]
    \textbf{Points clés de la correction :}
    \begin{itemize}
        \item Configuration doc : Taille $85 \times 55$ mm, Paysage, Marges 5 mm, Bleeds 3 mm. Unités mm.
        \item Maître : Pied de page appliqué aux 2 pages. Numérotation inutile (carte = 1 exemplaire unique), mais bon réflexe.
        \item Recto : Cadres positionnés aux coordonnées précises (X, Y via Propriétés). Alignement gauche logo/titre. Styles créés \textbf{avant} saisie.
        \item Verso : QR Code inséré comme image. Centré via \texttt{Aligner et Distribuer} (Centrer sur page).
        \item Export : Preflight OK. PDF/X-1a, Pages simples, Marques coupe 3mm, Polices subset, CMJN ISO Coated v2.
        \item \textbf{Piège évité :} Ne pas mettre de pli (pas de Facing pages pour une carte simple). Ne pas oublier le verso dans l'export (Pages : Toutes).
    \end{itemize}
\end{corrige}

\begin{corrige}[Exercice Affiche A3]
    \textbf{Points clés de la correction :}
    \begin{itemize}
        \item Photo fond : Cadre $303 \times 426$ mm (A3 + bleeds), position $-3, -3$. Image 4000 px / 297 mm $\approx 342$ DPI $\rightarrow$ OK. Envoyée au fond. Calque "Fond" verrouillé.
        \item Titre : Cadre texte large, style \texttt{TitreAffiche} (Blanc, Ombre portée via \texttt{Propriétés > Texte > Effets} ou duplication cadre décalé en noir derrière). Centré horizontalement, placé haut (Y $\approx$ 40 mm).
        \item Date/Lieu : Style \texttt{InfoAffiche} (Jaune Or C=0 M=20 J=100 N=0). Centré sous titre.
        \item Liste lauréats : 2 cadres texte chaînés, alignés sur guides colonnes (2 colonnes, gouttière 10 mm). Style \texttt{Laureat} (Blanc, 14 pt, Interlignage 18 pt, Aligné gauche). Habillage : Cadre ellipse invisible centre page (X=136.5, Y=213, Diam 80 mm) avec contour texte distance 5 mm.
        \item Logos partenaires : Calque "Logos". Cadres images alignés bas (Y $\approx$ 380 mm). \texttt{Aligner et Distribuer} : Aligner bas, Distribuer horizontalement centres.
        \item Export : 2 fichiers. 1. \texttt{Affiche\_Excellence\_IMPRIM.pdf} (PDF/X-1a, CMJN, Marques coupe). 2. \texttt{Affiche\_Excellence\_WEB.pdf} (PDF 1.5, RVB autorisé, Pas de marques coupe, Compression 150 DPI pour légèreté).
    \end{itemize}
\end{corrige}

\newpage

\coursec{Exercices}

\courssub{Je vérifie mes connaissances}

\begin{exercice}[QCM : Identification d'un logiciel de PAO]
Parmi les logiciels suivants, lequel est un logiciel de Publication Assistée par Ordinateur (PAO) dédié à la mise en page professionnelle ?
\begin{enumerate}
    \item Microsoft Word
    \item Microsoft Publisher
    \item Microsoft Excel
    \item Microsoft Paint
\end{enumerate}
Justifie ton choix en expliquant la différence fondamentale entre un traitement de texte et un logiciel de PAO.
\end{exercice}

\begin{exercice}[Vrai ou Faux justifié]
Détermine si les affirmations suivantes sont vraies ou fausses. Justifie chaque réponse.
\begin{enumerate}
    \item Dans un logiciel de PAO, le texte s'écrit directement sur la page comme dans un traitement de texte, sans objet conteneur.
    \item Les repères (marges, colonnes, gouttières) sont imprimés sur le document final.
    \item Un fichier \texttt{.pub} (format natif Publisher) peut être ouvert et modifié directement dans Microsoft Word sans conversion.
    \item L'exportation en PDF est le format recommandé pour l'impression professionnelle car il fige la mise en page.
\end{enumerate}
\end{exercice}

\begin{exercice}[Définitions et vocabulaire technique]
Donne une définition précise et concise des termes suivants dans le contexte de la PAO :
\begin{enumerate}
    \item \textbf{Gouttière} (ou espace inter-colonnes)
    \item \textbf{Repères de rognage} (traits de coupe)
    \item \textbf{Volet des pages} (dans l'interface de Publisher)
    \item \textbf{Zone de texte liée} (chaînage de zones de texte)
\end{enumerate}
\end{exercice}

\begin{exercice}[Calcul direct : Dimensions d'un dépliant]
Un dépliant publicitaire est réalisé sur une feuille A4 (21 cm $\times$ 29,7 cm) en orientation \textbf{paysage}. Il comporte 3 volets (2 plis) avec des gouttières de 0,5 cm entre les volets et des marges gauche et droite de 1 cm chacune.
Calcule la largeur utile (largeur d'impression) de \textbf{chaque volet} en centimètres.
\end{exercice}

\courssub{Je m'entraîne}

\begin{exercice}[Légender l'interface de Publisher]
L'image ci-dessous représente une capture d'écran simplifiée de l'interface de Microsoft Publisher au démarrage.
\begin{center}
\begin{tikzpicture}[font=\small]
% Fenêtre principale
\draw[thick, popdark] (0,0) rectangle (16,10);
% Barre de titre
\draw[popblue, fill=popblue!10] (0,9.5) rectangle (16,10) node[midway, above=2pt] {Barre de titre - Nom du fichier};
% Ruban
\draw[poporange, fill=poporange!10] (0,8.5) rectangle (16,9.5) node[midway] {RUBAN (Onglets : Accueil, Insertion, Mise en page, etc.)};
% Volet des pages (gauche)
\draw[popgreen, fill=popgreen!10] (0,0.5) rectangle (2.5,8.5) node[midway] {VOLET DES PAGES};
% Zone de travail (centre)
\draw[poppurple, fill=poppurple!10] (2.5,0.5) rectangle (13.5,8.5) node[midway] {ZONE DE TRAVAIL (Page courante)};
% Repères sur la page
\draw[dashed, popmuted] (3.5,1.5) rectangle (12.5,7.5) node[midway, align=center] {Marges intérieures\\ (Zone imprimable)};
\draw[dashed, popmuted] (3,1) rectangle (13,8) node[midway, align=center] {Marges extérieures / Repères de page};
% Barre d'état (bas)
\draw[poppink, fill=poppink!10] (0,0) rectangle (16,0.5) node[midway] {BARRE D'ÉTAT (Zoom, Nombre de pages, Mode d'affichage)};
\end{tikzpicture}
\end{center}
Légende les zones numérotées (1 à 5) en associant chaque numéro à son nom technique et à sa fonction principale.
\end{exercice}

\begin{exercice}[Démarche de création : Carte de visite depuis un modèle]
Tu dois créer une carte de visite pour le \textbf{Cybercafé "Connexion Rapide"} à Yaoundé (Propriétaire : M. Ngono, Tél : 6XX XX XX XX, Email : contact@connexion.cm).
Rédige la démarche complète, étape par étape, depuis le lancement de Publisher jusqu'à l'enregistrement du fichier, en utilisant \textbf{un modèle prédéfini} (Modèle intégré "Carte de visite").
Indique précisément :
\begin{enumerate}
    \item Comment accéder aux modèles intégrés.
    \item Comment choisir le bon format (orientation, dimensions standard).
    \item Comment remplacer les champs génériques (Nom, Société, Logo) par les informations réelles.
    \item Comment gérer le verso (recto-verso) si le modèle le propose.
\end{enumerate}
\end{exercice}

\begin{exercice}[Analyse critique : Affiche mal conçue]
On te présente une affiche de sensibilisation contre le paludisme réalisée par un élève. Voici les défauts observés :
\begin{itemize}
    \item Utilisation de \textbf{5 polices de caractères différentes} (Comic Sans, Times New Roman, Arial, Impact, Courier New).
    \item L'image principale (moustique) est \textbf{étirée déformée} (rapport hauteur/largeur non conservé).
    \item Un bloc de texte important dépasse dans la marge extérieure (hors zone imprimable).
    \item Les couleurs de fond (vert fluo) et du texte (jaune clair) offrent un \textbf{contraste très faible}.
    \item Aucun repère de rognage n'a été défini pour l'imprimeur.
\end{itemize}
Pour chaque défaut, explique \textbf{la conséquence} sur le document final (lisibilité, professionnalisme, impression) et propose \textbf{la correction technique} à appliquer dans Publisher (outils/menu exacts).
\end{exercice}

\begin{exercice}[Mise en page : Journal scolaire 4 pages]
Tu dois produire le gabarit (maquette vide) d'un journal scolaire de 4 pages (format A3 plié en 2 = 4 pages A4) dans Publisher.
\begin{enumerate}
    \item Quelle option de "Taille de page" choisis-tu au départ (A3 ou A4) ? Justifie.
    \item Comment configures-tu les \textbf{marges intérieures} (côté pliure) plus larges que les marges extérieures pour faciliter la lecture après pliage/agrafage ?
    \item Décris comment insérer un \textbf{numéro de page automatique} qui n'apparaît pas sur la page de garde (page 1) mais commence à 2 sur la page 2.
    \item Comment créer des \textbf{colonnes de largeur égale} avec une gouttière de 0,6 cm sur les pages intérieures (pages 2 et 3) ?
\end{enumerate}
\end{exercice}

\courssub{Je me prépare au Probatoire}

\begin{exercice}[Sujet type Probatoire : Production d'un faire-part de mariage -- Contexte complet]
\textbf{Situation :} La famille \textbf{Mvondo} prépare le mariage de leur fille aînée, \textbf{Marie}, avec \textbf{Jean}, le 15 août 2026 à 15h00 à la \textbf{Paroisse Saint-Jean-Baptiste de Mvolyé} (Yaoundé). La réception suivra à la \textbf{Salle des Fêtes de l'Hôtel de Ville de Yaoundé}. Le code vestimentaire est "Tenue blanche avec touche d'or". RSVP avant le 1er août au 6XX XX XX XX.

\textbf{Travail demandé :} Tu dois produire ce faire-part en utilisant Microsoft Publisher (ou équivalent LibreOffice Draw/Scribus) en respectant les contraintes techniques suivantes :
\begin{enumerate}
    \item \textbf{Format :} Carré 14 cm $\times$ 14 cm (format carte d'invitation), orientation portrait.
    \item \textbf{Marges :} 1,5 cm de chaque côté. Une gouttière centrale de 0,5 cm si mise en page 2 colonnes pour le texte.
    \item \textbf{Objets obligatoires :}
        \begin{itemize}
            \item Un cadre décoratif (forme rectangle arrondi) en bordure de page (marge 0,5 cm du bord physique).
            \item Une image vectorielle (icône alliances ou colombes) centrée en haut, taille 3 cm $\times$ 3 cm.
            \item Deux zones de texte liées (chaînées) pour le corps du texte (noms, date, lieux, RSVP) en police \texttt{Garamond} taille 12 pt, justifié, interligne 1,15.
            \item Un fond de page texturé (effet parchemin ou lin) clair.
        \end{itemize}
    \item \textbf{Contraintes de qualité :} Alignement parfait des blocs (magnétisme/alignement relatif), respect des zones de sécurité (pas de texte à moins de 3 mm du bord de coupe), couleurs CMJN (Bleu marine \#001F3F et Or \#FFD700).
\end{enumerate}

\textbf{Questions :}
\begin{enumerate}
    \item Dessine (schéma annoté au crayon) le \textbf{gabarit de la page} (repères, marges, gouttière, position approximative des 4 objets obligatoires) avec les côtes en cm.
    \item Rédige la \textbf{procédure pas à pas} (10 à 12 étapes max) pour réaliser ce document depuis une page blanche, en citant les onglets et groupes du Ruban utilisés (ex: \textit{Onglet Insertion > Groupe Illustrations > Formes}).
    \item Comment \textbf{vérifier} avant impression que le document respecte les marges de sécurité (aperçu avant impression, repères de rognage) ?
    \item Quelle \textbf{procédure d'exportation} choisis-tu pour envoyer le fichier à l'imprimeur (format, options : traits de coupe, résolution images, polices incorporées) ? Justifie.
    \item \textbf{Justification des choix (Savoir-être) :} Explique en 5 lignes pourquoi l'utilisation d'un logiciel de PAO (Publisher) est plus adaptée qu'un traitement de texte (Word) pour ce type de document "carte d'invitation haut de gamme".
\end{enumerate}
\end{exercice}

\begin{exercice}[Sujet type Probatoire : Dépliant touristique 3 volets -- Calculs et Mise en œuvre]
\textbf{Contexte :} L'Office du Tourisme de la région de l'Ouest (Bafoussam) veut un dépliant 3 volets (2 plis roulés) format A4 fermé (donc A4 paysage ouvert : 29,7 cm $\times$ 21 cm).

\textbf{Données techniques :}
\begin{itemize}
    \item Marges extérieures (gauche/droite) : 1,5 cm.
    \item Gouttières entre volets (espace pour le pli) : 0,6 cm (deux gouttières).
    \item Fond perdu (dépassement image pour coupe) : 3 mm de chaque côté du format fini.
\end{itemize}

\textbf{Questions :}
\begin{enumerate}
    \item \textbf{Calculs (à faire apparaître clairement) :}
        \begin{enumerate}
            \item Calculer la \textbf{largeur utile d'un volet} (zone d'impression texte/image).
            \item Calculer les \textbf{coordonnées X} (abscisses) des traits de pli (par rapport au bord gauche de la page ouverte) pour guider le massicotage.
            \item Donner le \textbf{format du fichier PDF avec fond perdu} (Largeur $\times$ Hauteur en cm) à fournir à l'imprimeur.
        \end{enumerate}
    \item \textbf{Configuration dans Publisher :} Décris comment configurer exactement ces valeurs dans la boîte de dialogue \textit{Mise en page > Marges > Marges personnalisées} et \textit{Mise en page > Colonnes}.
    \item \textbf{Organisation du contenu :} Le volet 1 (couverture) porte le titre "L'Ouest, Cœur du Cameroun". Le volet 6 (derrière, face cachée quand plié) contient les mentions légales et contacts. Propose une répartition logique des 4 volets intérieurs (2, 3, 4, 5) pour : "Sites touristiques", "Hébergement/Restauration", "Culture/Artisanat", "Accès/Transport". Justifie l'ordre de lecture (pliage roulé).
    \item \textbf{Gestion des images :} Une photo de la chefferie de Bafut (format original 4000 $\times$ 3000 px) doit occuper toute la hauteur du volet 2 (hauteur utile 18 cm) en fond perdu. Quelle \textbf{résolution en PPP (DPI)} aura l'image une fois placée ? Est-ce suffisant pour l'impression offset (300 DPI requis) ? Si non, que faire ?
\end{enumerate}
\end{exercice}

\begin{exercice}[Sujet type Probatoire : Affiche de sensibilisation -- Production pratique \& Export PDF]
\textbf{Mise en situation réelle (TP noté) :} En salle informatique, tu as 45 minutes pour produire une affiche A3 (29,7 $\times$ 42 cm) portrait : \textbf{"Zéro Palu : Je me protège, je protège ma famille"}.

\textbf{Cahier des charges strict :}
\begin{enumerate}
    \item \textbf{Arrière-plan :} Dégradé vertical du blanc (haut) vers un bleu clair \#D6EAF8 (bas).
    \item \textbf{En-tête (haut 15\%)} : Titre principal en \texttt{Montserrat ExtraBold} 48 pt, Blanc, Ombre portée noire (décalage 2 pt), centré. Sous-titre en \texttt{Montserrat Regular} 24 pt, Bleu marine \#1B4F72.
    \item \textbf{Corps central (70\%)} : Trois \textbf{blocs identiques} (largeur 30\% chacun, hauteur 80\% zone), espacés équitablement (alignement : distribution horizontale). Chaque bloc = 1 Icône SVG (Moustiquaire, Spray, Dépistage) 8 cm $\times$ 8 cm + Texte explicatif 3 lignes (Police \texttt{Open Sans} 16 pt, Gris foncé \#34495E, aligné centre).
    \item \textbf{Pied de page (bas 10\%)} : Logos MINSANTE / OMS / PNLP alignés horizontalement, hauteur 2 cm. Texte "Numéro vert : 1510" en \texttt{Montserrat Bold} 20 pt, Rouge \#E74C3C, centré.
    \item \textbf{Contraintes techniques :} Repères de rognage activés. Fond perdu 3 mm. Toutes les images vectorielles (SVG). Polices incorporées.
\end{enumerate}

\textbf{Questions d'examen (écrites après la production) :}
\begin{enumerate}
    \item Quelle \textbf{fonctionnalité de Publisher} as-tu utilisée pour que les 3 blocs centraux aient \textbf{exactement la même taille et le même alignement vertical} sans les redimensionner manuellement un par un ? (Nom précis de l'outil).
    \item Comment as-tu procédé pour \textbf{incorporer les polices} (Montserrat, Open Sans) dans le fichier Publisher afin qu'elles s'affichent correctement sur le poste de l'imprimeur qui ne les a pas installées ?
    \item Décris la procédure exacte d'\textbf{exportation en PDF "Prêt pour l'impression commerciale"} (chemin menus, options cochées : traits de coupe, fond perdu, résolution 300 PPP, standard PDF/X-1a si dispo).
    \item \textbf{Auto-évaluation :} Coche les critères respectés sur ta production :
    \begin{tabular}{|p{5cm}|c|c|}
    \hline
    \textbf{Critère} & \textbf{Oui} & \textbf{Non} \\ \hline
    Format page A3 + fond perdu 3 mm configuré & & \\ \hline
    3 blocs centraux alignés/distribués parfaitement & & \\ \hline
    Icônes vectorielles (SVG) non pixellisées au zoom 400\% & & \\ \hline
    Export PDF avec traits de coupe et polices incorporées & & \\ \hline
    \end{tabular}
\end{enumerate}
\end{exercice}

\newpage

\coursec{Corrigés des exercices}

\begin{corrige}[Exercice 1 — QCM : Identification d'un logiciel de PAO]
\textbf{Bonne réponse : 2. Microsoft Publisher.}

\textbf{Justification :}
\begin{itemize}
    \item \textbf{Microsoft Word} est un \cle{traitement de texte} : le texte s'écrit en continu, de haut en bas, et la mise en page suit le flux du texte. Il est conçu pour les lettres, rapports et dissertations.
    \item \textbf{Microsoft Publisher} est un logiciel de \cle{PAO} (Publication Assistée par Ordinateur) : chaque élément (texte, image, forme) est placé dans un \textbf{objet conteneur} (zone de texte, cadre image) que l'on positionne librement et précisément sur la page. Il est dédié à la mise en page professionnelle : affiches, dépliants, cartes de visite, faire-part.
    \item \textbf{Microsoft Excel} est un tableur (calculs, tableaux).
    \item \textbf{Microsoft Paint} est un logiciel de dessin simple (images matricielles).
\end{itemize}

\textbf{Différence fondamentale :} dans un traitement de texte, \emph{le texte pilote la page} (on tape, la page se remplit) ; dans un logiciel de PAO, \emph{la page pilote le contenu} (on dessine d'abord la maquette avec des cadres, puis on y place le contenu, au millimètre près).
\end{corrige}

\begin{corrige}[Exercice 2 — Vrai ou Faux justifié]
\begin{enumerate}
    \item \textbf{FAUX.} Dans un logiciel de PAO, tout texte doit être placé dans une \cle{zone de texte} (cadre) que l'on crée, dimensionne et déplace. On ne peut pas taper directement « à même la page » comme dans Word.
    \item \textbf{FAUX.} Les repères (marges, guides de colonnes, gouttières) sont des \textbf{aides à la conception} affichées à l'écran uniquement (lignes bleues/roses non imprimables). Seuls les \emph{repères de rognage} (traits de coupe), ajoutés volontairement à l'export PDF, apparaissent sur le document destiné à l'imprimeur — et ils sont situés hors du format fini.
    \item \textbf{FAUX.} Le format \texttt{.pub} est un format \textbf{propriétaire natif} de Publisher : Word ne sait pas l'ouvrir. Pour partager le document, il faut l'\textbf{exporter en PDF} (non modifiable) ou l'enregistrer dans un format d'échange.
    \item \textbf{VRAI.} Le PDF \cle{fige la mise en page} : polices incorporées, images, positions et couleurs sont conservées à l'identique sur tout poste. C'est le format standard exigé par les imprimeurs professionnels (avec traits de coupe et fond perdu).
\end{enumerate}
\end{corrige}

\begin{corrige}[Exercice 3 — Définitions et vocabulaire technique]
\begin{enumerate}
    \item \textbf{Gouttière :} espace vertical vide qui sépare deux colonnes de texte (ou deux volets d'un dépliant). Elle évite que les textes des colonnes voisines se touchent et ménage la place du pli sur un dépliant.
    \item \textbf{Repères de rognage (traits de coupe) :} petits traits fins placés aux quatre coins du document, hors de la zone imprimée, qui indiquent à l'imprimeur où \textbf{couper} le papier pour obtenir le format fini exact.
    \item \textbf{Volet des pages :} panneau latéral (à gauche par défaut) de l'interface de Publisher qui affiche les \textbf{miniatures} de toutes les pages de la publication. Il permet de naviguer, d'ajouter, de supprimer, de dupliquer ou de réordonner les pages.
    \item \textbf{Zone de texte liée (chaînage) :} liaison créée entre deux zones de texte de sorte que le texte qui déborde de la première \textbf{se poursuit automatiquement} dans la seconde. Modifier la taille de la première zone répercute le texte dans la zone liée (outil \textit{Créer un lien} / \textit{Chaîner}).
\end{enumerate}
\end{corrige}

\begin{corrige}[Exercice 4 — Calcul direct : Dimensions d'un dépliant]
\textbf{Données :} feuille A4 paysage $\Rightarrow$ largeur totale $L = 29,7$ cm ; marges gauche + droite : $1 + 1 = 2$ cm ; 3 volets $\Rightarrow$ 2 gouttières de $0,5$ cm, soit $1$ cm au total.

\textbf{Largeur disponible pour les volets :}
\[
L_{\text{dispo}} = 29,7 - (1 + 1) - (2 \times 0,5) = 29,7 - 2 - 1 = 26,7 \text{ cm}
\]

\textbf{Largeur utile de chaque volet :}
\[
L_{\text{volet}} = \frac{26,7}{3} = 8,9 \text{ cm}
\]

\textbf{Vérification :} $3 \times 8,9 + 2 \times 0,5 + 2 \times 1 = 26,7 + 1 + 2 = 29,7$ cm. \checkmark

\textbf{Conclusion :} chaque volet a une largeur utile d'impression de \textbf{8,9 cm}.
\end{corrige}

\begin{corrige}[Exercice 5 — Légender l'interface de Publisher]
\begin{center}
\begin{tabularx}{\linewidth}{|c|l|X|}
\hline
\rowcolor{popblueL} \textbf{N°} & \textbf{Nom technique} & \textbf{Fonction principale} \\ \hline
1 & Barre de titre & Affiche le nom du fichier et du logiciel ; contient les boutons Réduire / Agrandir / Fermer. \\ \hline
2 & Ruban & Regroupe toutes les commandes en onglets (Accueil, Insertion, Mise en page\ldots) et en groupes ; c'est l'outil principal de mise en forme. \\ \hline
3 & Volet des pages & Affiche les miniatures des pages ; permet de naviguer, ajouter, supprimer, réordonner les pages. \\ \hline
4 & Zone de travail (page courante) & Espace de conception où l'on place les objets (zones de texte, images, formes) ; les repères en pointillés matérialisent les marges (zone imprimable). \\ \hline
5 & Barre d'état & Affiche le numéro de page, le niveau de zoom, le mode d'affichage ; permet de zoomer rapidement. \\ \hline
\end{tabularx}
\end{center}

\textbf{Point de vigilance :} les lignes en pointillés (marges intérieures/extérieures) sont des \textbf{repères non imprimables} : ils guident le placement des objets mais n'apparaîtront jamais sur le papier.
\end{corrige}

\begin{corrige}[Exercice 6 — Démarche de création : Carte de visite depuis un modèle]
\textbf{Méthode : Création de la carte de visite « Connexion Rapide »}
\begin{enumerate}
    \item \textbf{Lancer Publisher} : menu Démarrer $\rightarrow$ Microsoft Publisher.
    \item \textbf{Accéder aux modèles intégrés} : dans l'écran d'accueil, cliquer sur \textit{Nouveau} $\rightarrow$ \textit{Modèles intégrés} $\rightarrow$ catégorie \textit{Cartes de visite}.
    \item \textbf{Choisir le format} : sélectionner un modèle en \textbf{orientation paysage}, dimensions standard d'une carte de visite (8,5 cm $\times$ 5,4 cm). Vérifier dans le volet de droite : \textit{Taille de page} et \textit{Jeu de couleurs} (choisir un jeu sobre, 2 couleurs maximum).
    \item \textbf{Créer} : cliquer sur le bouton \textit{Créer} ; le modèle s'ouvre avec des champs génériques.
    \item \textbf{Remplacer les champs texte} : cliquer dans chaque zone de texte générique et saisir les informations réelles :
    \begin{itemize}
        \item Nom de la société : \textbf{Cybercafé « Connexion Rapide »}
        \item Propriétaire : \textbf{M. Ngono}
        \item Adresse : \textbf{Yaoundé — Cameroun}
        \item Téléphone : \textbf{6XX XX XX XX} ; Email : \textbf{contact@connexion.cm}
    \end{itemize}
    \item \textbf{Remplacer le logo} : clic droit sur l'image générique $\rightarrow$ \textit{Changer l'image} $\rightarrow$ \textit{À partir d'un fichier}, sélectionner le logo du cybercafé. Vérifier que l'image n'est pas déformée (poignées d'angle uniquement).
    \item \textbf{Gérer le verso} : si le modèle propose 2 pages (recto/verso), utiliser le \textbf{volet des pages} pour passer à la page 2 et y saisir les informations complémentaires (services : impression, photocopie, formation bureautique). Sinon, \textit{Insertion $\rightarrow$ Page} pour créer un verso.
    \item \textbf{Vérifier} : contrôler l'orthographe, l'alignement (repères magnétiques) et l'aperçu (\textit{Fichier $\rightarrow$ Imprimer}).
    \item \textbf{Enregistrer} : \textit{Fichier $\rightarrow$ Enregistrer sous} $\rightarrow$ nom : \texttt{carte\_visite\_connexion\_rapide.pub}. Pour l'imprimeur : \textit{Fichier $\rightarrow$ Exporter $\rightarrow$ Créer un document PDF}.
\end{enumerate}

\textbf{Points de vigilance :} ne jamais étirer le logo par les poignées latérales (déformation) ; garder au maximum 2 polices ; laisser le texte à au moins 3 mm du bord de coupe.
\end{corrige}

\begin{corrige}[Exercice 7 — Analyse critique : Affiche mal conçue]
\begin{center}
\begin{tabularx}{\linewidth}{|l|X|X|}
\hline
\rowcolor{popblueL} \textbf{Défaut} & \textbf{Conséquence} & \textbf{Correction technique dans Publisher} \\ \hline
5 polices différentes & Document désordonné, amateur, fatigant à lire ; perte de hiérarchie visuelle. & Se limiter à \textbf{2 polices} (une pour les titres, une pour le texte). Sélectionner tout le texte (\texttt{Ctrl+A} dans chaque zone) et appliquer les polices via \textit{Onglet Accueil $\rightarrow$ Groupe Police}. \\ \hline
Image étirée/déformée & Moustique méconnaissable, rendu non professionnel. & Supprimer l'image, la réinsérer (\textit{Insertion $\rightarrow$ Images}) et la redimensionner \textbf{uniquement par les poignées d'angle} (ou \textit{Format de l'image $\rightarrow$ Taille} en cochant \textit{Conserver les proportions}). \\ \hline
Texte hors zone imprimable & Le texte sera \textbf{coupé} ou non imprimé ; perte d'information. & Déplacer la zone de texte à l'intérieur des repères de marge (laisser 3 mm de sécurité) ; réduire la taille de la zone ou du texte si nécessaire. \\ \hline
Contraste faible (vert fluo / jaune clair) & Affiche \textbf{illisible}, surtout à distance ; message de sensibilisation inefficace. & Choisir un couple de couleurs contrastées (texte foncé sur fond clair, ex. bleu marine sur blanc) : \textit{Accueil $\rightarrow$ Couleur de police} et \textit{Format de la forme $\rightarrow$ Remplissage}. \\ \hline
Pas de repères de rognage & L'imprimeur ne sait pas où couper ; format final imprécis, bords blancs possibles. & À l'export : \textit{Fichier $\rightarrow$ Exporter $\rightarrow$ PDF $\rightarrow$ Options} $\rightarrow$ cocher \textit{Repères de coupe} (et fond perdu 3 mm si l'image touche le bord). \\ \hline
\end{tabularx}
\end{center}

\textbf{Conclusion :} une bonne affiche respecte trois règles : \textbf{lisibilité} (contraste, peu de polices), \textbf{fidélité} (images non déformées), \textbf{imprimabilité} (texte dans les marges, traits de coupe).
\end{corrige}

\begin{corrige}[Exercice 8 — Mise en page : Journal scolaire 4 pages]
\begin{enumerate}
    \item \textbf{Taille de page : A4.} Publisher travaille page par page : on conçoit 4 pages A4 (la page de garde, deux pages intérieures, la dernière de couverture). C'est l'imprimeur (ou le photocopieur en mode « livret ») qui assemble les 4 pages A4 sur une feuille A3 pliée en deux. Choisir A3 obligerait à gérer soi-même l'imposition (ordre des pages), ce qui dépasse notre niveau.
    \item \textbf{Marges intérieures plus larges :} \textit{Onglet Mise en page $\rightarrow$ Marges $\rightarrow$ Marges personnalisées}. Définir par exemple : marge \textbf{intérieure (côté pliure/agrafes) : 2 cm} ; marge extérieure : 1,2 cm ; haut/bas : 1,5 cm. La marge intérieure plus large évite que le texte disparaisse dans la pliure après agrafage.
    \item \textbf{Numéro de page automatique sauf page 1 :}
    \begin{enumerate}
        \item Aller sur la page 2, puis \textit{Insertion $\rightarrow$ Numéro de page} $\rightarrow$ choisir la position (ex. bas de page, centré).
        \item Le numéro est inséré sur la \textbf{page de fond} (page maîtresse) et s'incrémente automatiquement.
        \item Pour que la page 1 n'affiche pas de numéro : créer une \textbf{page maîtresse distincte} sans numéro et l'appliquer à la page 1 (\textit{Mise en page $\rightarrow$ Pages maîtresses} $\rightarrow$ \textit{Appliquer la page maîtresse}), ou masquer l'objet numéro sur la page 1.
    \end{enumerate}
    \item \textbf{Colonnes égales avec gouttière 0,6 cm (pages 2 et 3) :} sélectionner la page concernée, \textit{Mise en page $\rightarrow$ Repères $\rightarrow$ Repères de grille et de ligne de base} (ou \textit{Colonnes} selon la version) $\rightarrow$ définir \textbf{Nombre de colonnes : 2 ou 3} et \textbf{Espacement (gouttière) : 0,6 cm}. Les colonnes créées ont automatiquement une largeur égale. On place ensuite une zone de texte par colonne, que l'on \textbf{chaîne} (lier) pour que le texte coule de l'une à l'autre.
\end{enumerate}

\textbf{Point de vigilance :} les repères de colonnes ne créent pas le texte : il faut ensuite dessiner les zones de texte sur ces repères et les lier.
\end{corrige}

\begin{corrige}[Exercice 9 — Sujet type Probatoire : Faire-part de mariage]
\textbf{1. Gabarit de la page (schéma annoté) :}

\begin{popfigure}
\begin{tikzpicture}[font=\small, scale=0.55]
% Page physique 14x14
\draw[thick, popdark] (0,0) rectangle (14,14);
\node[below] at (7,0) {14 cm};
\node[left, rotate=90] at (0,7) {14 cm};
% Cadre décoratif à 0.5 cm du bord
\draw[popgold, thick, rounded corners=8pt] (0.5,0.5) rectangle (13.5,13.5);
\node[popgold, anchor=west] at (13.6,13) {\footnotesize Cadre (0,5 cm du bord)};
% Marges 1.5 cm
\draw[dashed, popblue] (1.5,1.5) rectangle (12.5,12.5);
\node[popblue, anchor=west] at (12.6,1.2) {\footnotesize Marges 1,5 cm};
% Gouttière centrale
\draw[dashed, poporange] (6.75,1.5) -- (6.75,12.5);
\draw[dashed, poporange] (7.25,1.5) -- (7.25,12.5);
\node[poporange] at (7,1.0) {\footnotesize Gouttière 0,5 cm};
% Image centrée en haut
\draw[popgreen, fill=popgreenL] (5.5,9.5) rectangle (8.5,12.5);
\node[popgreen] at (7,11) {\footnotesize Image 3$\times$3 cm};
% Zones de texte liées
\draw[poppurple, fill=poppurpleL] (1.5,2) rectangle (6.75,8.5);
\node[poppurple, align=center] at (4.1,5.2) {\footnotesize Zone texte 1\\ \footnotesize (Garamond 12, justifié)};
\draw[poppurple, fill=poppurpleL] (7.25,2) rectangle (12.5,8.5);
\node[poppurple, align=center] at (9.9,5.2) {\footnotesize Zone texte 2\\ \footnotesize (liée à la zone 1)};
\draw[-{Stealth}, poppurple, thick] (6.75,4) -- (7.25,4);
\end{tikzpicture}
\legende{Gabarit du faire-part de mariage (format 14 $\times$ 14 cm) avec marges, gouttière, cadre décoratif, image centrée et zones de texte chaînées.}
\end{popfigure}

\textbf{2. Procédure pas à pas (10 étapes) :}
\begin{enumerate}
    \item Lancer Publisher $\rightarrow$ \textit{Nouveau} $\rightarrow$ \textit{Page vierge}.
    \item \textit{Onglet Mise en page $\rightarrow$ Taille $\rightarrow$ Autres tailles} : définir Largeur 14 cm, Hauteur 14 cm, orientation portrait.
    \item \textit{Mise en page $\rightarrow$ Marges $\rightarrow$ Marges personnalisées} : 1,5 cm partout.
    \item \textit{Mise en page $\rightarrow$ Arrière-plan de page} : choisir une texture claire (parchemin/lin).
    \item \textit{Insertion $\rightarrow$ Formes $\rightarrow$ Rectangle arrondi} : dessiner le cadre à 0,5 cm du bord ; \textit{Format de la forme} : remplissage « Aucun », contour Or \texttt{\#FFD700}, épaisseur 2 pt.
    \item \textit{Insertion $\rightarrow$ Images} : insérer l'icône (alliances/colombes), la dimensionner à 3 cm $\times$ 3 cm (\textit{Format $\rightarrow$ Taille}) et la centrer en haut (\textit{Format $\rightarrow$ Aligner $\rightarrow$ Centrer horizontalement}).
    \item \textit{Insertion $\rightarrow$ Zone de texte} : dessiner la zone 1 (colonne gauche) ; taper le texte (Marie \& Jean, 15 août 2026 à 15h, Paroisse Saint-Jean-Baptiste de Mvolyé, réception Hôtel de Ville, RSVP).
    \item Dessiner la zone 2 (colonne droite), sélectionner la zone 1 $\rightarrow$ \textit{Format de la zone de texte $\rightarrow$ Créer un lien} $\rightarrow$ cliquer sur la zone 2 (chaînage).
    \item Mettre en forme le texte : police \texttt{Garamond}, 12 pt, justifié (\textit{Accueil $\rightarrow$ Paragraphe}), interligne 1,15 ; couleur Bleu marine \texttt{\#001F3F}.
    \item Vérifier l'alignement (repères magnétiques, \textit{Aligner les objets}), enregistrer en \texttt{.pub}, puis exporter en PDF.
\end{enumerate}

\textbf{3. Vérification avant impression :} \textit{Fichier $\rightarrow$ Imprimer} pour afficher l'\textbf{aperçu avant impression} : contrôler qu'aucun texte ne s'approche à moins de 3 mm du bord de coupe (zone de sécurité). Zoomer à 100 \% (taille réelle) pour juger la lisibilité. Vérifier à l'écran que tous les objets restent à l'intérieur des repères de marge, et activer les \textbf{repères de rognage} dans les options d'export PDF pour visualiser l'emplacement des coupes.

\textbf{4. Exportation pour l'imprimeur :} \textit{Fichier $\rightarrow$ Exporter $\rightarrow$ Créer un document PDF/XPS $\rightarrow$ Options} :
\begin{itemize}
    \item cocher \textbf{Repères de coupe} (traits de rognage) ;
    \item choisir l'optimisation \textbf{« Impression professionnelle »} (images en 300 PPP, polices incorporées) ;
    \item si disponible, norme \textbf{PDF/X-1a} (couleurs CMJN, polices intégrées).
\end{itemize}
\textbf{Justification :} le PDF fige la mise en page, embarque les polices (Garamond) et conserve les couleurs CMJN exigées en imprimerie ; les traits de coupe garantissent un massicotage exact au format 14 $\times$ 14 cm.

\textbf{5. Justification (savoir-être) :} un faire-part haut de gamme exige un \textbf{placement précis au millimètre} de chaque objet (cadre, image centrée, colonnes chaînées), ce qu'un traitement de texte ne permet pas : dans Word, les objets « flottent » difficilement et le texte pousse la mise en page. Publisher offre les \textbf{repères magnétiques}, l'\textbf{alignement relatif des objets}, le \textbf{chaînage des zones de texte} et l'\textbf{export PDF professionnel} avec traits de coupe. La PAO sépare clairement la maquette du contenu : c'est l'outil du métier d'imprimeur.

\textbf{Barème indicatif (sur 20) :} gabarit coté correct : 4 pts ; procédure complète et ordonnée (onglets cités) : 6 pts ; vérification : 3 pts ; export PDF justifié : 4 pts ; argumentation PAO vs traitement de texte : 3 pts.

\textbf{Points de vigilance :} exiger les \textbf{côtes en cm} sur le schéma ; le candidat doit citer les \textbf{noms exacts des onglets et groupes} ; la justification doit mentionner au moins deux fonctionnalités propres à la PAO (chaînage, alignement, repères, export PDF).
\end{corrige}

\begin{corrige}[Exercice 10 — Sujet type Probatoire : Dépliant touristique 3 volets]
\textbf{1. Calculs :}

\textbf{a) Largeur utile d'un volet :}
\[
L_{\text{dispo}} = 29,7 - (2 \times 1,5) - (2 \times 0,6) = 29,7 - 3 - 1,2 = 25,5 \text{ cm}
\]
\[
L_{\text{volet}} = \frac{25,5}{3} = 8,5 \text{ cm}
\]

\textbf{b) Coordonnées X des traits de pli} (depuis le bord gauche de la page ouverte) :
\begin{align*}
\text{Volet 1} &: 0 \rightarrow 1,5 + 8,5 = 10,0 \text{ cm (fin volet 1)}\\
\text{Gouttière 1} &: 10,0 \rightarrow 10,6 \text{ cm} \quad \Rightarrow \quad \textbf{pli 1 : } X_1 = 10,0 \text{ cm (bord du volet)}\\
\text{Volet 2} &: 10,6 \rightarrow 19,1 \text{ cm}\\
\text{Gouttière 2} &: 19,1 \rightarrow 19,7 \text{ cm} \quad \Rightarrow \quad \textbf{pli 2 : } X_2 = 19,7 \text{ cm}\\
\text{Volet 3} &: 19,7 \rightarrow 28,2 \text{ cm, puis marge } 1,5 \Rightarrow 29,7 \text{ cm \checkmark}
\end{align*}
Les deux plis se situent donc aux abscisses \textbf{$X_1 = 10,0$ cm} et \textbf{$X_2 = 19,7$ cm} (bords des volets ; le pli se fait au milieu de chaque gouttière, soit 10,3 cm et 19,4 cm — les deux conventions sont acceptées si elles sont expliquées).

\textbf{Vérification :} $1,5 + 8,5 + 0,6 + 8,5 + 0,6 + 8,5 + 1,5 = 29,7$ cm. \checkmark

\textbf{c) Format PDF avec fond perdu :} le fond perdu de 3 mm s'ajoute de \textbf{chaque côté} :
\[
\text{Largeur} = 29,7 + 0,3 + 0,3 = 30,3 \text{ cm} \qquad \text{Hauteur} = 21 + 0,3 + 0,3 = 21,6 \text{ cm}
\]
Fichier à fournir : \textbf{30,3 cm $\times$ 21,6 cm} (avec traits de coupe indiquant le format fini 29,7 $\times$ 21 cm).

\textbf{2. Configuration dans Publisher :}
\begin{itemize}
    \item \textit{Mise en page $\rightarrow$ Taille} : A4, orientation \textbf{Paysage}.
    \item \textit{Mise en page $\rightarrow$ Marges $\rightarrow$ Marges personnalisées} : Gauche = 1.5cm ; Droite = 1.5cm ; Haut/Bas selon besoin (ex. 1,5 cm).
    \item \textit{Mise en page $\rightarrow$ Repères $\rightarrow$ Repères de grille} (ou \textit{Colonnes}) : \textbf{Nombre de colonnes = 3}, \textbf{Espacement (gouttière) = 0.6cm}. Publisher calcule alors automatiquement des colonnes de 8,5 cm.
    \item Pour le fond perdu : agrandir la taille de page personnalisée à 30,3 $\times$ 21,6 cm et placer des repères à 3 mm du bord, ou gérer le fond perdu à l'export PDF (\textit{Options $\rightarrow$ Fond perdu}).
\end{itemize}

\textbf{3. Répartition des volets (pliage roulé) :}
\begin{center}
\begin{tabularx}{\linewidth}{|c|l|X|}
\hline
\rowcolor{popblueL} \textbf{Volet} & \textbf{Contenu} & \textbf{Justification} \\ \hline
1 & Couverture : « L'Ouest, Cœur du Cameroun » & Seul volet visible quand le dépliant est fermé : titre + image d'accroche. \\ \hline
2 & Sites touristiques (chefferies, chutes de la Métché\ldots) & Premier volet découvert à l'ouverture : on commence par l'essentiel (quoi visiter). \\ \hline
3 & Culture / Artisanat & Poursuite naturelle de la découverte. \\ \hline
4 & Hébergement / Restauration & Informations pratiques après l'envie de visiter. \\ \hline
5 & Accès / Transport (carte, axes routiers) & Dernières informations logistiques avant le contact. \\ \hline
6 & Mentions légales, contacts, horaires & Dos du dépliant fermé : informations de service. \\ \hline
\end{tabularx}
\end{center}
\textbf{Justification de l'ordre :} avec un pliage roulé, le lecteur découvre les volets dans l'ordre 1 $\rightarrow$ 2 $\rightarrow$ 3 $\rightarrow$ 4 $\rightarrow$ 5 ; on suit donc la logique du voyageur : \emph{rêver (sites) $\rightarrow$ découvrir (culture) $\rightarrow$ se loger $\rightarrow$ venir (transport) $\rightarrow$ contacter}.

\textbf{4. Résolution de l'image :} l'image fait 3000 px de haut pour une hauteur placée de 18 cm :
\[
18 \text{ cm} = \frac{18}{2,54} \approx 7,09 \text{ pouces} \qquad \text{Résolution} = \frac{3000}{7,09} \approx 423 \text{ PPP}
\]
\textbf{Conclusion :} 423 PPP $>$ 300 PPP : la résolution est \textbf{largement suffisante} pour l'impression offset. Aucune action corrective n'est nécessaire ; il faut seulement veiller à ne pas \emph{agrandir} l'image au-delà de sa taille (sinon la résolution chuterait : à 25 cm de haut, on tomberait à $\frac{3000}{9,84} \approx 305$ PPP, limite acceptable).

\textbf{Barème indicatif (sur 20) :} calculs a, b, c détaillés et vérifiés : 6 pts ; configuration Publisher : 4 pts ; répartition justifiée : 5 pts ; calcul de résolution et conclusion : 5 pts.

\textbf{Points de vigilance :} ne pas oublier les \textbf{deux} gouttières ni les \textbf{deux} marges ; le fond perdu s'ajoute \textbf{des deux côtés} ($2 \times 3$ mm par dimension) ; la résolution se calcule en \textbf{pouces} (1 pouce = 2.54cm).
\end{corrige}

\begin{corrige}[Exercice 11 — Sujet type Probatoire : Affiche « Zéro Palu »]
\textbf{1. Outil pour des blocs identiques et alignés :}
\begin{itemize}
    \item \textbf{Duplication} du premier bloc : \texttt{Ctrl+D} (ou copier-coller) garantit trois blocs de \textbf{taille strictement identique} sans redimensionnement manuel.
    \item \textbf{Alignement et distribution} : sélectionner les 3 blocs (clic + \texttt{Maj}), puis \textit{Onglet Format $\rightarrow$ Groupe Organiser $\rightarrow$ Aligner $\rightarrow$ Aligner en haut} (même alignement vertical), puis \textit{Aligner $\rightarrow$ Distribuer horizontalement} (espacements égaux).
\end{itemize}
Le nom précis de l'outil attendu : \textbf{« Aligner les objets » / « Distribuer horizontalement »} (avec les repères magnétiques activés).

\textbf{2. Incorporation des polices :}
\begin{itemize}
    \item Dans Publisher : \textit{Fichier $\rightarrow$ Options $\rightarrow$ Enregistrement} $\rightarrow$ cocher \textbf{« Incorporer les polices dans le fichier »} (Embed fonts). Le fichier \texttt{.pub} transporte alors Montserrat et Open Sans.
    \item À l'export PDF : choisir \textit{Options $\rightarrow$ Optimiser pour : Impression professionnelle}, qui \textbf{incorpore automatiquement les polices} dans le PDF. C'est la solution la plus sûre : l'imprimeur n'a pas besoin d'installer les polices.
\end{itemize}

\textbf{3. Exportation PDF « prêt pour l'impression commerciale » :}
\begin{enumerate}
    \item \textit{Fichier $\rightarrow$ Exporter $\rightarrow$ Créer un document PDF/XPS}.
    \item Cliquer sur \textbf{Options} avant de publier.
    \item Cocher : \textbf{Repères de coupe} (traits de rognage) ; \textbf{Fond perdu} (3 mm déjà prévu dans la page ou à définir) ; \textbf{Incorporer les polices}.
    \item Optimisation : \textbf{« Impression professionnelle »} (images conservées à 300 PPP minimum).
    \item Si proposé : standard \textbf{PDF/X-1a} (garantit CMJN + polices incorporées + images haute résolution).
    \item Publier, puis \textbf{rouvrir le PDF} pour vérifier : zoom 400 \% sur les icônes (vectorielles = nettes), présence des traits de coupe, format 30,3 $\times$ 42,6 cm (A3 + fond perdu).
\end{enumerate}

\textbf{4. Auto-évaluation (grille attendue) :}
\begin{center}
\begin{tabular}{|p{6cm}|c|c|}
\hline
\rowcolor{popblueL} \textbf{Critère} & \textbf{Oui} & \textbf{Non} \\ \hline
Format page A3 + fond perdu 3 mm configuré & X & \\ \hline
3 blocs centraux alignés/distribués parfaitement & X & \\ \hline
Icônes vectorielles (SVG) non pixellisées au zoom 400 \% & X & \\ \hline
Export PDF avec traits de coupe et polices incorporées & X & \\ \hline
\end{tabular}
\end{center}
Un « Non » doit conduire à corriger \textbf{avant} l'export final : vérifier la taille de page personnalisée, refaire la distribution, réinsérer les SVG d'origine (jamais une capture d'écran), refaire l'export avec les bonnes options.

\textbf{Barème indicatif (sur 20) :} production conforme au cahier des charges : 10 pts (2 pts par zone : fond, en-tête, corps, pied, contraintes techniques) ; réponse outil d'alignement/distribution : 3 pts ; incorporation des polices : 2 pts ; procédure d'export PDF détaillée : 3 pts ; auto-évaluation honnête et corrective : 2 pts.

\textbf{Points de vigilance :}
\begin{itemize}
    \item « Distribuer horizontalement » exige la \textbf{sélection multiple} des 3 blocs (\texttt{Maj} + clic) : beaucoup d'élèves l'oublient.
    \item Les images \textbf{SVG restent vectorielles} : nettes à tout zoom ; une image PNG/JPG pixellise à 400 \% — c'est le critère de vérification.
    \item Le fond perdu impose que le dégradé d'arrière-plan \textbf{dépasse de 3 mm} au-delà des traits de coupe, sinon un liseré blanc apparaîtra après massicotage.
    \item Temps de TP limité (45 min) : commencer par la \textbf{configuration de la page} (A3 + fond perdu) avant tout contenu.
\end{itemize}
\end{corrige}

\newpage

\coursec{Fiche bilan}

\begin{popfigure}
\begin{tikzpicture}[
  mindmap,
  grow cyclic,
  every node/.style={concept, font=\small},
  root concept/.append style={concept color=popblue, font=\bfseries},
  level 1/.append style={level distance=3.4cm, sibling angle=60},
  level 1 concept/.append style={font=\footnotesize},
]
\node[root concept]{La PAO}
  child[concept color=poporange]{ node[align=center]{Définition\\ publier un document\\ (affiche, journal, flyer)} }
  child[concept color=popgreen]{ node[align=center]{Interface\\ ruban, zone de travail,\\ cadres, règles} }
  child[concept color=poppurple]{ node[align=center]{Page et structure\\ format, marges,\\ orientation, guides} }
  child[concept color=poppink]{ node[align=center]{Objets\\ zones de texte, images,\\ formes, tableaux} }
  child[concept color=popgold]{ node[align=center]{Mise en forme\\ polices, couleurs,\\ habillage, alignement} }
  child[concept color=popblue!70]{ node[align=center]{Finalisation\\ vérifier, exporter (PDF),\\ imprimer} };
\end{tikzpicture}
\legende{Carte mentale de la leçon : utiliser un logiciel de publication (PAO).}
\end{popfigure}

\begin{bilan}
\textbf{Définitions clés}
\begin{itemize}
  \item \cle{PAO} (Publication Assistée par Ordinateur) : ensemble des techniques de création de documents destinés à être imprimés ou diffusés (affiches, journaux scolaires, cartes d'invitation, dépliants).
  \item \cle{Logiciel de PAO} : logiciel spécialisé dans la mise en page (ex. \emph{Microsoft Publisher}, \emph{Scribus}, \emph{Canva}), contrairement au traitement de texte qui sert surtout à rédiger.
  \item \cle{Cadre} (ou zone) : conteneur d'un objet (texte, image) que l'on place, déplace et redimensionne librement sur la page.
  \item \cle{Habillage} : façon dont le texte contourne une image.
  \item \cle{Exportation} : enregistrement de la publication dans un autre format, souvent \cle{PDF}, pour la diffusion ou l'impression.
\end{itemize}

\textbf{Repères à connaître par cœur}
\begin{center}
\begin{tabularx}{\linewidth}{|l|X|}
\hline
\rowcolor{popblueL} \textbf{Notion} & \textbf{À retenir} \\
\hline
Lancement & Démarrer $\rightarrow$ logiciel de PAO, ou double-clic sur l'icône \\
\hline
Interface & Ruban (onglets), zone de travail (page), règles, barre d'état \\
\hline
Formats de page & A4 (21 $\times$ 29,7 cm), A5 ; orientation portrait ou paysage \\
\hline
Insertion & Onglet Insertion : zone de texte, image, forme, tableau \\
\hline
Mise en forme & Police, taille, couleur, alignement, bordures, habillage \\
\hline
Sortie & Enregistrer (.pub), Exporter en PDF, Imprimer (Ctrl+P) \\
\hline
\end{tabularx}
\end{center}

\textbf{Méthode express : produire une publication}
\begin{enumerate}
  \item Définir le besoin : quel document ? pour qui ? (ex. affiche de la fête de l'école).
  \item Choisir le format de page : taille, orientation, marges, guides.
  \item Insérer les objets : zones de texte, images, formes.
  \item Mettre en forme : polices lisibles, couleurs harmonieuses, alignements.
  \item Vérifier (orthographe, débordements), enregistrer, exporter en PDF, imprimer.
\end{enumerate}
\end{bilan}

\begin{aretenir}[Checklist avant le BEPC]
\begin{itemize}
  \item $\square$ Je sais définir la PAO et citer au moins deux logiciels de publication.
  \item $\square$ Je sais distinguer un logiciel de PAO d'un traitement de texte.
  \item $\square$ Je sais lancer le logiciel et repérer les zones de son interface (ruban, page, règles).
  \item $\square$ Je sais choisir le format, l'orientation et les marges d'une page.
  \item $\square$ Je sais insérer une zone de texte, une image et une forme.
  \item $\square$ Je sais déplacer, redimensionner et aligner les cadres.
  \item $\square$ Je sais mettre en forme le texte (police, taille, couleur, alignement).
  \item $\square$ Je sais ce qu'est l'habillage du texte autour d'une image.
  \item $\square$ Je sais enregistrer ma publication et l'exporter en PDF.
  \item $\square$ Je sais justifier mes choix de mise en page pour une situation concrète (affiche, journal scolaire, invitation).
\end{itemize}
\end{aretenir}

\findecours

\end{document}
